\documentclass[11pt,a4paper]{article}
\usepackage[T1]{fontenc}
\usepackage[utf8]{inputenc}
\usepackage[a4paper,margin=25mm]{geometry}
\usepackage{amsmath}
\usepackage{booktabs}
\usepackage{longtable}
\usepackage{enumitem}
\usepackage[hidelinks]{hyperref}
\setlist{itemsep=1pt,topsep=3pt}
\title{\textbf{Literature Review Dossier}\\[2mm]\large Quantifying Ion Energy Flux Across Plasma Sheaths}
\author{Team Astrophage \\ DIA Techathon 2.0 (Naari Shakti) \\ Track: Plasma Technologies for Industry \& Society}
\date{}
\begin{document}
\maketitle
\thispagestyle{empty}

\section*{1\quad Executive Summary}
This dossier documents the literature review carried out for the project \emph{Quantifying Ion Energy Flux Across Plasma Sheaths}. In total, \textbf{731} peer-reviewed articles and preprints were located through systematic search of public scientific databases (the arXiv preprint repository, Crossref, and web search) across \textbf{24} research themes that together span the physics, measurement, modelling, and application scope of the project. Each reference was recorded with its title, authors, year, source, and a resolvable link (DOI or arXiv identifier). The themes below trace a complete path: from the first-principles physics of the sheath and the definition of ion energy flux, through the four non-intrusive optical diagnostics (optical emission spectroscopy, laser-induced fluorescence, Thomson scattering, and interferometry) and the fluid/kinetic models that constrain them, to the inverse-problem and data-fusion methodology at the core of the framework, and finally to the three application domains (electric propulsion, semiconductor plasma etching, and fusion plasma-facing components) that motivate it.

\section*{2\quad Research Coverage Matrix}
\renewcommand{\arraystretch}{1.2}
\begin{longtable}{@{}p{6.6cm}c p{6.0cm}@{}}
\toprule
\textbf{Research theme} & \textbf{\# refs} & \textbf{Role in the project} \\
\midrule\endhead
Plasma sheath theory (Bohm criterion, presheath) & 37 & First-principles model that constrains the inversion \\
RF and capacitive sheaths & 36 & Time-dependent sheath physics for pulsed/RF sources \\
Ion energy flux and power balance at surfaces & 33 & The target quantity of the framework \\
Ion energy and velocity distributions at surfaces & 27 & Distribution-level input for energy integration \\
Sputtering and ion-induced erosion & 32 & Consequence linking flux to lifetime \\
Plasma-material interaction & 28 & Surface response and boundary conditions \\
Langmuir probe diagnostics & 33 & Intrusive baseline the framework replaces \\
Optical emission spectroscopy & 33 & Passive, non-perturbing temperature/species constraint \\
Collisional-radiative models and line-ratio methods & 32 & Atomic-physics model turning light into parameters \\
Laser-induced fluorescence diagnostics & 36 & Laser diagnostic for species-resolved dynamics \\
LIF ion velocity distribution measurements & 26 & Direct measurement of the ion velocity distribution \\
Thomson scattering diagnostics & 34 & Direct, local electron temperature and density anchor \\
Plasma interferometry and electron density & 36 & Fast, robust line-integrated density constraint \\
Hall thruster erosion and lifetime & 35 & Primary application: lifetime qualification \\
Ion and gridded thruster diagnostics & 30 & Application: grid/optics erosion and performance \\
Electric propulsion plasma diagnostics & 24 & Diagnostic context for propulsion plumes \\
Plasma etching and reactive ion etching ion energy & 26 & Application: sub-nm process control \\
Semiconductor plasma processing diagnostics & 23 & Process-monitoring context \\
Fusion divertor and plasma-facing components & 32 & Application: heat/erosion at plasma-facing walls \\
Tokamak edge and scrape-off-layer diagnostics & 37 & Edge-plasma diagnostic context \\
Particle-in-cell sheath simulation & 28 & Kinetic simulation providing synthetic ground truth \\
Kinetic modelling of plasma sheaths & 18 & Kinetic theory beyond fluid closure \\
Inverse problems and regularization in diagnostics & 29 & Mathematical core of the reconstruction \\
Bayesian inference and integrated data analysis & 26 & Multi-diagnostic fusion methodology \\
\midrule
\textbf{Total} & \textbf{731} & \\
\bottomrule
\end{longtable}

\section*{3\quad Plasma sheath theory (Bohm criterion, presheath) \textnormal{(\normalsize 37 references)}}
This body of work establishes the electrostatic structure of the plasma-wall transition, from the quasi-neutral presheath to the non-neutral sheath, and the Bohm criterion that fixes the ion flux entering it. Collisional, kinetic and finite-Debye-length refinements extend the classical picture to the realistic regimes encountered in thrusters, etch reactors and fusion edges. These models are the physical priors that regularize the ion-energy-flux reconstruction.
\begin{enumerate}[leftmargin=1.4em]
\item \textbf{Kinetic theory of the presheath and the Bohm criterion} (2011). S. D. Baalrud, C. C. Hegna. Plasma Sources Science and Technology. \url{https://doi.org/10.1088/0963-0252/20/2/025013}
\item \textbf{Reply to comment on 'Kinetic theory of the presheath and the Bohm criterion'} (2012). S. D. Baalrud, C. C. Hegna. Plasma Sources Science and Technology. \url{https://doi.org/10.1088/0963-0252/21/6/068002}
\item \textbf{Comment on 'Kinetic theory of the presheath and the Bohm criterion'} (2012). K.-U. Riemann. Plasma Sources Science and Technology. \url{https://doi.org/10.1088/0963-0252/21/6/068001}
\item \textbf{Extensions and Applications of the Bohm Criterion} (2014). Scott D. Baalrud, Brett Scheiner, Benjamin Yee, Matthew Hopkins, Edward Barnat. arXiv:1412.4826. \url{https://arxiv.org/abs/1412.4826}
\item \textbf{Bohm criterion of plasma sheaths away from asymptotic limit} (2022). Yuzhi Li, Bhuvana Srinivasan, Yanzeng Zhang, Xian-Zhu Tang. arXiv:2201.11191. \url{https://arxiv.org/abs/2201.11191}
\item \textbf{Theory of the Electron Sheath and Presheath} (2015). Brett Scheiner, Scott D. Baalrud, Benjamin T. Yee, Matthew M. Hopkins, Edward V. Barnat. arXiv:1510.03490. \url{https://arxiv.org/abs/1510.03490}
\item \textbf{On the electron sheath theory and its applications in plasma-surface interaction} (2021). Guang-Yu Sun, Zhang Shu, An-Bang Sun, Guan-Jun Zhang. arXiv:2107.10074. \url{https://arxiv.org/abs/2107.10074}
\item \textbf{On the importance of slow ions in the kinetic Bohm criterion} (2024). Alessandro Geraldini, Stephan Brunner. arXiv:2404.17655. \url{https://arxiv.org/abs/2404.17655}
\item \textbf{The exact form of the Bohm criterion for a collisional plasma} (2016). Tsanko Vaskov Tsankov, Uwe Czarnetzki. arXiv:1603.00515. \url{https://arxiv.org/abs/1603.00515}
\item \textbf{The Kinetic and Hydrodynamic Bohm Criterions for Plasma Sheath Formation} (2021). Masahiro Suzuki, Masahiro Takayama. arXiv:2106.15990. \url{https://arxiv.org/abs/2106.15990}
\item \textbf{Continuum Kinetic Simulations of Plasma Sheaths and Instabilities} (2018). Petr Cagas. arXiv:1809.06368. \url{https://arxiv.org/abs/1809.06368}
\item \textbf{How to Patch Active Plasma and Collisionless Sheath: Practical Guide} (2002). Igor D. Kaganovich. arXiv:physics/0208041. \url{https://arxiv.org/abs/physics/0208041}
\item \textbf{Plasma-wall transition and sheath formation} (2021). Giovanni Manfredi, Fabrice Valsaque. arXiv:2108.01889. \url{https://arxiv.org/abs/2108.01889}
\item \textbf{Presheath-like structures and effusive particle losses for biased probes at and near I-V electron saturation} (2023). Brett Scheiner. arXiv:2312.17465. \url{https://arxiv.org/abs/2312.17465}
\item \textbf{Why is there no generally accepted theory of the plasma-sheath transition and the Bohm criterion?} (2025). K.-U. Riemann. Physics of Plasmas. \url{https://pubs.aip.org/aip/pop/article/32/3/030501/3338398}
\item \textbf{The Bohm Plasma-Sheath Model and the Bohm Criterion Revisited} (2007). N. Sternberg, V. Godyak. IEEE Transactions on Plasma Science. \url{https://doi.org/10.1109/tps.2007.905944}
\item \textbf{The Bohm–Chodura plasma sheath criterion} (1995). P. C. Stangeby. Physics of Plasmas. \url{https://doi.org/10.1063/1.871483}
\item \textbf{Bohm Criterion in a Magnetized Plasma Sheath} (2011). Y. Ji, X. Zou, J. Liu, X. Qiu. Plasma Science and Technology. \url{https://doi.org/10.1088/1009-0630/13/5/02}
\item \textbf{Plasma-sheath transition and Bohm criterion} (1992). K.-U. Riemann. Contributions to Plasma Physics. \url{https://doi.org/10.1002/ctpp.2150320309}
\item \textbf{A note on the plasma sheath and the Bohm criterion} (2007). G. D. Severn. American Journal of Physics. \url{https://doi.org/10.1119/1.2372469}
\item \textbf{On the upper bound in the Bohm sheath criterion} (2016). A. D. Kotelnikov, D. I. Skovorodin. Plasma Physics Reports. \url{https://doi.org/10.1134/s1063780x16020045}
\item \textbf{The Bohm criterion and the field singularity at the sheath edge} (1989). K.-U. Riemann. Physics of Fluids B: Plasma Physics. \url{https://doi.org/10.1063/1.859020}
\item \textbf{Two-Fluid Theory of Presheath and Space Charge Sheath in Low Pressure Plasmas} (1991). F. Valentini. Contributions to Plasma Physics. \url{https://doi.org/10.1002/ctpp.2150310206}
\item \textbf{Sheath formation in low-pressure discharges, the Bohm criterion and the consequences of collisions} (2013). F. Valentini, D. Kaiser. Plasma Sources Science and Technology. \url{https://doi.org/10.1088/0963-0252/23/1/015004}
\item \textbf{Plasma, presheath, collisional sheath and collisionless sheath potential profiles in weakly ionized, weakly collisional plasma} (2005). L. Oksuz, N. Hershkowitz. Plasma Sources Science and Technology. \url{https://doi.org/10.1088/0963-0252/14/1/022}
\item \textbf{Sheath and presheath in plasma with warm ions} (2009). S. Robertson. Physics of Plasmas. \url{https://doi.org/10.1063/1.3247874}
\item \textbf{The Bohm-sheath criterion in plasmas containing electrons and multiply charged ions} (2012). M. M. Khoramabadi, H. Ghomi, P. K. Shukla. Journal of Plasma Physics. \url{https://doi.org/10.1017/s0022377812000955}
\item \textbf{The Bohm criterion for the dusty plasma sheath} (2003). Y. Liu, D. Wang, X. Wang, Y. Zhang, X. Zou, C. Zhang. Physics of Plasmas. \url{https://doi.org/10.1063/1.1600734}
\item \textbf{Presheath-sheath coupling for kinetic trajectory simulation of a magnetized plasma sheath} (2019). R. P. Adhikari, S. Basnet, R. Lamichhane, R. Khanal. AIP Advances. \url{https://doi.org/10.1063/1.5088107}
\item \textbf{Structure of presheath-sheath in magnetized electronegative plasma} (2009). K. Yasserian, M. Aslaninejad, M. Ghoranneviss. Physics of Plasmas. \url{https://doi.org/10.1063/1.3077304}
\item \textbf{Magnetized electronegative plasma presheath-sheath containing two positive ion species} (2012). M. Aslaninejad, K. Yasserian. Physics of Plasmas. \url{https://doi.org/10.1063/1.3690101}
\item \textbf{Plasma sheath and presheath development near a partially reflective surface} (2021). B. N. Breizman, D. I. Kiramov. Journal of Plasma Physics. \url{https://doi.org/10.1017/s0022377820001518}
\item \textbf{Asymptotically matched sheath and presheath in a dense plasma} (1998). J. P. Thornhill, S. Desai. Journal of Plasma Physics. \url{https://doi.org/10.1017/s0022377898006497}
\item \textbf{Sheath and presheath in a collisionless open-field plasma} (1989). N. Sato, G. Miyawaki, Y. Fukui. Physics of Fluids B: Plasma Physics. \url{https://doi.org/10.1063/1.859137}
\item \textbf{Generalized Bohm sheath criterion in dusty plasma of space environment} (2021). S. Davood Sadatian, R. Gharjeh ghiyaei. Advances in Space Research. \url{https://doi.org/10.1016/j.asr.2021.06.020}
\item \textbf{Experimental Studies of the Bohm Criterion in a Two-Ion-Species Plasma Using Laser-Induced Fluorescence} (2003). D. Lee, N. Hershkowitz, G. D. Severn. Physical Review Letters. \url{https://doi.org/10.1103/PhysRevLett.90.145001}
\item \textbf{Laser-induced fluorescence measurements of argon and xenon ion velocities near the sheath boundary in 3 ion species plasmas} (2016). G. D. Severn, others. Physics of Plasmas. \url{https://pubs.aip.org/aip/pop/article/23/5/050703/955804}
\end{enumerate}

\section*{4\quad RF and capacitive sheaths \textnormal{(\normalsize 36 references)}}
Radio-frequency and capacitively-coupled sheaths oscillate on nanosecond timescales, modulating the ion energy delivered to surfaces. The literature characterizes sheath dynamics, harmonic structure and the resulting bimodal ion energy distributions. This motivates the time-resolved acquisition central to the proposed framework.
\begin{enumerate}[leftmargin=1.4em]
\item \textbf{On the role of secondary electron emission in capacitively coupled radio frequency plasma sheath: a theoretical ground} (2019). Guang-Yu Sun, Han-Wei Li, An-Bang Sun, Yuan Li, Bai-Peng Song, Hai-Bao Mu, Xiao-Ran Li, Guan-Jun Zhang. arXiv:1905.11688. \url{http://arxiv.org/abs/1905.11688}
\item \textbf{Kinetic Simulation of He radio frequency capacitive coupled plasma} (2023). M. Shihab, A. Elbadawy, M. S. Afify, N. El-Siragy. arXiv:2309.16368. \url{http://arxiv.org/abs/2309.16368}
\item \textbf{Kinetic simulation of the sheath dynamics in the intermediate radio-frequency regime} (2013). M. Shihab, A. T. Elgendy, I. Korolov, A. Derzsi, J. Schulze, D. Eremin, T. Mussenbrock, Z. Donko, R. P. Brinkmann. arXiv:1305.6786. \url{http://arxiv.org/abs/1305.6786}
\item \textbf{The effect of the driving frequency on the confinement of beam electrons and plasma density in low pressure capacitive discharges} (2014). S. Wilczek, J. Trieschmann, J. Schulze, E. Schuengel, R. P. Brinkmann, A. Derzsi, I. Korolov, Z. Donko, T. Mussenbrock. arXiv:1410.5327. \url{http://arxiv.org/abs/1410.5327}
\item \textbf{Implicit and electrostatic Particle-in-cell/Monte Carlo model in two dimensional and axisymmetric geometry II: Self-bias voltage effects in capacitively coupled plasmas} (2010). Wei Jiang, Hong-yu Wang, Zhen-hua Bi, You-nian Wang. arXiv:1001.2974. \url{http://arxiv.org/abs/1001.2974}
\item \textbf{High frequency sheath modulation and higher harmonic generation in a low pressure very high frequency capacitively coupled plasma excited by sawtooth waveform} (2020). Sarveshwar Sharma, Nishant Sirse, Miles M. Turner. arXiv:2007.09624. \url{http://arxiv.org/abs/2007.09624}
\item \textbf{Unveiling the role of dielectric trap states on capacitively coupled radio-frequency plasma discharge: dynamic charging behaviors} (2020). Shu Zhang, Guang-Yu Sun, Volcokas Arnas, Guan-Jun Zhang, An-Bang Sun. arXiv:2009.07097. \url{http://arxiv.org/abs/2009.07097}
\item \textbf{Ionization by bulk heating of electrons in capacitive radio frequency atmospheric pressure microplasmas} (2012). T. Hemke, D. Eremin, T. Mussenbrock, A. Derzsi, Z. Donko, K. Dittmann, J. Meichsner, J. Schulze. arXiv:1208.6519. \url{http://arxiv.org/abs/1208.6519}
\item \textbf{Kinetic Interpretation of Resonance Phenomena in Low Pressure Capacitively Coupled Radio Frequency Plasmas} (2015). S. Wilczek, J. Trieschmann, D. Eremin, R. P. Brinkmann, J. Schulze, E. Schuengel, A. Derzsi, I. Korolov, P. Hartmann, Z. Donko, T. Mussenbrock. arXiv:1507.05505. \url{http://arxiv.org/abs/1507.05505}
\item \textbf{On the self-excitation mechanisms of Plasma Series Resonance oscillations in single- and multi-frequency capacitive discharges} (2016). Edmund Schuengel, Steven Brandt, Ihor Korolov, Aranka Derzsi, Zoltan Donko, Julian Schulze. arXiv:1602.03482. \url{http://arxiv.org/abs/1602.03482}
\item \textbf{Intense boundary emission destroys normal radio-frequency plasma sheath} (2019). Guang-Yu Sun, An-Bang Sun, Guan-Jun Zhang. arXiv:1909.00127. \url{http://arxiv.org/abs/1909.00127}
\item \textbf{Plasma sheath physics: A circuital description, amelioration, and application} (2025). . arXiv:2506.03935. \url{https://arxiv.org/abs/2506.03935}
\item \textbf{An electric circuital analysis of laboratory plasma sheath fluctuations and propagations} (2025). . arXiv:2505.19553. \url{https://arxiv.org/abs/2505.19553}
\item \textbf{A Lagrangian model for laser-induced fluorescence and its application to measurements of plasma ion temperature and electrostatic waves} (2017). . arXiv:1704.00242. \url{https://arxiv.org/abs/1704.00242}
\item \textbf{Anomalous Capacitive Sheath with Deep Radio-Frequency Electric-Field Penetration} (2002). Igor D. Kaganovich. Physical Review Letters. \url{https://doi.org/10.1103/physrevlett.89.265006}
\item \textbf{Dynamics of a collisional, capacitive RF sheath} (1989). M. A. Lieberman. IEEE Transactions on Plasma Science. \url{https://doi.org/10.1109/27.24645}
\item \textbf{Nonlinear plasma dynamics in capacitive radio frequency discharges} (2007). Thomas Mussenbrock, Ralf Peter Brinkmann. Plasma Sources Science and Technology. \url{https://doi.org/10.1088/0963-0252/16/2/022}
\item \textbf{Ion energy distributions in rf sheaths; review, analysis and simulation} (1999). E. Kawamura, V. Vahedi, M. A. Lieberman, C. K. Birdsall. Plasma Sources Science and Technology. \url{https://doi.org/10.1088/0963-0252/8/3/202}
\item \textbf{Capacitive and resistive double sheath model in an asymmetric radio frequency plasma discharge} (2015). E. Faudot. Physics of Plasmas. \url{https://doi.org/10.1063/1.4928417}
\item \textbf{Gamma-electron dynamics and radio frequency capacitive sheath diagnostics} (1999). K. E. Orlov, A. S. Smirnov. IEEE Transactions on Plasma Science. \url{https://doi.org/10.1109/27.799811}
\item \textbf{Radio Frequency Capacitive Sheath Diagnostics by Time Resolved Measurements of Fast Electrons}. A. S. Smirnov, K. E. Orlov. NATO Science Series B: Electron Kinetics and Applications of Glow Discharges. \url{https://doi.org/10.1007/0-306-47076-4_32}
\item \textbf{Excitation Mechanisms and Sheath Dynamics in Capacitively Coupled Radio-Frequency Oxygen Plasmas} (2010). K. Dittmann, K. Matyash. Contributions to Plasma Physics. \url{https://doi.org/10.1002/ctpp.201000038}
\item \textbf{Sheath dynamics in radio-frequency atmospheric glow discharges} (2005). J. J. Shi, M. G. Kong. IEEE Transactions on Plasma Science. \url{https://doi.org/10.1109/tps.2005.844999}
\item \textbf{Kinetic effects on the radio-frequency sheath dynamics} (2004). Nong Xiang, F. Waelbroeck. IEEE International Conference on Plasma Science (ICOPS 2004). \url{https://doi.org/10.1109/plasma.2004.1339846}
\item \textbf{Sheath reversal during transient radio-frequency bias} (2002). P. R. J. Barroy, A. Goodyear. IEEE Transactions on Plasma Science. \url{https://doi.org/10.1109/tps.2002.1003967}
\item \textbf{Modeling a collisional, capacitive sheath for surface modification applications in radio-frequency discharges} (1992). Karla Bornig. Applied Physics Letters. \url{https://doi.org/10.1063/1.107248}
\item \textbf{Near-electrode sheath dynamics, current anharmonicity and battery effect in asymmetrical, low-pressure radio frequency discharges} (1992). Yu. P. Raizer, M. N. Shneider. Plasma Sources Science and Technology. \url{https://doi.org/10.1088/0963-0252/1/2/005}
\item \textbf{Radio-frequency sheath-plasma interactions with magnetic field tangency points along the sheath surface} (2013). H. Kohno, J. R. Myra. Physics of Plasmas. \url{https://doi.org/10.1063/1.4818991}
\item \textbf{Magnetized radio-frequency sheath in low temperature plasma} (2022). Rupali Sahu, Albina Tropina. AIAA SCITECH 2022 Forum. \url{https://doi.org/10.2514/6.2022-0981}
\item \textbf{Characteristics of Collision, Capacitive Radio Frequency Sheath} (2005). Zhang Yu, Ding Wanyu, Wang Wenchun. Plasma Science and Technology. \url{https://doi.org/10.1088/1009-0630/7/2/012}
\item \textbf{Collisional Sheath in the Electronegative Radio-Frequency Plasma} (2007). Gan Baoxia, Deng Wenjuan. Plasma Science and Technology. \url{https://doi.org/10.1088/1009-0630/9/4/03}
\item \textbf{Effects of Low-Frequency Source on a Dual-Frequency Capacitive Sheath near a Concave Electrode} (2014). Meilan Hao, Zhongling Dai. Plasma Science and Technology. \url{https://doi.org/10.1088/1009-0630/16/4/04}
\item \textbf{Dynamics of dust particles in a collisional radio-frequency plasma sheath} (2014). Morad El Kaouini, Hassan Chatei. Physica Scripta. \url{https://doi.org/10.1088/0031-8949/2014/t161/014052}
\item \textbf{Hydrodynamic radio-frequency model of an ion sheath near a conductor in a plasma} (1991). Gilbert A. Morin, Keith G. Balmain. Radio Science. \url{https://doi.org/10.1029/90rs02757}
\item \textbf{Laser-induced fluorescence measurement of the dynamics of a pulsed planar sheath} (1994). . Physics of Plasmas. \url{https://pubs.aip.org/aip/pop/article/1/4/1064/445311}
\item \textbf{Revisiting the anomalous rf field penetration into a warm plasma} (2005). Igor D. Kaganovich, Oleg V. Polomarov, Constantine E. Theodosiou. arXiv:physics/0506135. \url{http://arxiv.org/abs/physics/0506135}
\end{enumerate}

\section*{5\quad Ion energy flux and power balance at surfaces \textnormal{(\normalsize 33 references)}}
These studies quantify the energy carried to bounding surfaces by ions and its role in the overall power balance of bounded plasmas. They connect sheath potential, ion flux and secondary processes to a deposited power density. They define precisely the observable the project seeks to reconstruct non-intrusively.
\begin{enumerate}[leftmargin=1.4em]
\item \textbf{Model of etching profiles for ion energy flux dependent etch rates in a collisionless plasma sheath} (1995). Barbara Abraham-Shrauner, Chungdar Daniel Wang. Journal of Applied Physics. \url{https://doi.org/10.1063/1.358637}
\item \textbf{Development of an ion energy flux sensor for plasma processing} (2001). K. C. Leou, S. C. Lai, C. H. Cheng, C. Lin, C. H. Tsai. IEEE Conference Record - PPPS-2001. \url{https://doi.org/10.1109/ppps.2001.961254}
\item \textbf{Ion flux and energy distributions at electrode surfaces in LAPPS} (2001). S. G. Walton, D. Leonhardt, D. D. Blackwell, D. P. Murphy, R. F. Fernsler, A. R. A. Meger. IEEE Conference Record - PPPS-2001. \url{https://doi.org/10.1109/ppps.2001.961107}
\item \textbf{The effect of ion–neutral collisions on the weakly collisional plasma-sheath and the reduction of the ion flux to the wall} (2004). Zoltan Sternovsky. Plasma Sources Science and Technology. \url{https://doi.org/10.1088/0963-0252/14/1/004}
\item \textbf{Direct measurements of sheath-accelerated secondary electrons for monitoring the incident ion flux in plasma immersion ion implantation} (2002). Keiji Nakamura, Mitsuaki Tanaka, Hideo Sugai. Plasma Sources Science and Technology. \url{https://doi.org/10.1088/0963-0252/11/2/306}
\item \textbf{Energy measurements of sheath-accelerated secondary electrons in plasma immersion ion implantation} (2003). Keiji Nakamura, Mitsuaki Tanaka, Hideo Sugai. Surface and Coatings Technology. \url{https://doi.org/10.1016/s0257-8972(03)00087-2}
\item \textbf{Plasma Ion Bombardment Induced Heat Flux on the Wafer Surface in Inductively Coupled Plasma Reactive Ion Etch} (2023). Sung Il Cho, Hyun Keun Park, Surin An, Sang Jeen Hong. Applied Sciences. \url{https://doi.org/10.3390/app13179533}
\item \textbf{Sheath potential difference studies in low-pressure high-frequency capacitive RF plasma as a fundamental research for ion energy impinging on the material surface} (2006). Zhi-Gang Li, YuCheng Ding. Surface and Coatings Technology. \url{https://doi.org/10.1016/j.surfcoat.2005.08.001}
\item \textbf{Sheath dynamics in plasma immersion ion implantation} (1996). R. Günzel, J. Brutscher. Surface and Coatings Technology. \url{https://doi.org/10.1016/0257-8972(96)02883-6}
\item \textbf{Fluid Model of Plasma Sheath Involving Ion Energy Spectrum} (2010). Josef Blazek, Petr Bartos, Petr Spatenka. IEEE Transactions on Plasma Science. \url{https://doi.org/10.1109/tps.2010.2055894}
\item \textbf{Ion dynamics of pulsed plasma source ion implantation in the sheath of a hemispherical bowl-shaped target} (2003). Chengsen Liu, Dezhen Wang. Surface and Coatings Technology. \url{https://doi.org/10.1016/s0257-8972(03)00251-2}
\item \textbf{Variations of Plasma Parameters and Sheath Thickness Due to the Extraction of Ion and Electron Flux in a Double Plasma Device} (2007). M. K. Mishra, A. Phukan, M. Chakraborty. Contributions to Plasma Physics. \url{https://doi.org/10.1002/ctpp.200710080}
\item \textbf{Effect of Ion-Neutral Collisions on the Plasma Sheath in the Presence of Surface-Produced Negative Ions} (2024). Sutapa Samanta, Rakesh Moulick, P. J. Bhuyan, B. J. Saikia. Contributions to Plasma Physics. \url{https://doi.org/10.1002/ctpp.202400088}
\item \textbf{The negative ion flux across a double sheath at the formation of a virtual cathode} (2010). R. McAdams, M. Bacal. Plasma Sources Science and Technology. \url{https://doi.org/10.1088/0963-0252/19/4/042001}
\item \textbf{The Study of Kinetic Energy of Ion and Sheath Thickness in Magnetized Plasma Sheath} (2015). Roshan Chalise, Raju Khanal. Journal of Materials Science and Engineering A. \url{https://doi.org/10.17265/2161-6213/2015.1-2.005}
\item \textbf{Plasma-sheath ion nitriding: a cost-efficient technology for the surface modification of SS304} (1998). S. Mukherjee, P. I. John. Surface and Coatings Technology. \url{https://doi.org/10.1016/s0257-8972(97)00372-1}
\item \textbf{Sheath Formation Criterion and Ion Flux for Non-Maxwellian Plasma} (1997). Hiroshi Amemiya. Journal of the Physical Society of Japan. \url{https://doi.org/10.1143/jpsj.66.1335}
\item \textbf{Positive ion speed at the plasma-sheath boundary of a negative ion-emitting electrode} (2020). Avnish K. Pandey, Shantanu K. Karkari. Contributions to Plasma Physics. \url{https://doi.org/10.1002/ctpp.201900116}
\item \textbf{Sheath structure and ion current onto a patterned surface immersed in a plasma at floating potential} (2002). Boris Briehl, Herbert M. Urbassek. Surface and Coatings Technology. \url{https://doi.org/10.1016/s0257-8972(02)00410-3}
\item \textbf{Power balance in a fusion reactor} (2007). . Plasma Physics and Fusion Energy (Cambridge University Press). \url{https://doi.org/10.1017/cbo9780511755705.006}
\item \textbf{Plasma-Sheath-Surface Dynamics} (1990). Charles K. Birdsall. Technical report (DTIC). \url{https://doi.org/10.21236/ada226086}
\item \textbf{Modeling and simulation of ion energy distribution functions in technological plasmas} (2011). Thomas Mussenbrock. arXiv:1109.4435. \url{https://arxiv.org/abs/1109.4435}
\item \textbf{Predicting PDF tails of flux in plasma sheath region} (2009). Johan Anderson, Eun-jin Kim. arXiv:0906.0481. \url{https://arxiv.org/abs/0906.0481}
\item \textbf{Enhanced entropy production in heat-flux-driven plasma sheath} (2016). Yuji Ohno, Zensho Yoshida. arXiv:1612.08836. \url{https://arxiv.org/abs/1612.08836}
\item \textbf{Magnetized Plasma Sheath in the Presence of Negative Ions} (2020). R. Paul, S. Adhikari, R. Moulick, S. S. Kausik, B. K. Saikia. arXiv:2002.05445. \url{https://arxiv.org/abs/2002.05445}
\item \textbf{Effect of two temperature electrons in a collisional magnetized plasma sheath} (2019). G. Sharma, S. Adhikari, R. Moulick, S. S. Kausik, B. K. Saikia. arXiv:1907.10280. \url{https://arxiv.org/abs/1907.10280}
\item \textbf{Influence of ionisation zone motion in high power impulse magnetron sputtering on angular ion flux and NbOx film growth} (2015). Robert Franz, César Clavero, Jonathan Kolbeck, André Anders. arXiv:1507.04498. \url{https://arxiv.org/abs/1507.04498}
\item \textbf{Data-driven discovery of a heat flux closure for electrostatic plasma phenomena} (2024). Emil R. Ingelsten, Madox C. McGrae-Menge, E. Paulo Alves, Istvan Pusztai. arXiv:2411.18358. \url{https://arxiv.org/abs/2411.18358}
\item \textbf{Interaction of Biased Electrodes and Plasmas: Sheaths, Double Layers and Fireballs} (2019). . arXiv:1911.06424. \url{https://arxiv.org/abs/1911.06424}
\item \textbf{Development of a Noninvasive Real-Time Ion Energy Distribution Monitoring System Applicable to Collisional Plasma Sheath} (2022). . Sensors (PMC). \url{https://www.ncbi.nlm.nih.gov/pmc/articles/PMC9413131/}
\item \textbf{Wafer Type Ion Energy Monitoring Sensor for Plasma Diagnosis} (2023). . Sensors (PMC). \url{https://www.ncbi.nlm.nih.gov/pmc/articles/PMC10007180/}
\item \textbf{Tailored ion energy distributions on plasma electrodes (Review Article)} (2022). . Journal of Vacuum Science \textbackslash{}\& Technology. \url{https://www.chee.uh.edu/sites/chbe/files/faculty/economou/tailored-ieds-60th-jvst.pdf}
\item \textbf{Ion energy distributions driven by the ultra-low frequency electrical asymmetry effect in biased inductively coupled plasmas} (2025). . Physics of Plasmas. \url{https://pubs.aip.org/aip/pop/article/32/8/083503/3357931/Ion-energy-distributions-driven-by-the-ultra-low}
\end{enumerate}

\section*{6\quad Ion energy and velocity distributions at surfaces \textnormal{(\normalsize 27 references)}}
Measurements and models of ion energy and velocity distribution functions at boundaries reveal the non-Maxwellian, drifting character of the ion population. The distribution, not merely a mean energy, determines erosion and heating. Recovering it (via LIF) is the key input to the energy-flux integral.
\begin{enumerate}[leftmargin=1.4em]
\item \textbf{Reconstruction of energy and angle distribution function of surface-emitted negative ions in hydrogen plasmas using mass spectrometry} (2019). D. Kogut, K. Achkasov, J. Dubois, R. Moussaoui, J. B. Faure, J. Layet, A. Simonin, G. Cartry. arXiv:1910.05157. \url{http://arxiv.org/abs/1910.05157v1}
\item \textbf{Ion energy distributions from the impact of an atmospheric dielectric barrier discharge plasma jet on surfaces} (2025). Daniel Henze, Laura Chauvet, Achim von Keudell. arXiv:2505.07771. \url{http://arxiv.org/abs/2505.07771v1}
\item \textbf{Driven dust-charge fluctuation and chaotic ion dynamics in the plasma sheath and pre-sheath regions} (2022). Mridusmita Das, Suniti Changmai, Madhurjya P. Bora. arXiv:2212.02813. \url{http://arxiv.org/abs/2212.02813v4}
\item \textbf{Planar Laser-Induced Fluorescence system for Space and Phase-resolved Ion Velocity Distribution Function Measurements} (2025). . arXiv:2504.02690. \url{https://arxiv.org/abs/2504.02690}
\item \textbf{Bayesian inference of ion velocity distribution function from laser-induced fluorescence spectra} (2021). . Scientific Reports. \url{https://www.nature.com/articles/s41598-021-00138-4}
\item \textbf{Temporally resolved ion velocity distribution measurements in a radio-frequency plasma sheath} (2011). . Physics of Plasmas. \url{https://pubs.aip.org/aip/pop/article-abstract/18/5/053503/811883/Temporally-resolved-ion-velocity-distribution}
\item \textbf{Data processing techniques for ion and electron-energy distribution functions} (2023). . Physics of Plasmas. \url{https://pubs.aip.org/aip/pop/article/30/4/040501/2882934/Data-processing-techniques-for-ion-and-electron}
\item \textbf{Ion-Energy Diagnostics in the Plasma Exhaust Plume of a Hall Thruster}. . Journal of Propulsion and Power. \url{https://arc.aiaa.org/doi/pdf/10.2514/2.5660}
\item \textbf{Ion thrusters for electric propulsion: Scientific issues developing a niche technology into a game changer} (2020). . Review of Scientific Instruments. \url{https://pubs.aip.org/aip/rsi/article/91/6/061101/957865/Ion-thrusters-for-electric-propulsion-Scientific}
\item \textbf{Origins of ion energy distribution function (IEDF) in high power impulse magnetron sputtering (HIPIMS) plasma discharge} (2008). Hecimovic, Burcalova, Ehiasarian. Journal of Physics D: Applied Physics. \url{https://doi.org/10.1088/0022-3727/41/9/095203}
\item \textbf{Langmuir Probe in Non-Maxwellian Plasma: Correlation Between the Electron Energy Distribution Function, the Ion Energy at the Sheath Edge and the Floating Voltage} (2014). Jauberteau, Jauberteau. Contributions to Plasma Physics. \url{https://doi.org/10.1002/ctpp.201400004}
\item \textbf{Measurement of parallel ion energy distribution function in PISCES plasma} (1987). Tynan, Goebel, Conn. Technical report (OSTI). \url{https://doi.org/10.2172/6268436}
\item \textbf{The Ion Velocity Distribution Function at the Plasma-Sheath Interface: New Boundary Conditions for the Fluid Model} (1992). Günther. Contributions to Plasma Physics. \url{https://doi.org/10.1002/ctpp.2150320317}
\item \textbf{Ion energy distribution in a radio frequency sheath of plasma with Kappa electron velocity distribution} (2020). Ou, Huang. Contributions to Plasma Physics. \url{https://doi.org/10.1002/ctpp.202000108}
\item \textbf{Evolution of the ion distribution function for an expanding sheath with collisions} (1997). Sheridan. Plasma Sources Science and Technology. \url{https://doi.org/10.1088/0963-0252/6/1/013}
\item \textbf{The effects of electron energy distribution function on the plasma sheath structure in the presence of charged nanoparticles} (2020). Khalilpour, Foroutan. Journal of Plasma Physics. \url{https://doi.org/10.1017/s0022377820000161}
\item \textbf{Plasma sheath model and ion energy distribution for all radio frequencies} (1999). Panagopoulos, Economou. Journal of Applied Physics. \url{https://doi.org/10.1063/1.369701}
\item \textbf{End-boundary sheath potential, electron and ion energy distribution in the low-pressure non-ambipolar electron plasma} (2013). Chen, Chen, Funk. Plasma Sources Science and Technology. \url{https://doi.org/10.1088/0963-0252/22/6/065015}
\item \textbf{Advancing ion energy diagnostics: measurement of IEDF in cold atmospheric plasma jets using retarding field energy analyzer} (2025). S, Sharma, Guha, Rajan, Anand. Journal of Instrumentation. \url{https://doi.org/10.1088/1748-0221/20/07/p07036}
\item \textbf{Verifying effects of instability enhanced ion-ion Coulomb collisions on ion velocity distribution functions near the sheath edge in low temperature plasmas} (2014). Yip, Hershkowitz, Severn. Plasma Sources Science and Technology. \url{https://doi.org/10.1088/0963-0252/24/1/015018}
\item \textbf{Simulation of heavy ion energy distribution function in the cathode sheath of noble gas mixture glow discharge} (1999). Bondarenko, Kristya. Vacuum. \url{https://doi.org/10.1016/s0042-207x(98)00415-1}
\item \textbf{Ion energy distribution in plasma immersion ion implantation} (1997). Mändl, Brutscher, Günzel, Möller. Surface and Coatings Technology. \url{https://doi.org/10.1016/s0257-8972(97)00051-0}
\item \textbf{Effect of the electron energy distribution function in plasma on the Bohm criterion and on the drop voltage through the sheath: Case of microwave expanding plasma} (2009). Jauberteau, Jauberteau. Journal of Applied Physics. \url{https://doi.org/10.1063/1.3211910}
\item \textbf{Reconstruction of the time-averaged sheath potential profile in an argon radiofrequency plasma using the ion energy distribution} (1995). Fivaz, Brunner, Schwarzenbach, Howling, Hollenstein. Plasma Sources Science and Technology. \url{https://doi.org/10.1088/0963-0252/4/3/007}
\item \textbf{Retarding Field Analyzer for Ion Energy Distribution Measurement Through a Radio-Frequency or Pulsed Biased Sheath} (2009). Gahan, Dolinaj, Hayden, Hopkins. Plasma Processes and Polymers. \url{https://doi.org/10.1002/ppap.200931607}
\item \textbf{Ion energy distribution at an r.f.-biased electrode in an inductively coupled plasma affected by collisions in a sheath} (2000). Mizutani, Hayashi. Thin Solid Films. \url{https://doi.org/10.1016/s0040-6090(00)01148-2}
\item \textbf{Measurement of Ion Energy Distribution Function of Fast Plasma Flow Driven by Plasma Focus Device Using Retarding Field Energy Analyzer} (2023). Oguchi, Matsuyama, Shigeta, Takezaki, Sasaki, Ito. Plasma and Fusion Research. \url{https://doi.org/10.1585/pfr.18.2401049}
\end{enumerate}

\section*{7\quad Sputtering and ion-induced erosion \textnormal{(\normalsize 32 references)}}
This literature relates incident ion energy and angle to sputter yield and net surface erosion, including energy thresholds and material dependence. It converts an ion-energy-flux measurement into predictions of wear and lifetime. It underpins the commercial value proposition for propulsion and fusion hardware.
\begin{enumerate}[leftmargin=1.4em]
\item \textbf{Erosion yields of carbon under various plasma conditions in Pilot-PSI} (2012). K. Bystrov, J. Westerhout, M. Matveeva, A. Litnovsky, L. Marot, E. Zoethout, G. De Temmerman. arXiv:1210.3502. \url{http://arxiv.org/abs/1210.3502v1}
\item \textbf{White Light Interferometry for Quantitative Surface Characterization in Ion Sputtering Experiments} (2012). S. V. Baryshev, A. V. Zinovev, C. E. Tripa, R. A. Erck, I. V. Veryovkin. arXiv:1204.0808. \url{http://arxiv.org/abs/1204.0808v1}
\item \textbf{On the stabilization of ion sputtered surfaces} (2007). Benny Davidovitch, Michael J. Aziz, Michael P. Brenner. arXiv:0706.2662. \url{http://arxiv.org/abs/0706.2662v1}
\item \textbf{Effect of bombarding steel with Xe+ ions on the surface nanostructure and on pulsed plasma nitriding process} (2014). S. Cucatti, E. A. Ochoa, M. Morales, R. Droppa, J. Garcia, H. C. Pinto, L. F. Zagonel, D. Wisnivesky, C. A. Figueroa, F. Alvarez. arXiv:1411.7281. \url{http://arxiv.org/abs/1411.7281v1}
\item \textbf{Preparation of atomically clean and flat Si(100) surfaces by low-energy ion sputtering and low-temperature annealing} (2003). J. C. Kim, J.-Y. Ji, J. S. Kline, J. R. Tucker, T.-C. Shen. arXiv:cond-mat/0306397. \url{http://arxiv.org/abs/cond-mat/0306397v1}
\item \textbf{Cs evaporation in a negative ion source and Cs cleaning tests by plasma sputtering} (2023). M. Barbisan, R. S. Delogu, A. Pimazzoni, C. Poggi, M. Ugoletti, M. Cavenago. arXiv:2308.15328. \url{http://arxiv.org/abs/2308.15328v1}
\item \textbf{On the Competition Between the Phenomena Involved in the Aerosol Dynamics in Sputtering Nonequilibrium Plasma} (2016). Armelle Michau, Khaled Hassouni, Cécile Arnas, Guillaume Lombardi. arXiv:1702.01019. \url{http://arxiv.org/abs/1702.01019v1}
\item \textbf{Advances in neutral tungsten ultraviolet spectroscopy for the potential benefit to gross erosion diagnosis} (2019). . arXiv:1912.07723. \url{https://arxiv.org/abs/1912.07723}
\item \textbf{Sputtering yield measurements during low energy xenon plasma bombardment} (2003). R. P. Doerner, D. G. Whyte, D. M. Goebel. Journal of Applied Physics. \url{https://doi.org/10.1063/1.1566474}
\item \textbf{Sputtering by ion bombardment: theoretical concepts} (1981). P. Sigmund. Topics in Applied Physics: Sputtering by Particle Bombardment I. \url{https://doi.org/10.1007/3540105212_7}
\item \textbf{Physical Sputtering of Solids at Ion Bombardment} (1986). J. Roth. Physics of Plasma-Wall Interactions in Controlled Fusion. \url{https://doi.org/10.1007/978-1-4757-0067-1_8}
\item \textbf{Sputtering—Surface changes induced by ion bombardment} (1976). B. Navinšek. Progress in Surface Science. \url{https://doi.org/10.1016/0079-6816(76)90001-0}
\item \textbf{Heavy ion sputtering induced surface topography development} (1983). G. Carter, B. Navinšek, J. L. Whitton. Topics in Applied Physics: Sputtering by Particle Bombardment II. \url{https://doi.org/10.1007/3-540-12593-0_6}
\item \textbf{Examples of Ion Bombardment Effects On Film Growth and Erosion Processes-Plasma and Beam Experiments} (1986). E. Kay. Erosion and Growth of Solids Stimulated by Atom and Ion Beams. \url{https://doi.org/10.1007/978-94-009-4422-0_14}
\item \textbf{Sputtering erosion of stratified media and by time dependent ion bombardment} (1983). G. W. Nobes, G. Carter, M. J. Nobes-Lewis, J. L. Whitton. Vacuum. \url{https://doi.org/10.1016/0042-207x(83)90657-7}
\item \textbf{Non-Linear Effects of Molecular Ion Bombardment on Sputtering Yield} (1982). H. Karashima. Japanese Journal of Applied Physics. \url{https://doi.org/10.1143/jjap.21.434}
\item \textbf{Angular dependence of the sputtering yield of water ice by 100 keV proton bombardment} (2005). J. Vidal, B. D. Teolis, R. A. Baragiola. Surface Science. \url{https://doi.org/10.1016/j.susc.2005.05.007}
\item \textbf{Molecular Dynamic Simulation of the Effect of Initial Surface Temperature on Arc Erosion Due to Ion Bombardment} (2023). Wang, Matharage, Wang. Plasma Physics and Technology. \url{https://doi.org/10.14311/ppt.2023.3.146}
\item \textbf{Sputtering yields for low-energy Ar+- and Ne+-ion bombardment} (2013). . Vacuum. \url{https://www.sciencedirect.com/science/article/abs/pii/S0042207X12002035}
\item \textbf{Materials surface modification by plasma bombardment under simultaneous erosion and redeposition conditions} (1987). . Nuclear Instruments and Methods in Physics Research B. \url{https://www.sciencedirect.com/science/article/abs/pii/0168583X8790471X}
\item \textbf{Status analysis on sputtering and erosion evaluation methods of ion optic systems} (2024). . Chinese Journal of Aeronautics. \url{https://www.sciencedirect.com/science/article/pii/S1000936124003170}
\item \textbf{Review Article: Unraveling synergistic effects in plasma-surface processes by means of beam experiments} (2017). . Journal of Vacuum Science \textbackslash{}\& Technology A. \url{https://pmc.ncbi.nlm.nih.gov/articles/PMC5648578/}
\item \textbf{Evaluation of 4-point probe and eddy current methods as plasma facing components erosion measurement methods} (2024). . Nuclear Fusion. \url{https://iopscience.iop.org/article/10.1088/1741-4326/ae03d7}
\item \textbf{Material testing facilities and programs for plasma-facing component testing} (2017). . Nuclear Fusion. \url{https://iopscience.iop.org/article/10.1088/1741-4326/aa4feb}
\item \textbf{Gross and net erosion balance of plasma-facing materials in full-W tokamaks} (2021). . Nuclear Fusion. \url{https://iopscience.iop.org/article/10.1088/1741-4326/ac22d2}
\item \textbf{The erosion of selected tungsten coatings by ion beam and plasma sources compared to calculated predictions} (2023). . Fusion Engineering and Design. \url{https://www.sciencedirect.com/science/article/abs/pii/S0920379623003630}
\item \textbf{Measurements of channel erosion of Hall thrusters by laser-induced fluorescence} (2020). . Journal of Applied Physics. \url{https://pubs.aip.org/aip/jap/article/128/18/183301/157720/Measurements-of-channel-erosion-of-Hall-thrusters}
\item \textbf{Review of Plasma-Induced Hall Thruster Erosion} (2020). . Applied Sciences. \url{https://www.mdpi.com/2076-3417/10/11/3775}
\item \textbf{Three-Dimensional Model for Erosion of a Hall-Effect Thruster Discharge Channel Wall} (2014). . Journal of Propulsion and Power. \url{https://arc.aiaa.org/doi/10.2514/1.B35098}
\item \textbf{Computing the Wall Erosion Rate of a Low-Power Hall Thruster Using Particle-in-Cell Numerical Simulations} (2024). . Brazilian Journal of Physics. \url{https://link.springer.com/article/10.1007/s13538-024-01692-8}
\item \textbf{Assessments of surface composition and sputtering yield of Au-Cu alloys for Ar+ ion bombardment} (1984). Kang, Kawatoh, Shimizu. Surface Science. \url{https://doi.org/10.1016/0039-6028(84)90116-x}
\item \textbf{Sputtering and secondary ion yields of metals subjected to oxygen ion bombardment} (1982). Taga, Inoue, Satta. Surface Science. \url{https://doi.org/10.1016/0039-6028(82)90181-9}
\end{enumerate}

\section*{8\quad Plasma-material interaction \textnormal{(\normalsize 28 references)}}
Broader plasma-material interaction studies cover heat loading, implantation, redeposition and surface evolution under plasma exposure. They frame the boundary conditions and damage mechanisms the diagnostic informs. They connect the measured flux to engineering outcomes.
\begin{enumerate}[leftmargin=1.4em]
\item \textbf{CPIC: a Curvilinear Particle-In-Cell code for plasma-material interaction studies} (2013). Gian Luca Delzanno, Enrico Camporeale, J. David Moulton, Joseph E. Borovsky, Elizabeth A. MacDonald, Michelle F. Thomsen. arXiv:1311.2286. \url{http://arxiv.org/abs/1311.2286v1}
\item \textbf{Self-consistent generation of the ambipolar electric field in collisionless plasmas via multi-mode electrostatics} (2025). Luca Barbieri. arXiv:2509.03142. \url{http://arxiv.org/abs/2509.03142v1}
\item \textbf{Planar Laser Induced Fluorescence Mapping of a Carbon Laser Produced Plasma} (2022). . arXiv:2207.10843. \url{https://arxiv.org/abs/2207.10843}
\item \textbf{Attosecond dispersion as a diagnostics tool for solid-density laser-generated plasmas} (2022). Andreas Sundstrom, Istvan Pusztai, Per Eng-Johnsson, Tunde Fulop. arXiv:2202.00406. \url{http://arxiv.org/abs/2202.00406v3}
\item \textbf{Developments of high heat flux components and plasma surface interaction studies in Japan} (1996). T. Hino, T. Yamashina. Vacuum. \url{https://doi.org/10.1016/0042-207x(96)00097-8}
\item \textbf{Evaluation of material erosion from plasma-facing surfaces in hard disruptions via simulated ablation due to heat flux in electrothermal discharges} (2013). J. R. Echols, A. L. Winfrey. 2013 IEEE 25th Symposium on Fusion Engineering (SOFE). \url{https://doi.org/10.1109/sofe.2013.6635503}
\item \textbf{Material erosion and surface temperature response to plasma-disruption heat load simulations} (1995). E. V. van Osch, J. G. van der Laan. Journal of Nuclear Materials. \url{https://doi.org/10.1016/0022-3115(94)00584-2}
\item \textbf{Mesoscale Erosion of Tungsten Surface due to Low-Energy High-Flux Helium/Deuterium Plasma Irradiation} (2005). Shuichi Takamura. Journal of Plasma and Fusion Research. \url{https://doi.org/10.1585/jspf.81.25}
\item \textbf{Nonlinear interactions between high heat flux plasma and electron-emissive hot material surface} (1996). M. Y. Ye, S. Masuzaki, K. Shiraishi, S. Takamura, N. Ohno. Physics of Plasmas. \url{https://doi.org/10.1063/1.871854}
\item \textbf{A Heat Flux Microsensor for Direct Measurements in Plasma Surface Interactions} (2011). Remi Dussart, Anne-Lise Thomann, Nadjib Semmar. Microsensors (InTech). \url{https://doi.org/10.5772/16862}
\item \textbf{Physical modeling of plasma material interaction and first numerical results of vapor shield and material erosion during tokamak plasma disruptions} (1993). B. Goel, W. Hobel, G. Piazza, H. Wurz, I. K. Konkashbaev, I. Landman. Fusion Technology 1992. \url{https://doi.org/10.1016/b978-0-444-89995-8.50046-1}
\item \textbf{Surface erosion from plasma materials interaction} (1979). R. Behrisch. Journal of Nuclear Materials. \url{https://doi.org/10.1016/0022-3115(79)90399-4}
\item \textbf{Heat flux near critical density in laser-plasma interaction} (1986). J. F. Luciani, P. Mora. Physics Letters A. \url{https://doi.org/10.1016/0375-9601(86)90140-4}
\item \textbf{Erosion behavior and surface characterization of beryllium under high-flux deuterium plasma bombardment} (1997). J. Won. Journal of Nuclear Materials. \url{https://doi.org/10.1016/s0022-3115(96)00682-4}
\item \textbf{Heat transfer with nonuniform surface heat flux and thermal-hydraulic design of the plasma-facing components of fusion reactors} (1991). Tomoaki Kunugi, Mohammad Z. Hasan. Fusion Technology 1990. \url{https://doi.org/10.1016/b978-0-444-88508-1.50180-8}
\item \textbf{Parametric studies of high-heat-flux-induced erosion of surfaces in an electrothermal accelerator} (1989). O. Hankins, M. Bourham, O. Auciello, J. Stock, J. Gilligan, B. Wehring. IEEE 1989 International Conference on Plasma Science. \url{https://doi.org/10.1109/plasma.1989.165969}
\item \textbf{Studies of material surface erosion exposed to plasmas in an electrothermal launcher} (1990). O. E. Hankins, M. A. Bourham, J. G. Gilligan, S. Tallavarjula. 1990 Plasma Science IEEE Conference Record. \url{https://doi.org/10.1109/plasma.1990.110591}
\item \textbf{Surface erosion studies in a plasma-propellant interaction experiment} (1994). M. A. Bourham, J. G. Gilligan, C. M. Edwards, M. L. Nahm. 1994 IEEE 21st International Conference on Plasma Sciences (ICOPS). \url{https://doi.org/10.1109/plasma.1994.588996}
\item \textbf{Early Career: Harnessing nanotechnology for fusion plasma-material interface research in an in-situ particle-surface interaction facility} (2014). Jean Paul Allain. OSTI Technical Report. \url{https://doi.org/10.2172/1169163}
\item \textbf{Erosion and modifications of tungsten-coated carbon and copper under high heat flux} (2003). Liu Xiang, S. Tamura, K. Tokunaga, N. Yoshida, Zhang Fu, Xu Zeng-yu, Ge Chang-chun, N. Noda. Plasma Science and Technology. \url{https://doi.org/10.1088/1009-0630/5/4/008}
\item \textbf{Simulation study of heat flux deposition pattern on the surface of HT-7 toroidal limiters} (2008). . Plasma Science and Technology. \url{https://doi.org/10.1088/1009-0630/10/2/02}
\item \textbf{Interaction of two-phase plasma flux of vacuum arc discharge with a surface} (1998). S. A. Mouboiadjian. Surface and Coatings Technology. \url{https://doi.org/10.1016/s0257-8972(97)00559-8}
\item \textbf{A novel optical emission spectroscopy method for diagnostics of contribution of different ionization mechanisms and flux of ions in different valences in the discharge channel of a Hall thruster} (2024). . Chinese Journal of Aeronautics. \url{https://www.sciencedirect.com/science/article/pii/S1000936124000050}
\item \textbf{Comprehensive Screening of Plasma-Facing Materials for Nuclear Fusion} (2024). . PRX Energy. \url{https://link.aps.org/doi/10.1103/PRXEnergy.3.043002}
\item \textbf{Challenges for plasma-facing components in nuclear fusion} (2019). . Matter and Radiation at Extremes (AIP). \url{https://pubs.aip.org/aip/mre/article/4/5/056201/253043/Challenges-for-plasma-facing-components-in-nuclear}
\item \textbf{Thermal damage of tungsten-armored plasma-facing components under high heat flux loads} (2020). . Scientific Reports. \url{https://www.nature.com/articles/s41598-020-57852-8}
\item \textbf{Melt layer erosion during ELM-like heat loading on molybdenum as an alternative plasma-facing material} (2017). . Nuclear Fusion (PMC). \url{https://www.ncbi.nlm.nih.gov/pmc/articles/PMC5612961/}
\item \textbf{Laser Thomson scattering for low-temperature plasmas}. . ResearchGate (review). \url{https://www.researchgate.net/publication/230905900_Laser_Thomson_scattering_for_low-temperature_plasmas}
\end{enumerate}

\section*{9\quad Langmuir probe diagnostics \textnormal{(\normalsize 33 references)}}
Electrostatic probe theory and practice provide the classical route to electron temperature, density and plasma potential. The same literature documents the perturbation, erosion and interpretation limits that make probes unsuitable near high-flux sheaths. This establishes the gap that non-intrusive optical fusion addresses.
\begin{enumerate}[leftmargin=1.4em]
\item \textbf{Analysis of the temperature influence on Langmuir probe measurements on the basis of gyrofluid simulations} (2011). Felix P. Gennrich, Alexander Kendl. arXiv:1109.3600. \url{http://arxiv.org/abs/1109.3600v1}
\item \textbf{New procedure to estimate plasma parameters through the q-Weibull distribution by using a Langmuir probe in a cold plasma} (2022). F. J. Gonzalez, J. I. Gonzalez, S. Soler, C. E. Repetto, B. J. Gomez, D. B. Berdichevsky. arXiv:2202.11870. \url{http://arxiv.org/abs/2202.11870v1}
\item \textbf{Spatially resolved determination of the electronic density and temperature by a visible spectrotomography diagnostic in a linear magnetized plasma} (2020). V. Gonzalez-Fernandez, P. David, R. Baude, A. Escarguel, Y. Camenen. arXiv:2005.06353. \url{http://arxiv.org/abs/2005.06353v1}
\item \textbf{Design and Implementation of a Fast-Sweeping Langmuir Probe Diagnostic for DC Arc Jet Environments} (2025). Sebastian V. Colom, Magnus A. Haw, Jocelino Rodrigues. arXiv:2510.26162. \url{http://arxiv.org/abs/2510.26162v1}
\item \textbf{Influence of temperature fluctuations on plasma turbulence investigations with Langmuir probes} (2011). . arXiv:1111.4091. \url{http://arxiv.org/abs/1111.4091v1}
\item \textbf{Plasma probe characteristics in low density hydrogen pulsed plasmas} (2014). D. I. Astakhov, W. J. Goedheer, C. J. Lee, V. V. Ivanov, V. M. Krivtsun, A. I. Zotovich, S. M. Zyryanov, D. V. Lopaev, F. Bijkerk. arXiv:1412.3036. \url{http://arxiv.org/abs/1412.3036v1}
\item \textbf{Sheath thickness measurements with the biased plasma impedance probe, agreement with Child-Langmuir scaling} (2026). John Whitlock Brooks, Richeek Dutta. arXiv:2602.08743. \url{http://arxiv.org/abs/2602.08743v1}
\item \textbf{Langmuir probe diagnostic studies of pulsed hydrogen plasmas in planar microwave reactors} (2002). A. Rousseau, E. Teboul, N. Lang, M. Hannemann, J. Ropcke. arXiv:physics/0206032. \url{http://arxiv.org/abs/physics/0206032v2}
\item \textbf{Thomson scattering from high-temperature high-density plasmas revisited} (2009). Jian Zheng, C. X. Yu. arXiv:0901.0346. \url{http://arxiv.org/abs/0901.0346v1}
\item \textbf{Correlation length scalings in fusion edge plasma turbulence computations} (2011). Stefan J. Konzett, Dirk Reiser, Alexander Kendl. arXiv:1109.1997. \url{http://arxiv.org/abs/1109.1997v1}
\item \textbf{Langmuir-probe measurement of electron temperature in a flowing high-density plasma} (1970). R. M. Clements, P. R. Smy. Electronics Letters. \url{https://doi.org/10.1049/el:19700373}
\item \textbf{Langmuir probe diagnostics of microwave electron cyclotron resonance (ECR) plasma} (2008). S. B. Singh, N. Chand, D. S. Patil. Vacuum. \url{https://doi.org/10.1016/j.vacuum.2008.05.030}
\item \textbf{Comparison of Langmuir probe and spectroscopic electron temperature measurements} (1971). M. Numano, O. Furukawa, I. Michiyoshi. Plasma Physics. \url{https://doi.org/10.1088/0032-1028/13/10/014}
\item \textbf{Measurement of Electron Density and Temperature and Their Fluctuations Using Modified Triple Langmuir Probe Grounded through Finite Resistance} (2007). M. Takeuchi, K. Toi, R. Ikeda, C. Suzuki, CHS Experimental Group. Plasma and Fusion Research. \url{https://doi.org/10.1585/pfr.2.s1091}
\item \textbf{Comparison of Langmuir probe and laser Thomson scattering for plasma density and electron temperature measurements in HiPIMS plasma} (2019). Peter J. Ryan, James W. Bradley, Mark D. Bowden. Physics of Plasmas. \url{https://doi.org/10.1063/1.5094602}
\item \textbf{Ionospheric Langmuir Probe Electron Temperature Asymmetry and Magnetic Field Connectivity} (2017). Richard Marchand. IEEE Transactions on Plasma Science. \url{https://doi.org/10.1109/tps.2016.2619668}
\item \textbf{Langmuir Probe Diagnostics of a Low Temperature Non-Isothermal Plasma in a Weak Magnetic Field} (1997). P. Kudrna, E. Passoth. Contributions to Plasma Physics. \url{https://doi.org/10.1002/ctpp.2150370504}
\item \textbf{Langmuir probe measurements of electron temperature and electron density behind the skimmer of an inductively coupled plasma mass spectrometer} (1991). H. S. Niu, Ke Hu, R. S. Houk. Spectrochimica Acta Part B: Atomic Spectroscopy. \url{https://doi.org/10.1016/0584-8547(91)80082-e}
\item \textbf{Langmuir Probe for the Measurement of Electron Density and Electron Temperature in the Ionosphere} (1967). Satya Prakash, B. H. Subbaraya. Review of Scientific Instruments. \url{https://doi.org/10.1063/1.1721035}
\item \textbf{Electron energy distribution function, plasma potential and electron density measured by Langmuir probe in tokamak edge plasma} (2009). Tsv. K. Popov, P. Ivanova, J. Stockel, R. Dejarnac. Plasma Physics and Controlled Fusion. \url{https://doi.org/10.1088/0741-3335/51/6/065014}
\item \textbf{Langmuir probe measurements of the ion extraction process in inductively coupled plasma mass spectrometry-I. Spatially resolved determination of electron density and electron temperature} (1994). Hongsen Niu, R. S. Houk. Spectrochimica Acta Part B: Atomic Spectroscopy. \url{https://doi.org/10.1016/0584-8547(94)80109-6}
\item \textbf{Comparison of Langmuir probe and Thomson scattering measurements of electron temperature and density} (1974). Larry A. Jones. Journal of Applied Physics. \url{https://doi.org/10.1063/1.1663216}
\item \textbf{Measurements of electron density in different low pressure plasmas using the multipole resonance probe, Langmuir probe and optical emission spectroscopy} (2016). Moritz Oberberg, Stefan Ries, Nikita Bibinov, Peter Awakowicz. 2016 IEEE International Conference on Plasma Science (ICOPS). \url{https://doi.org/10.1109/plasma.2016.7534183}
\item \textbf{New Method for the Determination of the Electron Temperature from Langmuir Probe Characteristics} (2010). M. Hannemann. Contributions to Plasma Physics. \url{https://doi.org/10.1002/ctpp.201010135}
\item \textbf{Determination of the electron density in current-less argon plasma using Langmuir probe measurements} (2004). Tsv. K. Popov, M. Dimitrova, F. M. Dias. Vacuum. \url{https://doi.org/10.1016/j.vacuum.2004.07.079}
\item \textbf{Spatial distribution of electron temperature and ion density in laser induced ruby (Al2O3:Cr3+) plasma using Langmuir probe} (2012). S. Kumari, A. Kushwaha, A. Khare. Journal of Instrumentation. \url{https://doi.org/10.1088/1748-0221/7/05/c05017}
\item \textbf{Measurement of Spatial Distributions of Electron Density and Electron Temperature in Direct Current Glow Discharge by Double Langmuir Probes} (2002). Wei-Han Tao, Hirotsugu K. Yasuda. Plasma Chemistry and Plasma Processing. \url{https://doi.org/10.1023/a:1014803711271}
\item \textbf{Comparison of Langmuir Probe and rf Capacitance Probe Electron Density Measurements in Low-Density Plasmas} (1972). Jack Smith. Journal of Applied Physics. \url{https://doi.org/10.1063/1.1661653}
\item \textbf{Time-resolved Langmuir probe measurements of plasma density and temperature in electromagnetic shock-tubes} (1965). W. Fucks, R. Theenhaus. Journal of Nuclear Energy Part C: Plasma Physics. \url{https://doi.org/10.1088/0368-3281/7/2/408}
\item \textbf{Diagnostics of Tin Droplet-Based Laser Produced Plasma by Triple Langmuir Probe} (2024). Zegong Chen, Duluo Zuo, Xinbing Wang. SSRN preprint. \url{https://doi.org/10.2139/ssrn.4925658}
\item \textbf{Sweep Langmuir Probe and Triple Probe Diagnostics for Transient Plasma Produced by Hypervelocity Impact} (2012). Enling Tang, Shenghai Xiang, Minghai Yang, Lexin Li. Plasma Science and Technology. \url{https://doi.org/10.1088/1009-0630/14/8/12}
\item \textbf{Comparative analyses of plasma probe diagnostics techniques} (2015). . Journal of Applied Physics. \url{https://pubs.aip.org/aip/jap/article/118/23/233302/141689/Comparative-analyses-of-plasma-probe-diagnostics}
\item \textbf{Recommended Practice for Use of Langmuir Probes in Electric Propulsion Testing} (2017). Daniel L. Brown, Michael L. R. Walker, James Szabo, Wensheng Huang, John E. Foster. Journal of Propulsion and Power. \url{https://arc.aiaa.org/doi/10.2514/1.B35531}
\end{enumerate}

\section*{10\quad Optical emission spectroscopy \textnormal{(\normalsize 33 references)}}
OES exploits the plasma's own emission to infer electron temperature and excited-species densities without perturbation. The literature spans line-ratio methods, actinometry and imaging spectroscopy. OES supplies one of the four complementary constraints in the fusion.
\begin{enumerate}[leftmargin=1.4em]
\item \textbf{First hydrogen operation of NIO1: characterization of the source plasma by means of an optical emission spectroscopy diagnostic} (2022). M. Barbisan, C. Baltador, B. Zaniol, M. Cavenago, U. Fantz, R. Pasqualotto, G. Serianni, L. Vialetto, D. Wünderlich. arXiv:2211.03360. \url{http://arxiv.org/abs/2211.03360v1}
\item \textbf{Hydrogen electron cyclotron resonance ion sources plasma characterization based on simple optical emission spectroscopy} (2025). Jorge Feuchtwanger, Víctor Etxebarria, Joaquín Portilla, Josu Jugo, Inari Badillo, Iñigo Arredondo. arXiv:2502.05080. \url{http://arxiv.org/abs/2502.05080v1}
\item \textbf{Automated Estimation of Plasma Temperature and Density from Emission Spectroscopy} (2023). Todd A. Oliver, Craig Michoski, Samuel Langendorf, Andrew LaJoie. arXiv:2312.12674. \url{http://arxiv.org/abs/2312.12674v1}
\item \textbf{Experimental Characterization and Simulation-Based Analysis of the RIT-2.5 Radiofrequency Ion Thruster} (2025). N. Joshi, O. Meusel. arXiv:2509.10738. \url{http://arxiv.org/abs/2509.10738}
\item \textbf{Optical emission spectroscopy for simultaneous measurement of plasma electron density and temperature in a low-pressure microwave induced plasma} (2009). N. Konjević, S. Jovićević, M. Ivković. Physics of Plasmas. \url{https://doi.org/10.1063/1.3240325}
\item \textbf{An optical emission spectroscopy method for determining the electron temperature and density in low-temperature xenon plasma by using a collisional-radiative model considering the hyperfine structure of emission line into metastable state} (2023). Yan-Fei Wang, Xi-Ming Zhu. Spectrochimica Acta Part B: Atomic Spectroscopy. \url{https://doi.org/10.1016/j.sab.2023.106777}
\item \textbf{Optical emission spectroscopy of an industrial thermal atmospheric pressure plasma jet: Parametric study of electron temperature} (2023). Nikolaos Giannakaris, Gustav Gürtler, Thomas Stehrer. Spectrochimica Acta Part B: Atomic Spectroscopy. \url{https://doi.org/10.1016/j.sab.2023.106736}
\item \textbf{Electron temperature characterization of H2 processing plasma by optical emission spectroscopy} (2024). Shota Nunomura, Naomichi Ezumi. Applied Physics Express. \url{https://doi.org/10.35848/1882-0786/ad87aa}
\item \textbf{Determination of the electron temperature by optical emission spectroscopy in a 13.56 MHz dusty methane plasma: Influence of the power} (2009). Véronique Massereau-Guilbaud, Isabelle Géraud-Grenier, André Plain. Journal of Applied Physics. \url{https://doi.org/10.1063/1.3267292}
\item \textbf{Study on electron temperature in an ablative pulsed plasma thruster by optical emission spectroscopy} (2013). Weiwei Wang, Feng Liu, Xue Wang. EPL (Europhysics Letters). \url{https://doi.org/10.1209/0295-5075/101/55001}
\item \textbf{Optical Emission Spectroscopic (OES) analysis for diagnostics of electron density and temperature in non-equilibrium argon plasma based on collisional-radiative model} (2019). . Advances in Physics: X (Taylor \textbackslash{}\& Francis). \url{https://doi.org/10.1080/23746149.2019.1592707}
\item \textbf{Synthetic diagnostic for the beam emission spectroscopy diagnostic using a full optical integration} (2017). L. Hausammann, R. M. Churchill, L. Shi. Plasma Physics and Controlled Fusion. \url{https://doi.org/10.1088/1361-6587/aa5234}
\item \textbf{An Investigation Into Electron Temperature, Number Density, and Optical Emission Spectroscopy of the DC Hollow-Cathode Plasma Discharge in N2O} (2012). Farook Bashir Yousif, Aidee Berenice Mondragon. IEEE Transactions on Plasma Science. \url{https://doi.org/10.1109/tps.2012.2192453}
\item \textbf{The Relationship Between Characteristics of DLC Film and Electron Temperature Measured by Optical Emission Spectroscopy} (2016). K. Takizawa, S. Kunitsugu. IEEE Transactions on Plasma Science. \url{https://doi.org/10.1109/tps.2016.2619681}
\item \textbf{Plasma Diagnostic by Optical Emission Spectroscopy on Dolomite and Cross-Validation Using SEM-EDX and Electron Probe Micro-Analysis} (2020). M. Fahad, M. Abrar, K. H. Shah. Plasma Physics Reports. \url{https://doi.org/10.1134/s1063780x2003006x}
\item \textbf{Optical emission spectroscopy of helium helps deduce plasma electron density and temperature} (2020). Adam Liebendorfer. Scilight (AIP). \url{https://doi.org/10.1063/10.0000853}
\item \textbf{Investigation of the temperature and electron density of the gliding arc discharge plasma by optical emission spectroscopy} (2024). Mohammed A. Ismail, Saba J. Kadhem. Journal of Optics (Springer). \url{https://doi.org/10.1007/s12596-024-02048-w}
\item \textbf{Optical emission spectroscopy as a real time diagnostic tool for plasma-assisted deposition of TiN} (1997). Mohamed Boumerzoug, Marcel Boudreau, Peter Mascher. Plasma Chemistry and Plasma Processing. \url{https://doi.org/10.1007/bf02766814}
\item \textbf{Study of argon ions density and electron temperature and density in magnetron plasma by optical emission spectroscopy and collisional-radiative model} (2017). Kirill E. Evdokimov, Maxim Konischev, Vladimir Pichugin. Resource-Efficient Technologies. \url{https://doi.org/10.18799/24056529/2017/2/123}
\item \textbf{Effect of an inductively coupled plasma mass spectrometry sampler interface on electron temperature, electron number density, gas-kinetic temperature and analyte emission intensity upstream in the plasma} (2002). Scott A. Lehn, Kelly A. Warner, Mao Huang. Spectrochimica Acta Part B: Atomic Spectroscopy. \url{https://doi.org/10.1016/s0584-8547(02)00147-7}
\item \textbf{Plasma diagnostic on a low-flow plasma for inductively coupled plasma optical emission spectrometry} (2008). Carsten Engelhard, George C.-Y. Chan, Gerardo Gamez. Spectrochimica Acta Part B: Atomic Spectroscopy. \url{https://doi.org/10.1016/j.sab.2008.03.015}
\item \textbf{Optical emission spectroscopy of electron-cyclotron-resonance-heated helium mirror plasmas} (1991). K. L. Junck, M. L. Brake, W. D. Getty. Plasma Chemistry and Plasma Processing. \url{https://doi.org/10.1007/bf01447032}
\item \textbf{Determination of electron temperature in ion plating discharges by optical emission spectroscopy} (1991). K. S. Fancey, A. Matthews. Vacuum. \url{https://doi.org/10.1016/0042-207x(91)90010-g}
\item \textbf{Temporally resolved measurements of gas temperature using optical emission spectroscopy enhanced with probing nanosecond plasma pulse} (2020). Xingxing Wang, Alexey Shashurin. AIAA Scitech 2020 Forum. \url{https://doi.org/10.2514/6.2020-0439}
\item \textbf{Concept of a Diagnostic System for Measuring the Electron Temperature Profile of Plasma from the Intensity of Electron Cyclotron Emission for the TRT Facility} (2024). N. A. Solovev, D. E. Dias Mikhaylova. Plasma Physics Reports. \url{https://doi.org/10.1134/s1063780x24600257}
\item \textbf{Measurements of electron number density and temperature in a supersonic plasma jet by optical emission spectroscopy} (2021). A. S. Pashchina. Journal of Physics: Conference Series. \url{https://doi.org/10.1088/1742-6596/2100/1/012003}
\item \textbf{Measuring electron density, electric field and temperature of a micro-discharge air plasma jet using optical emission spectroscopy} (2015). Hamed Abbasi, Majid Nazeri, Shahriar Mirpour. 2015 2nd International Conference on Knowledge-Based Engineering and Innovation (KBEI). \url{https://doi.org/10.1109/kbei.2015.7436207}
\item \textbf{The Study on the Arc Plasma Temperature Measurement by Optical Emission Spectroscopy with Fiber Optical Transmission} (1998). Junde Wang, Baoming Li, Binghe Gu. Spectroscopy Letters. \url{https://doi.org/10.1080/00387019808006775}
\item \textbf{Electron Temperature Measurement of Argon Capacitive Coupling Plasma by Optical-Emission Spectroscopy} (2024). Yasir Aljawadi. International Academic Journal of Science and Engineering. \url{https://doi.org/10.9756/iajse/v11i1/iajse1121}
\item \textbf{Fast Optical Emission Spectroscopy Measurements of Time-Resolved Electron Temperature in a Hall Thruster} (2024). . AIAA SciTech 2024 Forum. \url{https://doi.org/10.2514/6.2024-2165}
\item \textbf{Tomographic characterization of electron properties in xenon Hall thruster plasmas using collisional–radiative optical emission spectroscopy} (2024). . Plasma Sources Science and Technology. \url{https://doi.org/10.1088/1361-6595/ae3ad4}
\item \textbf{Combination of optical emission spectroscopy and multivariate data analysis techniques as a versatile non-invasive tool for characterizing xenon/krypton mixed gas plasma inside operating ion thrusters} (2022). . Journal of Applied Physics. \url{https://pubs.aip.org/aip/jap/article/131/5/053301/2836281}
\item \textbf{Time-resolved electron temperature oscillations in Hall thrusters} (2025). . Journal of Applied Physics. \url{https://pubs.aip.org/aip/jap/article/137/24/243303/3350951}
\end{enumerate}

\section*{11\quad Collisional-radiative models and line-ratio methods \textnormal{(\normalsize 32 references)}}
Collisional-radiative modelling links measured emission to plasma parameters by solving population balances of excited states. It is the interpretive backbone that makes OES quantitative. It is a required physics model in the inversion chain.
\begin{enumerate}[leftmargin=1.4em]
\item \textbf{Spectroscopic measurement of electron temperature and density in argon plasmas based on collisional-radiative model} (2000). K. Kano, M. Suzuki, H. Akatsuka. Plasma Sources Science and Technology. \url{https://doi.org/10.1088/0963-0252/9/3/309}
\item \textbf{Optical emission spectroscopy in low-temperature plasmas containing argon and nitrogen: determination of the electron temperature and density by the line-ratio method} (2010). X-M Zhu, Y-K Pu. Journal of Physics D: Applied Physics. \url{https://iopscience.iop.org/article/10.1088/0022-3727/43/40/403001}
\item \textbf{Evaluation of a collisional radiative model for electron temperature determination in hydrogen plasma} (2022). Vinoth S., others. Review of Scientific Instruments. \url{https://pubs.aip.org/aip/rsi/article-abstract/93/9/093503/2848928/Evaluation-of-a-collisional-radiative-model-for}
\item \textbf{Electron temperature measurement by a helium line intensity ratio method in helicon plasmas} (2001). Y. Kawachi et al.. Physics of Plasmas. \url{https://pubs.aip.org/aip/pop/article-abstract/8/12/5303/264310/Electron-temperature-measurement-by-a-helium-line}
\item \textbf{Validation of collisional radiative modelling of emission line ratios for helium beam plasma diagnostic} (2009). . Journal of Quantitative Spectroscopy and Radiative Transfer. \url{https://www.sciencedirect.com/science/article/abs/pii/S0022407309000028}
\item \textbf{Optical emission spectroscopy and collisional-radiative modeling for low temperature Ar plasmas} (2018). . Journal of Quantitative Spectroscopy and Radiative Transfer. \url{https://www.sciencedirect.com/science/article/abs/pii/S0022407318307453}
\item \textbf{Comparison of He I line intensity ratio method and electrostatic probe for electron density and temperature measurements in NAGDIS-II} (2006). . Journal Article (OSTI). \url{https://www.osti.gov/biblio/20782471}
\item \textbf{Diagnosis of electron density and temperature by using collisional radiative model in capacitively coupled Ar plasmas I: triple-frequency discharges} (2020). Hao Zheng, Jidun Wu, Qilu Cao, Jiaojiao Zhang, Xiaojiang Huang. arXiv:2010.10714. \url{http://arxiv.org/abs/2010.10714}
\item \textbf{Ionisation state, excited populations and emission of impurities in dynamic finite density plasmas. I: The generalised collisional-radiative model for light elements} (2005). H. P. Summers, W. J. Dickson, M. G. O'Mullane, N. R. Badnell, A. D. Whiteford, D. H. Brooks, J. Lang, S. D. Loch, D. C. Griffin. arXiv:astro-ph/0511561. \url{http://arxiv.org/abs/astro-ph/0511561}
\item \textbf{Non-Maxwellian rate coefficients for electron and ion collisions in Rydberg plasmas: implications for excitation and ionization} (2020). Daniel Vrinceanu, Roberto Onofrio, Hossein R. Sadeghpour. arXiv:2012.04069. \url{http://arxiv.org/abs/2012.04069}
\item \textbf{Determination of the Electron Temperature of Atmospheric Pressure Argon Plasmas by Absolute Line Intensities and a Collisional Radiative Model} (2014). Leila Taghizadeh, Anton Nikiforov, Rino Morent, Joost van der Mullen, Christophe Leys. Plasma Processes and Polymers. \url{https://doi.org/10.1002/ppap.201400022}
\item \textbf{Collisional-radiative model of hydrogen low-temperature plasma: Processes and cross sections of electron-molecule collisions} (2011). V. A. Shakhatov, Yu. A. Lebedev. High Temperature. \url{https://doi.org/10.1134/s0018151x11020131}
\item \textbf{Spectroscopic Measurement of Electron Temperature and Density in an Argon Plasma Jet Based on Collisional-Radiative Model} (2001). K. Kano, M. Suzuki, H. Akatsuka. Contributions to Plasma Physics. \url{https://doi.org/10.1002/1521-3986(200101)41:1<91::aid-ctpp91>3.0.co;2-9}
\item \textbf{Time-resolved argon theta-pinch plasma properties by line ratio method with collisional-radiative model} (2014). Warner C. Meeks, Joshua L. Rovey. 2014 IEEE 41st International Conference on Plasma Sciences (ICOPS). \url{https://doi.org/10.1109/plasma.2014.7012325}
\item \textbf{Detailed Collisional Radiative Model for Low-Temperature Neon Plasma} (2020). S. Gupta, S. S. Baghel, R. Srivastava, R. K. Gangwar. 2020 IEEE International Conference on Plasma Science (ICOPS). \url{https://doi.org/10.1109/icops37625.2020.9717743}
\item \textbf{Bayesian Based Inference of Electron Temperature and Density by Using a Collisional-Radiative Model in Low Density Helium Plasmas} (2018). B. Ahn, Y. G. Lim, Y.-c. Ghim. 2018 IEEE International Conference on Plasma Science (ICOPS). \url{https://doi.org/10.1109/icops35962.2018.9575863}
\item \textbf{Spectroscopic Measurements of Electron Temperature and Electron Density in Electron Beam Plasma Generator Based on Collisional Radiative Model} (2007). H. Goktas, A. Demir, E. Kacar, H. Hegazy, R. Turan, G. Oke, A. Seyhan. Spectroscopy Letters. \url{https://doi.org/10.1080/00387010600934469}
\item \textbf{OES diagnostics of atmospheric pressure argon plasma: Electron temperature and density assessment through visible bremsstrahlung inversion method and collisional-radiative model} (2024). Keren Lin, Thijs van der Gaag, Wataru Kikuchi, Hiroshi Akatsuka, Motoshi Goto. Applied Physics Letters. \url{https://doi.org/10.1063/5.0198632}
\item \textbf{Development of a collisional radiative model for interpreting the spectroscopic measurements of Ar II line emission} (2003). I. V. Moskalenko, D. A. Shuvaev. Plasma Physics Reports. \url{https://doi.org/10.1134/1.1625994}
\item \textbf{Analysis of Detached Recombining Plasmas by Collisional-Radiative Model with Energetic Electron Component} (2002). N. Ohno, M. Motoyama, S. Takamura. Contributions to Plasma Physics. \url{https://doi.org/10.1002/1521-3986(200204)42:2/4<419::aid-ctpp419>3.0.co;2-l}
\item \textbf{X-ray spectra of high temperature tungsten plasma calculated with collisional radiative model} (2013). Jun Wang, Hong Zhang, Xin-Lu Cheng. Chinese Physics B. \url{https://doi.org/10.1088/1674-1056/22/8/085201}
\item \textbf{Development of diagnostics of electron density and temperature for atmospheric-pressure helium plasma based on optical emission spectroscopy analysis and a collisional-radiative model} (2023). Keren Lin, Atsushi Nezu, Hiroshi Akatsuka. Japanese Journal of Applied Physics. \url{https://doi.org/10.35848/1347-4065/accc93}
\item \textbf{Determination of Non-Maxwellian Electron Energy Distributions in Low-Pressure Plasmas by Using the Optical Emission Spectroscopy and a Collisional-Radiative Model} (2011). Ximing Zhu, Yikang Pu. Plasma Science and Technology. \url{https://doi.org/10.1088/1009-0630/13/3/02}
\item \textbf{Petrov-Galerkin Model Reduction for Collisional-Radiative Argon Plasma} (2025). Ivan Zanardi, Alessandro Meini, Alberto Padovan, Daniel Bodony, Marco Panesi. SSRN preprint. \url{https://doi.org/10.2139/ssrn.5290830}
\item \textbf{Diagnosing the lightning plasma temperature based on the collisional-radiative model} (2020). W. Zhang, J. Jiang, P. Yuan. Journal of Physics: Conference Series. \url{https://doi.org/10.1088/1742-6596/1412/19/192004}
\item \textbf{Collisional-Radiative Model of Helium Plasma. I. Description of the Model} (1991). J. Slavik. Contributions to Plasma Physics. \url{https://doi.org/10.1002/ctpp.2150310605}
\item \textbf{Determination of the Electron Density and Electron Temperature in a Magnetron Discharge Plasma Using Optical Spectroscopy and the Collisional-Radiative Model of Argon} (2017). K. E. Evdokimov, M. E. Konishchev, V. F. Pichugin, A. A. Pustovalova, N. M. Ivanova, Ch. Sun. Russian Physics Journal. \url{https://doi.org/10.1007/s11182-017-1137-0}
\item \textbf{Generalized Collisional Radiative Model Using Screened Hydrogenic Levels} (2016). H.-K. Chung, S. B. Hansen, H. A. Scott. Modern Methods in Collisional-Radiative Modeling of Plasmas (Springer). \url{https://doi.org/10.1007/978-3-319-27514-7_3}
\item \textbf{Collisional-radiative model for neutral helium in plasma revisited} (2003). M. Goto. Journal of Quantitative Spectroscopy and Radiative Transfer. \url{https://doi.org/10.1016/s0022-4073(02)00060-2}
\item \textbf{Theoretical investigation of iodine plasma through detailed electron impact excitation cross-section calculations and collisional-radiative model analysis} (2024). Ayushi Agrawal, Shivam Gupta, Lalita Sharma, Rajesh Srivastava. Plasma Sources Science and Technology. \url{https://doi.org/10.1088/1361-6595/ad8df9}
\item \textbf{Electron energy distribution function in the collisional-radiative model of an argon plasma} (1985). J. Vlcek, V. Pelikan. Journal of Physics D: Applied Physics. \url{https://doi.org/10.1088/0022-3727/18/3/005}
\item \textbf{Collisional-radiative model of hydrogen: Balmer-line-ratio electron temperature diagnostics (PPPL/PPST report)} (2022). Vinoth S. et al.. Princeton Plasma Physics Laboratory (PPST). \url{https://w3.pppl.gov/ppst/docs/vinoth.pdf}
\end{enumerate}

\section*{12\quad Laser-induced fluorescence diagnostics \textnormal{(\normalsize 36 references)}}
LIF uses tuned laser excitation and fluorescence detection to probe specific species non-intrusively. The literature covers schemes, calibration and application across low-temperature and high-density plasmas. It is the enabling technique for ion-resolved measurement.
\begin{enumerate}[leftmargin=1.4em]
\item \textbf{Laser induced fluorescence diagnostic for velocity distribution functions: applications, physics, methods and developments} (2021). Chi-Shung Yip, Di Jiang. Plasma Science and Technology. \url{https://doi.org/10.1088/2058-6272/abec62}
\item \textbf{Development of the megahertz planar laser-induced fluorescence diagnostic for plasma turbulence visualization} (2004). Aleksey Kuritsyn, Fred M. Levinton. Review of Scientific Instruments. \url{https://doi.org/10.1063/1.1787148}
\item \textbf{A new pure ion plasma device with laser induced fluorescence diagnostic} (1997). F. Anderegg, X.-P. Huang, E. Sarid, C. F. Driscoll. Review of Scientific Instruments. \url{https://doi.org/10.1063/1.1148119}
\item \textbf{Development of a new time-resolved laser-induced fluorescence diagnostic} (2012). C. J. Durot, A. D. Gallimore. IEEE International Conference on Plasma Science. \url{https://doi.org/10.1109/plasma.2012.6384014}
\item \textbf{Development of laser induced fluorescence diagnostic for measuring the parameters of plasma containing rare gas species} (2010). A. V. Gorbunov, N. A. Molodtsov, I. V. Moskalenko, D. A. Shcheglov. Review of Scientific Instruments. \url{https://doi.org/10.1063/1.3475799}
\item \textbf{Novel techniques for plasma diagnostics: Electron paramagnetic resonance and laser-induced fluorescence spectroscopy} (1981). Terry A. Miller. Plasma Chemistry and Plasma Processing. \url{https://doi.org/10.1007/bf00566372}
\item \textbf{Plasma electron density measurements by the laser- and collision-induced fluorescence method} (1983). K. Tsuchida, S. Miyake, K. Kadota, J. Fujita. Plasma Physics. \url{https://doi.org/10.1088/0032-1028/25/9/003}
\item \textbf{Plasma Boundary Diagnostics by Laser Induced Fluorescence} (1982). E. Hintz. Physica Scripta. \url{https://doi.org/10.1088/0031-8949/1982/t2b/023}
\item \textbf{Laser Induced Fluorescence} (2003). Francis F. Chen, Jane P. Chang. Lecture Notes on Principles of Plasma Processing. \url{https://doi.org/10.1007/978-1-4615-0181-7_29}
\item \textbf{Coherent detection of the complete linear electrostatic plasma response of plasma ions using laser-induced fluorescence} (2002). F. Skiff. IEEE Transactions on Plasma Science. \url{https://doi.org/10.1109/tps.2002.1003906}
\item \textbf{Laser-Induced Fluorescence of Atoms in a Power-Modulated Inductively Coupled Plasma for Trace Detection and Diagnostic Study} (1992). Gregory C. Turk. Applied Spectroscopy. \url{https://doi.org/10.1366/0003702924123908}
\item \textbf{Effects of synchronization of laser induced fluorescence and corona induced fluorescence in a pulsed corona reactor} (2005). Xudong Hu, S. F. Legowski, S. V. B. J. Garikipati, Gui-Bing Zhao, M. Radosz. IEEE Transactions on Plasma Science. \url{https://doi.org/10.1109/tps.2005.851968}
\item \textbf{The Application of Planar Laser-Induced Fluorescence as a Diagnostic in Plasma Processing} (1987). J. K. Partin, J. R. Fincke, C. L. Jeffery, S. C. Snyder, J. A. Batdorf. MRS Proceedings. \url{https://doi.org/10.1557/proc-98-149}
\item \textbf{Laser induced fluorescence study of plasma-insulator wall interaction involving secondary electron emission}. Valentin Pigeon. Thesis / dissertation. \url{https://doi.org/10.70675/1fdc38c0z6e32z416aza1d5z996dd7741327}
\item \textbf{Two-Photon Absorption Laser Induced Fluorescence Measurements of Neutral Density in Helicon Plasma}. Matthew E. Galante. Thesis / dissertation (West Virginia University). \url{https://doi.org/10.33915/etd.3653}
\item \textbf{Laser Induced Fluorescence Studies of Electrostatic Double Layers in an Expanding Helicon Plasma}. Jerry Carr Jr.. Thesis / dissertation (West Virginia University). \url{https://doi.org/10.33915/etd.3648}
\item \textbf{Laser induced fluorescence in a high power argon helicon plasma} (2010). John Scharer, Matt Wiebold. IEEE International Conference on Plasma Science. \url{https://doi.org/10.1109/plasma.2010.5534328}
\item \textbf{Applications of Laser-Induced Fluorescence Spectroscopy for Combustion and Plasma Diagnostics} (2017). Robert P. Lucht. Laser Spectroscopy and Its Applications (book chapter). \url{https://doi.org/10.1201/9780203749104-9}
\item \textbf{Analytical and Diagnostic Applications of Laser Induced Fluorescence in Flames and Plasmas} (1985). N. Omenetto. NATO ASI Series: Analytical Laser Spectroscopy. \url{https://doi.org/10.1007/978-1-4613-2441-6_7}
\item \textbf{Time-resolved Thermometry in a Condensing Laser-ablated Copper Plasma Plume by Doppler-resolved Laser Induced Fluorescence} (1994). David K. Zerkle, Andrew D. Sappey. Laser Applications to Chemical Analysis (OSA). \url{https://doi.org/10.1364/laca.1994.thb.2}
\item \textbf{Laboratory Simulation of Magnetospheric Plasma Phenomena Using Laser Induced Fluorescence as a Diagnostic} (2013). Nathan Rynn. Geophysical Monograph Series: Ion Acceleration in the Magnetosphere and Ionosphere. \url{https://doi.org/10.1029/gm038p0235}
\item \textbf{A confocal laser-induced fluorescence diagnostic with a ring-shaped laser beam} (2023). I. Romadanov, Y. Raitses. arXiv:2303.12580. \url{https://arxiv.org/abs/2303.12580}
\item \textbf{FPGA based demodulation of laser induced fluorescence in plasmas} (2018). Sean Mattingly, Fred Skiff. arXiv:1808.04692. \url{https://arxiv.org/abs/1808.04692}
\item \textbf{The formation of atomic oxygen and hydrogen in atmospheric pressure plasmas containing humidity: picosecond two-photon absorption laser induced fluorescence and numerical simulations} (2020). Sandra Schroter, Jerome Bredin, Andrew R. Gibson, Andrew West, James P. Dedrick, Erik Wagenaars, Kari Niemi, Timo Gans, Deborah O'Connell. arXiv:2004.01558. \url{https://arxiv.org/abs/2004.01558}
\item \textbf{Quantitative evaluation of laser-induced fluorescence in magnetized plasma accounting for disalignment effect} (2020). Roman Bergert, Slobodan Mitic, Markus H. Thoma. arXiv:2002.09218. \url{https://arxiv.org/abs/2002.09218}
\item \textbf{Laboratory plasma physics experiments using merging supersonic plasma jets} (2014). S. C. Hsu, A. L. Moser, E. C. Merritt, C. S. Adams, J. P. Dunn, S. Brockington, A. Case, M. Gilmore, A. G. Lynn, S. J. Messer, F. D. Witherspoon. arXiv:1408.0323. \url{https://arxiv.org/abs/1408.0323}
\item \textbf{Time-evolution of the ion velocity distribution function in the discharge of a Hall effect thruster} (2009). . arXiv:0902.1432. \url{https://arxiv.org/abs/0902.1432}
\item \textbf{Ultrafast ps-TALIF and streak camera diagnostics of atomic hydrogen in a helium microplasma jet} (2025). . arXiv:2507.11586. \url{https://arxiv.org/abs/2507.11586}
\item \textbf{Time-resolved laser-induced fluorescence measurement of ion and neutral dynamics in a Hall thruster during ionization oscillations} (2015). . Journal of Applied Physics. \url{https://pubs.aip.org/aip/jap/article/118/23/233301/141682}
\item \textbf{On limitations of laser-induced fluorescence diagnostics for xenon ion velocity distribution function measurements in Hall thrusters} (2018). . Physics of Plasmas. \url{https://pubs.aip.org/aip/pop/article/25/3/033501/1060373}
\item \textbf{Laser-induced fluorescence study of a xenon Hall thruster}. . Applied Physics B. \url{https://link.springer.com/article/10.1007/s003400050297}
\item \textbf{Laser-induced fluorescence spectroscopy on neutrals for plasma studies in Hall thrusters} (2025). . Journal of Electric Propulsion. \url{https://link.springer.com/article/10.1007/s44205-025-00125-1}
\item \textbf{Two-photon absorption laser induced fluorescence: rate and density-matrix regimes for plasma diagnostics} (2020). . Plasma Sources Science and Technology. \url{https://iopscience.iop.org/article/10.1088/1361-6595/ab85d0}
\item \textbf{Two-photon absorption laser induced fluorescence (TALIF) detection of atomic iodine in low-temperature plasmas and a revision of the energy levels of I I} (2023). . Journal of Physics B: Atomic, Molecular and Optical Physics. \url{https://iopscience.iop.org/article/10.1088/1361-6455/acb7b6}
\item \textbf{Two-photon absorption laser induced fluorescence measurements of absolute neutral deuterium density, temperature, and bulk flow in Proto-MPEX} (2021). . Physics of Plasmas. \url{https://pubs.aip.org/aip/pop/article/28/8/082501/107022}
\item \textbf{Planar laser-induced fluorescence system for spatiotemporal ion velocity distribution function measurements}. . Journal article (OSTI.GOV record). \url{https://www.osti.gov/biblio/2999373}
\end{enumerate}

\section*{13\quad LIF ion velocity distribution measurements \textnormal{(\normalsize 26 references)}}
Doppler-resolved LIF reconstructs ion velocity distribution functions in and near sheaths, capturing drift and heating. These measurements are the most direct optical access to ion energetics at boundaries. They are the primary observable feeding the energy-flux reconstruction.
\begin{enumerate}[leftmargin=1.4em]
\item \textbf{3D ion velocity distribution function measurement in an electric thruster using laser induced fluorescence tomography} (2017). P. Q. Elias, J. Jarrige, E. Cucchetti. Review of Scientific Instruments. \url{https://doi.org/10.1063/1.5001304}
\item \textbf{Laser-induced-fluorescence observation of ion velocity distribution functions in a plasma sheath} (2006). N. Claire, G. Bachet, U. Stroth, F. Doveil. Physics of Plasmas. \url{https://doi.org/10.1063/1.2206786}
\item \textbf{Laser-Induced Fluorescence Measurement of the Ion-Energy-Distribution Function in a Collisionless Reconnection Experiment} (2005). A. Stark, W. Fox, J. Egedal, O. Grulke, T. Klinger. Physical Review Letters. \url{https://doi.org/10.1103/physrevlett.95.235005}
\item \textbf{A time-resolved laser induced fluorescence study on the ion velocity distribution function in a Hall thruster after a fast current disruption} (2009). S. Mazouffre, D. Gawron, N. Sadeghi. Physics of Plasmas. \url{https://doi.org/10.1063/1.3112704}
\item \textbf{Laser-induced fluorescence saturation effects on ion velocity distribution functions in the vicinity of reflecting surfaces} (2019). V. Pigeon, N. Claire, C. Arnas, S. Mazouffre. Physics of Plasmas. \url{https://doi.org/10.1063/1.5077047}
\item \textbf{Time-resolved laser-induced fluorescence spectroscopy using CW diode laser for diagnostics of argon-ion velocity distribution near AC-biased electrode} (2024). R. Takahashi, S. Kito, K. Eriguchi. Review of Scientific Instruments. \url{https://doi.org/10.1063/5.0200138}
\item \textbf{Azimuthal velocity measurement in the ion beam of a gridded ion thruster using laser-induced fluorescence spectroscopy} (2018). R. Tsukizaki, Y. Yamamoto, D. Koda, H. Kuninaka. Plasma Sources Science and Technology. \url{https://doi.org/10.1088/1361-6595/aa9f9a}
\item \textbf{Laser-Induced-Fluorescence-Derived Hall Effect Thruster Ion Velocity Distribution Visualization} (2011). W. A. Hargus. IEEE Transactions on Plasma Science. \url{https://doi.org/10.1109/tps.2011.2132149}
\item \textbf{Laser Induced Fluorescence Observation of the Ion Velocity Distribution Function Modifications Induced by Electrostatic Ion Shock Waves} (1996). G. Bachet, L. Chérigier, C. Arnas-Capeau, F. Doveil. Journal de Physique III. \url{https://doi.org/10.1051/jp3:1996175}
\item \textbf{Laser-induced fluorescence measurement of plasma ion distribution functions: Errors due to spatially inhomogeneous laser intensities} (1989). M. J. Goeckner, J. Goree, T. E. Sheridan. AIP Conference Proceedings. \url{https://doi.org/10.1063/1.38579}
\item \textbf{Laser-Induced Fluorescence Measurements of Ion Distribution Functions} (1990). J. Goree, M. J. Goeckner. Plasma-Surface Interactions and Processing of Materials. \url{https://doi.org/10.1007/978-94-009-1946-4_10}
\item \textbf{Spatial Profile of Ion Velocity Distribution Function in Helicon High-Density Plasma by Laser Induced Fluorescence Method} (2016). Y. Tanida, D. Kuwahara, S. Shinohara. Transactions of the Japan Society for Aeronautical and Space Sciences. \url{https://doi.org/10.2322/tastj.14.pb_7}
\item \textbf{Parallel Ion Flow Velocity Measurement Using Laser Induced Fluorescence Method in an Electron Cyclotron Resonance Plasma} (2010). S. Yoshimura, A. Okamoto, K. Terasaka, M. Y. Tanaka. Plasma and Fusion Research. \url{https://doi.org/10.1585/pfr.5.s2052}
\item \textbf{Measurement of wave-particle interaction and metastable lifetime using laser-induced fluorescence} (2019). F. Chu, F. Skiff. arXiv:1902.10701. \url{https://arxiv.org/abs/1902.10701}
\item \textbf{Ion dynamic characterization using phase-resolved laser-induced fluorescence spectroscopy in a Hall effect thruster} (2023). Y. Dancheva, P. Coniglio, M. Da Valle, F. Scortecci. arXiv:2302.07310. \url{https://arxiv.org/abs/2302.07310}
\item \textbf{Measurement of XeI and XeII velocity in the near exit plane of a low-power Hall effect thruster by light induced fluorescence spectroscopy} (2013). Y. Dancheva, V. Biancalana, D. Pagano, F. Scortecci. arXiv:1305.5138. \url{https://arxiv.org/abs/1305.5138}
\item \textbf{Diode Laser-Induced Fluorescence of Xenon Ion Velocity Distributions} (2005). T. Smith, B. Ngom, J. Linnell, A. Gallimore. 41st AIAA/ASME/SAE/ASEE Joint Propulsion Conference \textbackslash{}\& Exhibit. \url{https://doi.org/10.2514/6.2005-4406}
\item \textbf{Laser Induced Fluorescence Measurements of Xenon Ion Velocity Distributions near Ceramic Surfaces} (2015). S. Walsh, A. P. Yalin. 53rd AIAA Aerospace Sciences Meeting. \url{https://doi.org/10.2514/6.2015-1828}
\item \textbf{Measurement of the density and velocity distribution of sputtered Al in EBT-S by laser-induced fluorescence} (1982). E. Dullni, E. Hintz, J. B. Roberto. Journal of Nuclear Materials. \url{https://doi.org/10.1016/0022-3115(82)90182-9}
\item \textbf{Understanding Plasmas through Ion Velocity Distribution Function measurements}. S. Chakraborty Thakur. Doctoral dissertation (West Virginia University). \url{https://doi.org/10.33915/etd.3196}
\item \textbf{Ion energy distribution function measurements by laser-induced fluorescence in a dual radio frequency sheath} (2016). N. B. Moore, W. Gekelman, P. Pribyl. Journal of Vacuum Science \textbackslash{}\& Technology A. \url{https://pubs.aip.org/avs/jva/article/34/2/021303/246200}
\item \textbf{2-dimensional ion velocity distributions measured by laser-induced fluorescence above a radio-frequency biased silicon wafer} (2013). N. B. Moore, W. Gekelman, P. Pribyl, Y. Zhang, M. J. Kushner. Physics of Plasmas. \url{https://pubs.aip.org/aip/pop/article-abstract/20/8/083506}
\item \textbf{Report on laser-induced fluorescence transitions relevant for the microelectronics industry and sustainability applications} (2024). . Journal of Vacuum Science \textbackslash{}\& Technology A. \url{https://pubs.aip.org/avs/jva/article-abstract/42/6/063005}
\item \textbf{Single- and two-photon absorption laser-induced fluorescence spectroscopy in rare gases for gridded ion thruster diagnostics} (2022). . EPJ Techniques and Instrumentation. \url{https://doi.org/10.1140/epjti/s40485-022-00077-y}
\item \textbf{Quantitative two-photon laser-induced fluorescence measurements of atomic hydrogen densities, temperatures, and velocities in an expanding thermal plasma} (2002). A. J. Wagner, B. Brumfield, others. Review of Scientific Instruments. \url{https://pubs.aip.org/aip/rsi/article-abstract/73/1/73/436579}
\item \textbf{Measurement and simulation of laser-induced fluorescence from non-equilibrium ultracold neutral plasmas} (2009). A. Denning, S. D. Bergeson, F. Robicheaux. arXiv:0906.4724. \url{https://arxiv.org/abs/0906.4724}
\end{enumerate}

\section*{14\quad Thomson scattering diagnostics \textnormal{(\normalsize 34 references)}}
Thomson scattering provides model-light, spatially resolved electron temperature and density from laser light scattered by electrons. The literature addresses low-density sensitivity, stray-light rejection and incoherent/coherent regimes. It is the high-accuracy anchor added at scale-up.
\begin{enumerate}[leftmargin=1.4em]
\item \textbf{Advances in Thomson scattering diagnostics of plasmas used for chemical analysis} (2020). . Spectrochimica Acta Part B: Atomic Spectroscopy. \url{https://www.sciencedirect.com/science/article/abs/pii/S0584854720304845}
\item \textbf{Laser Thomson scattering measurements of electron temperature and density in a Hall-effect plasma} (2013). Robert L. Washeleski. Michigan Tech Dissertation. \url{https://digitalcommons.mtu.edu/etds/422/}
\item \textbf{Spatially resolved Thomson scattering measurements of electron properties across the acceleration region of a high-power magnetically shielded Hall effect thruster} (2024). . Journal of Applied Physics. \url{https://pubs.aip.org/aip/jap/article/136/11/113303/3313537/Spatially-resolved-Thomson-scattering-measurements}
\item \textbf{Portable Thomson scattering system for temporally resolved plasma measurements under low density conditions} (2024). . Review of Scientific Instruments. \url{https://pubs.aip.org/aip/rsi/article-abstract/95/3/033502/3269894/Portable-Thomson-scattering-system-for-temporally}
\item \textbf{Time resolved electron density and temperature measurements via Thomson scattering in an atmospheric nanosecond pulsed discharge} (2020). Jared Miles, Chase Murray, Ajani Ross, Kristina Lemmer, Jacob Russell, Steven Adams. Plasma Sources Science and Technology. \url{https://doi.org/10.1088/1361-6595/aba114}
\item \textbf{Polarization of incoherent Thomson scattering for electron temperature measurement} (2017). V V Mirnov, D J Den Hartog. Plasma Physics and Controlled Fusion. \url{https://doi.org/10.1088/1361-6587/aa61e2}
\item \textbf{Ponderomotive perturbations of low density low-temperature plasma under laser Thomson scattering diagnostics} (2017). Mikhail N. Shneider. arXiv:1707.07224. \url{http://arxiv.org/abs/1707.07224v2}
\item \textbf{Particle-in-Cell Modeling of Laser Thomson Scattering in Low-Density Plasmas at Elevated Laser Intensities} (2018). Andrew T. Powis, Mikhail N. Shneider. arXiv:1803.06278. \url{http://arxiv.org/abs/1803.06278v1}
\item \textbf{Thomson microwave scattering for electron number density diagnostics of miniature plasmas at low pressure} (2019). Xingxing Wang, Apoorv Ranjan, Animesh Sharma, Mikhail N. Shneider, Alexey Shashurin. arXiv:1905.06430. \url{http://arxiv.org/abs/1905.06430v2}
\item \textbf{Thomson scattering on the Large Plasma Device} (2022). Sofiya Ghazaryan, Marietta Kaloyan, Walter Gekelman, Zoltan Lucky, Stephen Vincena, Shreekrishna P. Tripathi, Patrick Pribyl, Christoph Niemann. arXiv:2204.12616. \url{http://arxiv.org/abs/2204.12616v1}
\item \textbf{Conceptual design of Thomson scattering system with high wavelength resolution in magnetically confined plasmas for electron phase-space measurements} (2025). Kentaro Sakai, Kentaro Tomita, Takeo Hoshi, Akito Nakano, Motoshi Goto, Kenichi Nagaoka, Ryo Yasuhara. arXiv:2511.06330. \url{http://arxiv.org/abs/2511.06330v2}
\item \textbf{On the accuracy of measurements of electron temperature by Thomson scattering diagnostic in the plasma core of the ITER tokamak} (2026). M. Yu. Kantor. arXiv:2606.23386. \url{http://arxiv.org/abs/2606.23386v1}
\item \textbf{Design of a Thomson scattering diagnostic for the SMART tokamak} (2024). M. Kaur, A. Diallo, B. LeBlanc, J. Segado-Fernandez, E. Viezzer, Roger Huxford, A. Mancini, D. J. Cruz-Zabala, M. Podesta, J. W. Berkery, M. Garcia-Munoz. arXiv:2407.13720. \url{http://arxiv.org/abs/2407.13720v1}
\item \textbf{Status of the JET LIDAR Thomson Scattering diagnostic} (2013). M. Maslov, M. N. A. Beurskens, M. Kempenaars, J. Flanagan, JET EFDA contributors. arXiv:1312.3109. \url{http://arxiv.org/abs/1312.3109v1}
\item \textbf{Thomson scattering diagnostic at the gas dynamic trap} (2022). A. Lizunov, T. Berbasova, A. Khilchenko, A. Kvashnin, E. Puryga, A. Sandomirsky, P. Zubarev. arXiv:2208.13478. \url{http://arxiv.org/abs/2208.13478v1}
\item \textbf{High resolution Thomson scattering diagnostic for measurements of radial profiles of electron temperature and density in the gas dynamic trap} (2023). A. Lizunov, T. Berbasova, A. Khilchenko, A. Kvashnin, E. Puryga, A. Sandomirsky, P. Zubarev. Review of Scientific Instruments. \url{https://doi.org/10.1063/5.0123329}
\item \textbf{Collective Thomson scattering diagnostic for the GDT experiment} (2019). A. G. Shalashov, E. D. Gospodchikov, T. A. Khusainov, L. V. Lubyako, O. B. Smolyakova, A. L. Solomakhin. arXiv:1912.03521. \url{http://arxiv.org/abs/1912.03521v3}
\item \textbf{A polarization-based Thomson scattering technique for burning plasmas} (2013). E. Parke, V. V. Mirnov, D. J. Den Hartog. arXiv:1309.4858. \url{http://arxiv.org/abs/1309.4858v1}
\item \textbf{The Thomson scattering cross section in a magnetized, high density plasma} (2019). Archie F. A. Bott, Gianluca Gregori. arXiv:1901.00820. \url{http://arxiv.org/abs/1901.00820v1}
\item \textbf{Thomson scattering in dense plasmas with density and temperature gradients} (2009). C. Fortmann, R. Thiele, R. R. Fäustlin, Th. Bornath, B. Holst, W.-D. Kraeft, V. Schwarz, S. Toleikis, Th. Tschentscher, R. Redmer. arXiv:0903.0466. \url{http://arxiv.org/abs/0903.0466v1}
\item \textbf{Thomson scattering from near-solid density plasmas using soft x-ray free electron lasers} (2006). A. Hoell. arXiv:physics/0611274. \url{http://arxiv.org/abs/physics/0611274v1}
\item \textbf{The effect of bound states on X-ray Thomson scattering for partially ionized plasmas} (2012). J. Nilsen, W. R. Johnson, K. T. Cheng. arXiv:1212.5972. \url{http://arxiv.org/abs/1212.5972v2}
\item \textbf{Thomson scattering measurements of electron temperature and density in an electron cyclotron resonance plasma} (1993). M. D. Bowden, T. Okamoto, F. Kimura, H. Muta, K. Uchino, K. Muraoka, T. Sakoda, M. Maeda, Y. Manabe, M. Kitagawa, T. Kimura. Journal of Applied Physics. \url{https://doi.org/10.1063/1.353046}
\item \textbf{Measurements of Electron Temperature and Density of a Micro-Discharge Plasma Using Laser Thomson Scattering} (2001). Yasuyuki Noguchi, Akira Matsuoka, Mark D. Bowden, Kiichiro Uchino, Katsunori Muraoka. Japanese Journal of Applied Physics. \url{https://doi.org/10.1143/jjap.40.326}
\item \textbf{Temporal evolution of electron density and electron temperature profiles in a non-thermal atmospheric-pressure plasma measured by laser Thomson scattering} (2014). Nima Bolouki, Kentaro Tomita, Safwat Hassaballa, Yukihiko Yamagata, Kiichiro Uchino. Japanese Journal of Applied Physics. \url{https://doi.org/10.7567/jjap.54.016101}
\item \textbf{The spatial-temporal distribution of the electron density and temperature in nanosecond pulsed discharge plasma measured by laser Thomson scattering} (2019). J. Shi, Y. Wang, Y. Zhao, C. Li, C. Feng, H. Ding. Journal of Instrumentation. \url{https://doi.org/10.1088/1748-0221/14/12/c12017}
\item \textbf{Measurement of Electron Density and Temperature Using Laser Thomson Scattering in PANTA} (2017). Kentaro Tomita, Yuta Sato, Nima Bolouki, Tatsuro Shiraishi, Kiichiro Uchino, Yudai Miwa, Tatsuya Kobayashi, Shigeru Inagaki. Plasma and Fusion Research. \url{https://doi.org/10.1585/pfr.12.1401018}
\item \textbf{Dependence on the scattering angle of the electron temperature and electron density in Thomson-scattering measurements on an atmospheric-pressure plasma jet} (2000). S. C. Snyder, D. M. Crawford, J. R. Fincke. Physical Review E. \url{https://doi.org/10.1103/physreve.61.1920}
\item \textbf{An apparatus for the space- and time-resolved measurement of electron temperature and density in a plasma by Thomson scattering of laser light} (1971). S. E. Segre, M. Martone. Il Nuovo Cimento B. \url{https://doi.org/10.1007/bf02899940}
\item \textbf{Thomson scattering diagnostic for a low-density, high-temperature, steady-state plasma} (1982). J. A. Cobble, L. Bighel. Review of Scientific Instruments. \url{https://doi.org/10.1063/1.1136985}
\item \textbf{A New Method for the Determination of Plasma Electron Temperature and Density from Thomson Scattering of an Optical Maser Beam} (1962). T. P. Hughes. Nature. \url{https://doi.org/10.1038/194268b0}
\item \textbf{Isocontour maps of electron temperature, electron number density and gas kinetic temperature in the Ar inductively coupled plasma obtained by laser-light Thomson and Rayleigh scattering} (1992). Mao Huang, D. S. Hanselman, Pengyuan Yang, G. M. Hieftje. Spectrochimica Acta Part B. \url{https://doi.org/10.1016/0584-8547(92)80072-o}
\item \textbf{Electron density and temperature contour plots from a laser-produced plasma using collective ultraviolet Thomson scattering} (1996). S. M. Cameron, M. D. Tracy, J. F. Camacho. IEEE Transactions on Plasma Science. \url{https://doi.org/10.1109/27.491683}
\item \textbf{Electron temperature and density measurement by Thomson scattering with a high repetition rate laser of 20 kHz on LHD} (2022). . Scientific Reports. \url{https://www.nature.com/articles/s41598-022-19328-9}
\end{enumerate}

\section*{15\quad Plasma interferometry and electron density \textnormal{(\normalsize 36 references)}}
Interferometric and refractive methods measure line-integrated electron density from optical phase shift. The literature covers multi-wavelength, heterodyne and dispersion techniques. It supplies a fast, robust density constraint to the fusion.
\begin{enumerate}[leftmargin=1.4em]
\item \textbf{Nanosecond-resolved 266 nm Mach-Zehnder interferometry for electron-density measurements of dense plasmas generated in supercritical fluids} (2026). Kyusang Cho, Juho Lee, Gunsu Yun. arXiv:2606.25327. \url{http://arxiv.org/abs/2606.25327v1}
\item \textbf{Electron density structure measurements with scattered intense laser beam} (2025). K. Sakai, K. Himeno, S. J. Tanaka, T. Asai, others. arXiv:2505.08711. \url{http://arxiv.org/abs/2505.08711v1}
\item \textbf{Deflection of Interferometry Beams due to Transverse Refractive Index Gradient in SST-1} (2025). . arXiv:2501.18193. \url{https://arxiv.org/abs/2501.18193}
\item \textbf{Optical diagnostics of laser-produced plasmas} (2022). . arXiv:2201.08783. \url{https://arxiv.org/abs/2201.08783}
\item \textbf{Laser-plasma electron-density measurement using x-ray interferometry} (2002). Xiquan Fu, Hong Guo. Physical Review E. \url{https://doi.org/10.1103/physreve.65.067401}
\item \textbf{A New Method of Electron Density Measurement by Fabry-Perot Interferometry} (2004). Hitoshi Hojo, Atsushi Mase. Journal of Plasma and Fusion Research. \url{https://doi.org/10.1585/jspf.80.358}
\item \textbf{Electron density measurement in a plasma by three-wave interferometry} (1991). L. P. Says, A. Mokhtari. Optics Communications. \url{https://doi.org/10.1016/0030-4018(91)90187-i}
\item \textbf{Measurement of two-dimensional electron density profile in a low current spark using interferometry} (2005). A. S. Patra, T. D. Phukan, A. Khare. IEEE Transactions on Plasma Science. \url{https://doi.org/10.1109/tps.2005.856412}
\item \textbf{Time Resolved Microwave Interferometry Measurement of the Electron Density in a Pulsed 1,3-Butadiene Discharge} (2010). Ashish Kumar Jindal, Lawrence Overzet, Matthew Goeckner. Plasma Chemistry and Plasma Processing. \url{https://doi.org/10.1007/s11090-010-9213-z}
\item \textbf{Electron density measurement in a pulsed-power plasma by FIR laser beam deflection and/or interferometry} (2010). P. Kempkes, F. Mackel, H. Soltwisch. Journal of Physics: Conference Series. \url{https://doi.org/10.1088/1742-6596/227/1/012046}
\item \textbf{Measurement of the Time Varying Electron Density Profile of a Plasma Column by a Modified Laser Interferometry} (1976). Akio Yasuda, Susumu Takeda. Japanese Journal of Applied Physics. \url{https://doi.org/10.1143/jjap.15.675}
\item \textbf{Research on measurement and processing of plasma electron density in Tokamak device} (2013). Shuangbao Shu, Jiarong Luo, Erbing Xue. Journal of Electronic Measurement and Instrument. \url{https://doi.org/10.3724/sp.j.1187.2012.00851}
\item \textbf{Measurement of electron density in atmospheric pressure small-scale plasmas using CO2-laser heterodyne interferometry} (2009). Joon-Young Choi, Nobuhiko Takano, Keiichiro Urabe, Kunihide Tachibana. Plasma Sources Science and Technology. \url{https://doi.org/10.1088/0963-0252/18/3/035013}
\item \textbf{Maxwell-Schrodinger Equation for X-Ray Laser Propagation and Interferometry Measurement of Plasma Electron Density} (2001). Liu Timon Cheng-Yi, Guo Hong, Fu Xi-Quan, Hu Wei, Yu Song. Chinese Physics Letters. \url{https://doi.org/10.1088/0256-307x/18/11/321}
\item \textbf{Measurement of Low Electron Density in Plasmas Using Laser Interferometry} (1965). Herbert Malamud. Review of Scientific Instruments. \url{https://doi.org/10.1063/1.1719375}
\item \textbf{Determination of electron density of laser-produced Y-Ba-Cu-O plasma with laser interferometry} (1992). Shuliang Jiao, Da Liu, Zaiguang Li. SPIE Proceedings (Laser Interferometry IV). \url{https://doi.org/10.1117/12.135349}
\item \textbf{Electron-beam generated air plasma: Measurement of ozone and electron density} (2008). Robert J. Vidmar, Kenneth R. Stalder. 2008 IEEE 35th International Conference on Plasma Science. \url{https://doi.org/10.1109/plasma.2008.4590661}
\item \textbf{Development of continuum interferometry for dynamic measurement of plasma density in wire array z-pinches} (2009). A. Haboub, V. V. Ivanov, G. S. Sarkisov. 2009 IEEE International Conference on Plasma Science. \url{https://doi.org/10.1109/plasma.2009.5227506}
\item \textbf{Electron Density Oscillations in CC-RF Oxygen Plasma Investigated by Gaussian Beam Microwave Interferometry} (2011). Christian Kullig, Kristian Dittmann, Jurgen Meichsner. IEEE Transactions on Plasma Science. \url{https://doi.org/10.1109/tps.2011.2156430}
\item \textbf{Modified relation for laser-plasma electron density measurement using x-ray interferometry} (2003). Hong Guo, Hua Tang, Xiquan Fu, Yunli Qiu, Mingwei Liu, Song Yu, Guangshan Lu, Dongmei Deng. SPIE Proceedings (Laser Resonators and Beam Control VI). \url{https://doi.org/10.1117/12.478980}
\item \textbf{Electron density measurement by differential interferometry} (2006). W. X. Ding, D. L. Brower, B. H. Deng, T. Yates. Review of Scientific Instruments. \url{https://doi.org/10.1063/1.2229217}
\item \textbf{Intensity interferometry for measurement of electron temperature fluctuations in fusion plasmas} (1993). S. Sattler, H. J. Hartfuss. Plasma Physics and Controlled Fusion. \url{https://doi.org/10.1088/0741-3335/35/9/016}
\item \textbf{Microscopic heterodyne interferometry for determination of electron density in high-pressure microplasma} (2014). Keiichiro Urabe, Hitoshi Muneoka, Sven Stauss, Kazuo Terashima. Plasma Sources Science and Technology. \url{https://doi.org/10.1088/0963-0252/23/6/064007}
\item \textbf{Electron Density Determination in an Arc Plasma by Laser Interferometry} (1975). M. Brunger, M. Kock. Zeitschrift fur Naturforschung A. \url{https://doi.org/10.1515/zna-1975-1208}
\item \textbf{Polarimetry-Interferometry Diagnostic for Measurement of Magnetic and Density Perturbations in the High-Temperature Plasma Core} (2011). Liang Lin, W. X. Ding, W. F. Bergerson, D. L. Brower. IEEE Transactions on Plasma Science. \url{https://doi.org/10.1109/tps.2011.2158556}
\item \textbf{Method of electron density measurement by Fabry-Perot interferometry} (2005). H. Hojo, A. Mase. 2005 Joint 30th Int. Conf. Infrared and Millimeter Waves / 13th Int. Conf. Terahertz Electronics. \url{https://doi.org/10.1109/icimw.2005.1572389}
\item \textbf{Direct measurement of the electron density in electron beam irradiated Ar/HCl gas mixtures by time-resolved interferometry} (1990). M. Lando, Z. Rozenberg, M. Rokni. Journal of Applied Physics. \url{https://doi.org/10.1063/1.346485}
\item \textbf{Measurement of a density profile of a hot-electron plasma in RT-1 with three-chord interferometry} (2015). H. Saitoh, Y. Yano, Z. Yoshida, M. Nishiura, J. Morikawa, Y. Kawazura, T. Nogami, M. Yamasaki. Physics of Plasmas. \url{https://doi.org/10.1063/1.4908550}
\item \textbf{Novel Plasma Diagnostic Measurement of Electron Temperature and Electron Density Using Tone Burst Wave} (2021). Takeshi Katahira, Hiromu Kawana, Mikio Ohuchi, Shuichi Sato. IEEE Transactions on Plasma Science. \url{https://doi.org/10.1109/tps.2021.3084138}
\item \textbf{A fiber-coupled dispersion interferometer for density measurements of pulsed power transmission line electron sheaths on Sandia's Z machine} (2022). . Review of Scientific Instruments. \url{https://pubs.aip.org/aip/rsi/article/93/11/113505/2849047/A-fiber-coupled-dispersion-interferometer-for}
\item \textbf{Combining laser interferometry and plasma spectroscopy for spatially resolved high-sensitivity plasma density measurements in discharge capillaries} (2021). . Review of Scientific Instruments. \url{https://pubs.aip.org/aip/rsi/article/92/1/013505/1018688/Combining-laser-interferometry-and-plasma}
\item \textbf{Optical diagnostics for density measurement in high-quality laser-plasma electron accelerators} (2022). . High Power Laser Science and Engineering. \url{https://www.cambridge.org/core/journals/high-power-laser-science-and-engineering/article/optical-diagnostics-for-density-measurement-in-highquality-laserplasma-electron-accelerators/49FDF23E13EF5F480656D26114D7717E}
\item \textbf{Plasma diagnostic techniques based on terahertz radiation} (2025). . Applied Physics Reviews. \url{https://pubs.aip.org/aip/apr/article/12/3/031314/3358497/Plasma-diagnostic-techniques-based-on-terahertz}
\item \textbf{Wavefront-sensor-based electron density measurements for laser-plasma accelerators} (2010). OSTI technical report. OSTI / technical report. \url{https://www.osti.gov/biblio/986199}
\item \textbf{Interferometry in Femtosecond Laser Plasma Diagnostics} (2015). . Physics Procedia (ScienceDirect). \url{https://www.sciencedirect.com/science/article/pii/S1875389215011293}
\item \textbf{Laser interferometry and photon scattering in plasma diagnostics}. . IEEE Xplore. \url{https://ieeexplore.ieee.org/iel6/8/25443/01138861.pdf}
\end{enumerate}

\section*{16\quad Hall thruster erosion and lifetime \textnormal{(\normalsize 35 references)}}
Studies of Hall-thruster channel erosion connect ion bombardment to the dominant life-limiting wear mechanism. They document measurement campaigns, models and mitigation strategies. They define the first commercial pull for a non-intrusive flux diagnostic.
\begin{enumerate}[leftmargin=1.4em]
\item \textbf{2D (r–z) numerical simulations of the plasma and channel erosion in a 100 kW class nested Hall thruster} (2018). Ioannis G. Mikellides, Alejandro Lopez Ortega. Plasma Sources Science and Technology. \url{https://doi.org/10.1088/1361-6595/aacacb}
\item \textbf{Plasma flow and plasma–wall transition in Hall thruster channel} (2001). M. Keidar, I. D. Boyd. Physics of Plasmas. \url{https://doi.org/10.1063/1.1421370}
\item \textbf{Kinetic particle simulation of discharge and wall erosion of a Hall thruster} (2013). Shinatora Cho, Kimiya Komurasaki, others. Physics of Plasmas. \url{https://doi.org/10.1063/1.4810798}
\item \textbf{Development of a cavity ring-down spectroscopy sensor for boron nitride sputter erosion in Hall thrusters} (2010). Lei Tao, Naoji Yamamoto, Alec D. Gallimore, Azer P. Yalin. arXiv:1005.0592. \url{http://arxiv.org/abs/1005.0592}
\item \textbf{Sheath Analysis of Conducting Channel Walls in Anode-Layer Hall Thruster with Acceleration Zone Shifted Downstream} (2023). Rei Kawashima. arXiv:2312.09569. \url{http://arxiv.org/abs/2312.09569}
\item \textbf{Use of emission spectroscopy for real-time assessment of relative wall erosion rate of BHT-200 Hall thruster for various regimes of operation} (2010). Murat Celik, Oleg Batishchev. Vacuum. \url{https://doi.org/10.1016/j.vacuum.2010.01.031}
\item \textbf{Spectroscopic Investigation of a Hall Thruster Ceramic Acceleration Channel Erosion Rate} (2006). Evgeny Dyshlyuk, Oleg Gorshkov. 42nd AIAA/ASME/SAE/ASEE Joint Propulsion Conference \textbackslash{}\& Exhibit. \url{https://doi.org/10.2514/6.2006-4660}
\item \textbf{Extending service life of Hall thrusters: recent progress and future challenges} (2019). . Reviews of Modern Plasma Physics. \url{https://link.springer.com/article/10.1007/s41614-019-0036-y}
\item \textbf{Volumetric erosion rate reduction of Hall thruster channel wall during ion sputtering process} (2007). Daren Yu, Yuquan Li, others. Journal of Physics D: Applied Physics. \url{https://doi.org/10.1088/0022-3727/40/8/017}
\item \textbf{Simulating Plasma-Induced Hall Thruster Wall Erosion With a Two-Dimensional Hybrid Model} (2007). Emmanuelle Sommier, Michelle K. Scharfe, others. IEEE Transactions on Plasma Science. \url{https://doi.org/10.1109/tps.2007.905943}
\item \textbf{Wall Erosion in 2D Hall Thruster Simulations} (2006). Emmanuelle Sommier, Michelle Allis, others. 42nd AIAA/ASME/SAE/ASEE Joint Propulsion Conference \textbackslash{}\& Exhibit. \url{https://doi.org/10.2514/6.2006-4656}
\item \textbf{An Investigation of Factors Involved in Hall Thruster Wall Erosion Modeling} (2006). John Yim, Michael Keidar, Iain D. Boyd. 42nd AIAA/ASME/SAE/ASEE Joint Propulsion Conference \textbackslash{}\& Exhibit. \url{https://doi.org/10.2514/6.2006-4657}
\item \textbf{Numerical prediction of wall erosion on a Hall thruster} (2008). Hiroyuki Koizumi, Kimiya Komurasaki, others. Vacuum. \url{https://doi.org/10.1016/j.vacuum.2008.03.096}
\item \textbf{Multilayer Coating Method for Investigating Channel-Wall Erosion in a Hall Thruster} (2013). Shinatora Cho, Shigeru Yokota, others. Journal of Propulsion and Power. \url{https://doi.org/10.2514/1.b34555}
\item \textbf{Hall Thruster Channel Wall Erosion Rate Measurement Method Using Multilayer Coating Chip} (2010). Shinatora Cho, Shigeru Yokota, others. 46th AIAA/ASME/SAE/ASEE Joint Propulsion Conference \textbackslash{}\& Exhibit. \url{https://doi.org/10.2514/6.2010-6697}
\item \textbf{Channel Wall Erosion Modeling of a SPT-Type Hall Thruster} (2012). Shinatora Cho, Shigeru Yokota, et al.. Transactions of the Japan Society for Aeronautical and Space Sciences, Aerospace Technology Japan. \url{https://doi.org/10.2322/tastj.10.pb_25}
\item \textbf{Predicting Hall Thruster Operational Lifetime} (2004). David Manzella, John Yim, Iain Boyd. 40th AIAA/ASME/SAE/ASEE Joint Propulsion Conference \textbackslash{}\& Exhibit. \url{https://doi.org/10.2514/6.2004-3953}
\item \textbf{Predicting Hall thruster operational lifetime with computational models} (2012). Yongjun Choi, S. Mahalingam, others. 2012 IEEE Aerospace Conference. \url{https://doi.org/10.1109/aero.2012.6187323}
\item \textbf{Influence of the erosion of the discharge channel wall on the efficiency of a stationary plasma thruster} (2012). A. A. Shagaida, O. A. Gorshkov. Technical Physics. \url{https://doi.org/10.1134/s1063784212080221}
\item \textbf{Influence of neutral gas supply position on wall erosion of Hall thruster studied by particle-in-cell simulation} (2023). Xifeng Cao, Guojun Xia, et al.. Plasma Science and Technology. \url{https://doi.org/10.1088/2058-6272/acaf25}
\item \textbf{Investigation of discharge channel wall material influence on lifetime of Hall effect thruster with high specific impulse} (2011). V. V. Abashkin, M. B. Belikov, others. Progress in Propulsion Physics. \url{https://doi.org/10.1051/eucass/201102617}
\item \textbf{Ion-induced wall erosion of Hall-effect thruster on magnetic field aligned mesh using hybrid model} (2026). Gwanyong Jung, Hong-Gye Sung. Aerospace Science and Technology. \url{https://doi.org/10.1016/j.ast.2026.112264}
\item \textbf{Reduced-order particle-in-cell simulations of a high-power magnetically shielded Hall thruster} (2023). Maryam Reza, Farbod Faraji, Aaron Knoll, Antonio Piragino, Tommaso Andreussi, Tommaso Misuri. arXiv:2304.06563. \url{http://arxiv.org/abs/2304.06563}
\item \textbf{Near-Wall Pathways of Anomalous Electron Transport in Hall Thrusters Revealed by 3D PIC Simulations} (2026). Zhe Liu, Zhongping Zhao, Yinjian Zhao. arXiv:2603.14849. \url{http://arxiv.org/abs/2603.14849}
\item \textbf{Wall Ion Loss Reduction by Acceleration Zone Shifting in Anode-Layer Hall Thruster} (2022). Rei Kawashima, Yushi Hamada, Shu Kawabata, Kimiya Komurasaki, Hiroyuki Koizumi. arXiv:2202.02921. \url{http://arxiv.org/abs/2202.02921}
\item \textbf{Self-consistent kinetic simulations of SPT and HEMP thrusters including the near-field plume region} (2009). K. Matyash, R. Schneider, A. Mutzke, O. Kalentev, F. Taccogna, N. Koch, M. Schirra. arXiv:0912.0470. \url{http://arxiv.org/abs/0912.0470}
\item \textbf{Plasma Diagnostic and Performance of a Permanent Magnet Hall Thruster} (2004). Jose Leonardo Ferreira, Joao Henrique Campos De Souza, Israel Da Silveira Rego, Ivan Soares Ferreira. arXiv:physics/0410170. \url{http://arxiv.org/abs/physics/0410170}
\item \textbf{A comparison of Fourier and POD mode decomposition methods for high-speed Hall thruster video} (2022). J. W. Brooks, M. S. McDonald, A. A. Kaptanoglu. arXiv:2205.14207. \url{http://arxiv.org/abs/2205.14207}
\item \textbf{Assessment of Pole Erosion in a Magnetically Shielded Hall Thruster} (2014). Ioannis G. Mikellides, Alejandro Lopez Ortega, others. 50th AIAA/ASME/SAE/ASEE Joint Propulsion Conference (AIAA Propulsion and Energy Forum). \url{https://arc.aiaa.org/doi/10.2514/6.2014-3897}
\item \textbf{Modeling of Hall thruster lifetime and erosion mechanisms (IEPC-2007-250)} (2007). . 30th International Electric Propulsion Conference (IEPC). \url{https://electricrocket.org/IEPC/IEPC-2007-250.pdf}
\item \textbf{In-Situ Diagnostic of Channel Erosion in Hall Thrusters} (2024). Ivan Romadanov, others. International Electric Propulsion Conference (PPPL HTX). \url{https://htx.pppl.gov/publication/Conference/romadanov-et-al-2024-in-situ-diagnostic-of-channel-erosion-in-hall-thrusters.pdf}
\item \textbf{Wear Testing of a Magnetically Shielded Hall Thruster at 2000 s Specific Impulse (IEPC-2015-155)} (2015). . 34th International Electric Propulsion Conference (IEPC). \url{https://electricrocket.org/IEPC/IEPC-2015-155_ISTS-2015-b-155.pdf}
\item \textbf{Investigation of Hall Effect Thruster Channel Wall Erosion Mechanisms (PhD dissertation)} (2016). Aaron M. Schinder. Georgia Institute of Technology (HPEPL). \url{https://hpepl.ae.gatech.edu/papers/Schinder_Dissertation.pdf}
\item \textbf{Channel length and wall recombination effects in the Hall thruster} (2000). A. Cohen-Zur, A. Gany, others. 36th AIAA/ASME/SAE/ASEE Joint Propulsion Conference and Exhibit. \url{https://doi.org/10.2514/6.2000-3654}
\item \textbf{Coaxial Hall Thruster with Diamond Inner Channel Wall} (2006). Nicolas Gascon, Scott Crawford, others. 42nd AIAA/ASME/SAE/ASEE Joint Propulsion Conference \textbackslash{}\& Exhibit. \url{https://doi.org/10.2514/6.2006-4995}
\end{enumerate}

\section*{17\quad Ion and gridded thruster diagnostics \textnormal{(\normalsize 30 references)}}
Gridded ion-thruster literature examines grid erosion, beam optics and discharge diagnostics. Ion energetics govern grid life and performance. This is a second propulsion application domain for the framework.
\begin{enumerate}[leftmargin=1.4em]
\item \textbf{Erosion rate diagnostics in ion thrusters using laser-induced fluorescence} (1997). . Journal of Propulsion and Power. \url{https://doi.org/10.2514/3.23631}
\item \textbf{In situ erosion measurement tools for electric propulsion thrusters: triangular laser head and telemicroscope} (2022). . EPJ Techniques and Instrumentation. \url{https://link.springer.com/article/10.1140/epjti/s40485-022-00076-z}
\item \textbf{A review of the impact of ground test-related facility effects on gridded ion thruster operation and performance} (2024). . Physics of Plasmas. \url{https://pubs.aip.org/aip/pop/article/31/3/030501/3268073/A-review-of-the-impact-of-ground-test-related}
\item \textbf{Ion Optics Model for Rapid Assessment of Gridded Ion Thruster Mission Performance and Life} (2024). . AIAA SciTech Forum. \url{https://arc.aiaa.org/doi/10.2514/6.2024-1550}
\item \textbf{Modeling grid erosion in the NEXT ion thruster using the CEX2D and CEX3D codes} (2023). . Journal of Electric Propulsion. \url{https://link.springer.com/article/10.1007/s44205-023-00043-0}
\item \textbf{Ion thruster accelerator grid erosion mechanism under extreme conditions of cylindrical erosion and chamfer erosion} (2023). . Chinese Journal of Aeronautics. \url{https://www.sciencedirect.com/science/article/pii/S1000936123000973}
\item \textbf{Real-time plasma diagnostic tool for the investigation of azimuthal rotating instabilities in Hall Effect thrusters} (2024). . Journal of Electric Propulsion. \url{https://link.springer.com/article/10.1007/s44205-024-00098-7}
\item \textbf{Stationary plasma thruster ion velocity distribution} (1994). . AIAA/ASME/SAE/ASEE Joint Propulsion Conference. \url{https://doi.org/10.2514/6.1994-3141}
\item \textbf{Perspectives on Physics of ExB Discharges Relevant to Plasma Propulsion and Similar Technologies} (2020). . arXiv:2007.09194. \url{https://arxiv.org/abs/2007.09194}
\item \textbf{Incoherent Thomson scattering diagnostic development for plasma propulsion investigations} (2018). . conference/preprint. \url{https://www.researchgate.net/publication/325120693_Incoherent_Thomson_scattering_diagnostic_development_for_plasma_propulsion_investigations}
\item \textbf{Application of Coherent Microwave Scattering and Multiphoton Ionization for Diagnostics of Electric Propulsion Systems} (2023). Adam R. Patel, Sashin L. B. Karunarathne, Nicholas Babusis, Alexey Shashurin. arXiv:2301.13023. \url{https://arxiv.org/abs/2301.13023}
\item \textbf{New Methods for Computer Tomography Based Ion Thruster Diagnostics and Simulation} (2024). Jörn Krenzer, Felix Reichenbach, Jochen Schein. arXiv:2412.04214. \url{https://arxiv.org/abs/2412.04214}
\item \textbf{CT-imaging in Electrostatic Thruster Ion-Optics} (2024). Jörn Krenzer, Felix Reichenbach, Jochen Schein. arXiv:2412.03426. \url{https://arxiv.org/abs/2412.03426}
\item \textbf{On the microscopic mechanism of ion-extraction of a gridded ion propulsion thruster} (2013). Danny Kirmse. arXiv:1312.5629. \url{https://arxiv.org/abs/1312.5629}
\item \textbf{Analysis and Particle-in-Cell Simulation of Gridded ICRH Plasma Thruster} (2022). Ishaan Mishra. arXiv:2211.09364. \url{https://arxiv.org/abs/2211.09364}
\item \textbf{Understanding In-Chamber Plasma Behavior Using a Dimensionally Scaled Gridded Ion Thruster in Three-Dimensional Kinetic Particle-in-Cell Simulations} (2025). Gyuha Lim, Deborah Levin. arXiv:2511.14036. \url{https://arxiv.org/abs/2511.14036}
\item \textbf{Three-Dimensional Kinetic Simulation of an Ion Thruster Plume with Carbon Backsputtering in a Vacuum Chamber} (2023). Keita Nishii, Deborah A. Levin. arXiv:2303.07496. \url{https://arxiv.org/abs/2303.07496}
\item \textbf{Kinetic Simulation of Ion Thruster Plume Neutralization in a Vacuum Chamber} (2023). Keita Nishii, Deborah Levin. arXiv:2306.16693. \url{https://arxiv.org/abs/2306.16693}
\item \textbf{Electrospray Propulsion Time-of-Flight Secondary Ion Mass Spectrometry Diagnostic} (2024). Giuliana Caramella Hofheins, Zach Ulibarri, Elaine M. Petro. arXiv:2411.09474. \url{https://arxiv.org/abs/2411.09474}
\item \textbf{Dormant Cathode Erosion in a Multiple-Cathode Gridded Ion Thruster} (2008). Joshua L. Rovey, Alec D. Gallimore. Journal of Propulsion and Power. \url{https://doi.org/10.2514/1.37031}
\item \textbf{Modeling and Analysis of Sputtering and Erosion in an Ion Gridded Thruster} (2026). Cesar Cardenas, Luca Massa, Carlos A Collazos. AIAA AVIATION 2026 Forum. \url{https://doi.org/10.2514/6.2026-4258}
\item \textbf{Internal erosion rates of a 10-kW xenon ion thruster} (1988). Vincent Rawlin. 24th Joint Propulsion Conference. \url{https://doi.org/10.2514/6.1988-2912}
\item \textbf{Effect of facility background gases on internal erosion of the 30-cm Hg ion thruster} (1978). V. Rawlin, M. Mantenieks. 13th International Electric Propulsion Conference. \url{https://doi.org/10.2514/6.1978-665}
\item \textbf{Sputter erosion and deposition in the discharge chamber of a small mercury ion thruster} (1973). J. Power. 10th Electric Propulsion Conference. \url{https://doi.org/10.2514/6.1973-1109}
\item \textbf{Plasma diagnostics of the RF-ion thruster RIT-10} (1972). R. Walther, M. Schaefer, J. Freisinger. 9th Electric Propulsion Conference. \url{https://doi.org/10.2514/6.1972-472}
\item \textbf{Direct thrust measurements and beam diagnostics on the ion thruster ESKA 18-P} (1972). H. Junge, U. Sprengel. 9th Electric Propulsion Conference. \url{https://doi.org/10.2514/6.1972-433}
\item \textbf{Advanced thruster diagnostics development at Aerospace - A mass/velocity analyzer for electric thruster plumes} (1987). D. Lichtin, J. Pollard, R. Cohen. 19th International Electric Propulsion Conference. \url{https://doi.org/10.2514/6.1987-1064}
\item \textbf{Modeling Neutral Densities Downstream of a Gridded Ion Thruster} (2011). George C. Soulas. Journal of Propulsion and Power. \url{https://doi.org/10.2514/1.b34094}
\item \textbf{Neutralizer-free gridded ion thruster} (2014). Dmytro Rafalskyi, Ane Aanesland. 50th AIAA/ASME/SAE/ASEE Joint Propulsion Conference. \url{https://doi.org/10.2514/6.2014-3423}
\item \textbf{Auxiliary propulsion thruster performance with ion machined accelerator grids} (1975). W. Hudson. 11th Electric Propulsion Conference. \url{https://doi.org/10.2514/6.1975-425}
\end{enumerate}

\section*{18\quad Electric propulsion plasma diagnostics \textnormal{(\normalsize 24 references)}}
This literature surveys the diagnostic toolset applied to electric-propulsion plasmas and plumes, including optical and probe methods. It situates the proposed framework within established practice. It identifies where non-intrusive fusion adds value.
\begin{enumerate}[leftmargin=1.4em]
\item \textbf{Ion diagnostics of a discharge in crossed electric and magnetic fields for electric propulsion} (2009). S. Mazouffre. Plasma Sources Science and Technology. \url{https://doi.org/10.1088/0963-0252/18/3/034022}
\item \textbf{Introduction: Recommended Practices for Performance and Diagnostic Measurement in Electric Propulsion} (2018). . Journal of Propulsion and Power. \url{https://doi.org/10.2514/1.B36751}
\item \textbf{Recommended Practice for Use of Emissive Probes in Electric Propulsion Testing} (2017). . Journal of Propulsion and Power. \url{https://arc.aiaa.org/doi/abs/10.2514/1.B35697}
\item \textbf{Incoherent Thomson scattering measurements of electron properties in a conventional and magnetically-shielded Hall thruster} (2020). . Plasma Sources Science and Technology. \url{https://doi.org/10.1088/1361-6595/ab6c42}
\item \textbf{Laser Thomson scattering measurements indicate non-isothermal magnetic field lines in magnetically shielded Hall effect thrusters} (2024). . Physics of Plasmas. \url{https://pubs.aip.org/aip/pop/article/31/11/113106/3320976}
\item \textbf{Laser-induced fluorescence detection of the 4d 4D7/2 ... state for plasma diagnostics} (2023). . arXiv:2311.02479. \url{https://arxiv.org/abs/2311.02479}
\item \textbf{Electron Density Measurement Inside a Hall Thruster Using Microwave Interferometry} (2021). . Journal of Propulsion and Power. \url{https://doi.org/10.2514/1.B38163}
\item \textbf{An in-flight plasma diagnostic package for spacecraft with electric propulsion} (2021). . EPJ Techniques and Instrumentation. \url{https://doi.org/10.1140/epjti/s40485-021-00073-8}
\item \textbf{Electrostatic probe apparatus for measurements in the near-anode region of Hall thrusters} (2004). . Review of Scientific Instruments. \url{https://pubs.aip.org/aip/rsi/article-abstract/75/5/1255/784030}
\item \textbf{A high sensitivity momentum flux measuring instrument for plasma thruster exhausts and diffusive plasmas} (2009). . Review of Scientific Instruments. \url{https://pubmed.ncbi.nlm.nih.gov/19485509/}
\item \textbf{Helicon-type radiofrequency plasma thrusters and magnetic plasma nozzles} (2019). . Reviews of Modern Plasma Physics. \url{https://doi.org/10.1007/s41614-019-0024-2}
\item \textbf{Perspectives on Physics of E×B Discharges Relevant to Plasma Propulsion and Similar Technologies} (2020). . arXiv:2007.09194. \url{https://arxiv.org/abs/2007.09194}
\item \textbf{AI-assisted Advanced Propellant Development for Electric Propulsion} (2025). Angel Pan Du, Miguel Arana-Catania, Enric Grustan Gutiérrez. arXiv:2509.26567. \url{https://arxiv.org/abs/2509.26567}
\item \textbf{Quantum-optical spectroscopy for plasma electric field measurements and diagnostics} (2017). David A. Anderson, Georg Raithel, Matthew Simons, Christopher L. Holloway. arXiv:1712.08717. \url{https://arxiv.org/abs/1712.08717}
\item \textbf{Tomography for Plasma Imaging: a Unifying Framework for Bayesian Inference} (2025). D. Hamm, C. Theiler, M. Simeoni, B. P. Duval, T. Debarre, L. Simons, J. R. Queralt. arXiv:2506.20232. \url{https://arxiv.org/abs/2506.20232}
\item \textbf{First Breakthrough for Future Air-Breathing Magneto-Plasma Propulsion Systems} (2016). Berkant Goksel, Igor Mashek. arXiv:1609.04054. \url{https://arxiv.org/abs/1609.04054}
\item \textbf{Development of an Improved Probe for Diagnostics of the Charge Composition of Ions in Electric Propulsion Plasma} (2026). D. A. Maistrenko. Technical Physics Letters. \url{https://doi.org/10.1134/s1063785026700343}
\item \textbf{Diagnostics Of Plasma Propulsion Devices} (1994). M. A. Cappelli. 1994 IEEE 21st International Conference on Plasma Sciences (ICOPS). \url{https://doi.org/10.1109/plasma.1994.588875}
\item \textbf{Principles of spectroscopic diagnostics of plasma with oscillating electric fields} (1986). E. A. Oks. Measurement Techniques. \url{https://doi.org/10.1007/bf00863283}
\item \textbf{Probe for the measurement of ion velocity distributions in low density plasma beams} (1967). L. Liebing. 6th Electric Propulsion and Plasmadynamics Conference. \url{https://doi.org/10.2514/6.1967-706}
\item \textbf{Diagnostics of hollow cathode high power pulsed plasma thruster} (1976). V. Naso. 12th International Electric Propulsion Conference. \url{https://doi.org/10.2514/6.1976-1005}
\item \textbf{Development of a microwave plasma diagnostic applied to electric propulsion systems}. Federico Boni. \url{https://doi.org/10.70675/0baadd37ze0f4z4102zaf29zdd24929bc822}
\item \textbf{Measurement challenges: electric and hybrid propulsion} (2022). J. Blakesley. NPL. \url{https://doi.org/10.47120/npl.9360}
\item \textbf{Diagnostics of a self-field accelerator with electric probes} (1973). F. Maisenhaelder. 10th Electric Propulsion Conference. \url{https://doi.org/10.2514/6.1973-1095}
\end{enumerate}

\section*{19\quad Plasma etching and reactive ion etching ion energy \textnormal{(\normalsize 26 references)}}
Etching studies tie ion energy and angular distributions at the wafer to anisotropy, selectivity and damage. At advanced nodes, ion-energy control is decisive for yield. This is a large industrial application for in-situ flux measurement.
\begin{enumerate}[leftmargin=1.4em]
\item \textbf{Plasma sheath tailoring by a magnetic field for three-dimensional plasma etching} (2023). E. Jüngling, S. Wilczek, T. Mussenbrock, M. Böke, A. von Keudell. arXiv:2311.08916. \url{http://arxiv.org/abs/2311.08916}
\item \textbf{Electron dynamics and SiO2 etching profile evolution in capacitive Ar/CHF3 discharges driven by sawtooth-tailored voltage waveforms} (2024). Wan Dong, Liu-Qin Song, Yi-Fan Zhang, Li Wang, Yuan-Hong Song, Julian Schulze. arXiv:2411.07839. \url{http://arxiv.org/abs/2411.07839}
\item \textbf{Etching-to-deposition transition in SiO2/Si3N4 using CHxFy ion-based plasma etching: An atomistic study with neural network potentials} (2025). Hyungmin An, Sangmin Oh, Dongheon Lee, Jae-hyeon Ko, Dongyean Oh, Changho Hong, Seungwu Han. arXiv:2508.00327. \url{http://arxiv.org/abs/2508.00327}
\item \textbf{Reduced etch lag and high aspect ratios by deep reactive ion etching (DRIE)} (2021). M. S. Gerlt, N. F. Läubli, M. Manser, B. J. Nelson, J. Dual. arXiv:2104.02763. \url{http://arxiv.org/abs/2104.02763}
\item \textbf{Crystallographic Orientation Dependent Reactive Ion Etch in Single Crystal Diamond} (2020). Ling Xie, Tony X. Zhou, Rainer J. Stöhr, Amir Yacoby. arXiv:2004.03423. \url{http://arxiv.org/abs/2004.03423}
\item \textbf{Nanostructuring SiC by sequential plasma oxidation and reactive ion etching} (2026). Joonas Isometsä, Auguste Bieleviciute, Xiaolong Liu, Ville Vähänissi, Hele Savin. arXiv:2603.16589. \url{http://arxiv.org/abs/2603.16589}
\item \textbf{Numerical and experimental studies of the carbon etching in EUV-induced plasma} (2015). D. I. Astakhov, W. J. Goedheer, C. J. Lee, V. V. Ivanov, V. M. Krivtsun, O. Yakushev, K. N. Koshelev, D. V. Lopaev, F. Bijkerk. arXiv:1507.02705. \url{http://arxiv.org/abs/1507.02705}
\item \textbf{Development of hard masks for reactive ion beam angled etching of diamond} (2022). Cleaven Chia, Bartholomeus Machielse, Amirhassan Shams-Ansari, Marko Loncar. arXiv:2201.06416. \url{http://arxiv.org/abs/2201.06416}
\item \textbf{Towards High Throughput Large Area Metalens Fabrication using UV-Nanoimprint lithography and Bosch Deep Reactive Ion Etching} (2020). Christopher A. Dirdal, Geir Uri Jensen, Hallvard Angelskår, Paul Conrad Vaagen Thrane, Jo Gjessing, Daniel Alfred Ordnung. arXiv:2004.11133. \url{http://arxiv.org/abs/2004.11133}
\item \textbf{Profile Control of a Borosilicate-Glass Groove Formed by Deep Reactive Ion Etching} (2008). T. Akashi, Y. Yoshimura. arXiv:0802.3085. \url{http://arxiv.org/abs/0802.3085}
\item \textbf{A road to hydrogenating graphene by a reactive ion etching plasma} (2011). M. Wojtaszek, N. Tombros, A. Caretta, P. H. M. van Loosdrecht, B. J. van Wees. arXiv:1109.1684. \url{http://arxiv.org/abs/1109.1684}
\item \textbf{Etching of Graphene Devices with a Helium Ion Beam} (2009). M. C. Lemme, D. C. Bell, J. R. Williams, L. A. Stern, B. W. H. Baugher, P. Jarillo-Herrero, C. M. Marcus. arXiv:0905.4409. \url{http://arxiv.org/abs/0905.4409}
\item \textbf{Ion and Radical Characteristics (Mass/Energy Distribution) of a Capacitively Coupled Plasma Source Using Plasma Process Gases (CxFy)} (2021). . Coatings. \url{https://doi.org/10.3390/coatings11080993}
\item \textbf{Effect of Hydrogen Ion Energy in the Process of Reactive Ion Etching of Sn Thin Films by Hydrogen Plasmas} (2021). Mengran Ji, Ryo Nagata. Plasma and Fusion Research. \url{https://doi.org/10.1585/pfr.16.1406003}
\item \textbf{Low-temperature reactive ion etching and microwave plasma etching of silicon} (1988). Shinichi Tachi, Kazunori Tsujimoto. Applied Physics Letters. \url{https://doi.org/10.1063/1.99382}
\item \textbf{Effect of Positive Photoresist on Silicon Etching by Reactive Ion Etching Process} (2010). Muhammad M. Morshed, Stephen M. Daniels. IEEE Transactions on Plasma Science. \url{https://doi.org/10.1109/tps.2010.2046500}
\item \textbf{A novel etching technology with reactive ion etching system for GaAs via-hole etching applications} (2003). C. S. Lin, Y. K. Fang. IEEE Transactions on Semiconductor Manufacturing. \url{https://doi.org/10.1109/tsm.2002.807738}
\item \textbf{Plasma Characterization in Chlorine-Based Reactive Ion Beam Etching and Chemically Assisted Ion Beam Etching} (1998). Akihiro Matsutani, Fumio Koyama. Japanese Journal of Applied Physics. \url{https://doi.org/10.1143/jjap.37.2747}
\item \textbf{RF Matching of a Reactive Ion Etching (RIE) Plasma Reactor} (2005). Roland Gesche. Frequenz. \url{https://doi.org/10.1515/freq.2005.59.3-4.73}
\item \textbf{Kinetics of Reactive Ion Etching of Polymers in an Oxygen Plasma: The Importance of Direct Reactive Ion Etching} (1993). Sandra W. Graham, Christoph Steinbrüchel. MRS Proceedings. \url{https://doi.org/10.1557/proc-334-433}
\item \textbf{Reactive ion etching and plasma etching of tungsten} (1993). P. Verdonck, G. Brasseur. Microelectronic Engineering. \url{https://doi.org/10.1016/0167-9317(93)90084-i}
\item \textbf{Etching characteristics of LiNbO3 crystal by fluorine gas plasma reactive ion etching} (2003). Masashi Tamura, Shinzo Yoshikado. Surface and Coatings Technology. \url{https://doi.org/10.1016/s0257-8972(03)00070-7}
\item \textbf{Laser diagnostic techniques for reactive ion etching: plasma understanding to process control} (1989). . Vacuum. \url{https://doi.org/10.1016/0042-207x(89)90531-9}
\item \textbf{Plasma and Reactive Ion Etching} (2001). J. W. Coburn. Encyclopedia of Materials: Science and Technology. \url{https://doi.org/10.1016/b0-08-043152-6/01244-4}
\item \textbf{Technology of reactive ion etching} (1998). S. Handa. Plasma Etching: Fundamentals and Applications. \url{https://doi.org/10.1093/oso/9780198562870.003.0005}
\item \textbf{Reactive ion etching of PbSe thin films in CH4/H2/Ar plasma atmosphere} (2021). Gang Yang, Binbin Weng. Materials Science in Semiconductor Processing. \url{https://doi.org/10.1016/j.mssp.2020.105596}
\end{enumerate}

\section*{20\quad Semiconductor plasma processing diagnostics \textnormal{(\normalsize 23 references)}}
Process-diagnostics literature covers endpoint detection, sensor integration and control in plasma tools. It frames how a flux diagnostic would enter fab process control. It defines integration and robustness requirements.
\begin{enumerate}[leftmargin=1.4em]
\item \textbf{Control of ion energy distribution at substrates during plasma processing} (2000). Shu-Bao Wang, Amy E. Wendt. Journal of Applied Physics. \url{https://doi.org/10.1063/1.373715}
\item \textbf{Ion energy distribution measurements in rf and pulsed dc plasma discharges} (2012). . Plasma Sources Science and Technology. \url{https://iopscience.iop.org/article/10.1088/0963-0252/21/2/024004}
\item \textbf{Novel Method of Ion Energy and Ion Energy Distribution Function Control at the Substrate during Plasma Processing} (1999). H. Wang, A. E. Wendt. IEEE International Conference on Plasma Science (ICOPS). \url{https://doi.org/10.1109/plasma.1999.829298}
\item \textbf{Independent control of ion density and ion bombardment energy in a dual RF excitation plasma} (1993). H. H. Goto, H.-D. Lowe, T. Ohmi. IEEE Transactions on Semiconductor Manufacturing. \url{https://doi.org/10.1109/66.210658}
\item \textbf{Improving Ion Control for Semiconductor Processing with Repetitive Pulsed Power Phase I} (2022). J. R. Prager, T. M. Ziemba, A. M. Lietz. IEEE International Conference on Plasma Science (ICOPS). \url{https://doi.org/10.1109/icops45751.2022.9813233}
\item \textbf{Optical diagnostics of plasma chemistries and chamber conditions in gate oxide stack etch} (2013). Szetsen Lee, Yu-Chung Tien, Chi-Jung Ni. Materials Science in Semiconductor Processing. \url{https://doi.org/10.1016/j.mssp.2012.10.011}
\item \textbf{Laser-induced fluorescence diagnosis of plasma processing sources} (2003). . Thin Solid Films. \url{https://www.sciencedirect.com/science/article/abs/pii/S0040609003003584}
\item \textbf{Time-resolved ion energy measurements using a retarding potential analyzer for electric propulsion applications} (2024). Austen Thomas, Kristina Lemmer. Review of Scientific Instruments. \url{https://doi.org/10.1063/5.0176167}
\item \textbf{Development of a high dynamic range retarding potential analyzer for electric propulsion plume diagnosis} (2022). . Review of Scientific Instruments. \url{https://pubmed.ncbi.nlm.nih.gov/35922286/}
\item \textbf{Plasma diagnostics for semiconductor processing} (2000). H. M. Anderson. IEEE LEOS Summer Topical Meetings Digest. \url{https://doi.org/10.1109/leosst.2000.869713}
\item \textbf{A proposal for both plasma ion- and electron-temperature diagnostics under simultaneous incidence of particles and x-rays into a semiconductor detector} (2003). T. Numakura, T. Cho, J. Kohagura. Plasma Physics and Controlled Fusion. \url{https://doi.org/10.1088/0741-3335/45/5/321}
\item \textbf{Semiconductor applications of plasma immersion ion implantation} (2003). Paul K. Chu. Plasma Physics and Controlled Fusion. \url{https://doi.org/10.1088/0741-3335/45/5/304}
\item \textbf{Plasma immersion ion implantation for semiconductor processing} (1996). N. W. Cheung. IEEE International Conference on Plasma Science (ICOPS). \url{https://doi.org/10.1109/plasma.1996.550204}
\item \textbf{Diagnostics of Energy Fluxes in Dusty Plasmas} (2017). Horst R. Maurer, Holger Kersten. Plasma Processing of Nanomaterials (book chapter). \url{https://doi.org/10.1201/b11473-12}
\item \textbf{Plasma diagnostics for monitoring and control} (1995). N. St. J. Braithwaite. IEE Colloquium on Applications of Plasma Technology to Surface Processing. \url{https://doi.org/10.1049/ic:19950908}
\item \textbf{Low Energy Plasma Beams for Semiconductor Technology} (1997). H. Oechsner. Plasma Processing of Semiconductors (book). \url{https://doi.org/10.1007/978-94-011-5884-8_26}
\item \textbf{Plasma Immersion Ion Implantation: A Perspective} (1992). Crid Yu, Nathan W. Cheung. Crucial Issues in Semiconductor Materials and Processing Technologies (book). \url{https://doi.org/10.1007/978-94-011-2714-1_25}
\item \textbf{In-Situ Diagnostics For Plasma Processing} (1989). P. Banks, W. Pilz, I. Hussla. SPIE Proceedings - Monitoring and Control of Plasma-Enhanced Processing of Semiconductors. \url{https://doi.org/10.1117/12.951013}
\item \textbf{Optimization of GaAs ECR etching in chemically assisted ion beam process using Cl2/Ar plasma} (2000). S. F. Yoon, T. K. Ng, H. Q. Zheng. Materials Science in Semiconductor Processing. \url{https://doi.org/10.1016/s1369-8001(00)00034-2}
\item \textbf{Control of polysilicon nanowires conductivity by angle-dependent ion implantation} (2018). Shahar Aziza, Alex Ripp, Dror Horvitz. Materials Science in Semiconductor Processing. \url{https://doi.org/10.1016/j.mssp.2017.11.001}
\item \textbf{Insights into the mechanism of in-plasma photo-assisted etching using optical emission spectroscopy} (2016). . Journal of Vacuum Science \textbackslash{}\& Technology A. \url{https://pubs.aip.org/avs/jva/article/34/6/061303/244961/Insights-into-the-mechanism-of-in-plasma-photo}
\item \textbf{Characteristics of High Power Ion Beam Produced in Plasma Focus and Its Application to Semiconductor Processing} (2002). Katsumi Masugata. AIP Conference Proceedings. \url{https://doi.org/10.1063/1.1530884}
\item \textbf{Diagnostics in Thermal Plasma Processing} (1989). P. Fauchais, J. F. Coudert, M. Vardelle. Plasma Diagnostics (book). \url{https://doi.org/10.1016/b978-0-12-067635-4.50012-0}
\end{enumerate}

\section*{21\quad Fusion divertor and plasma-facing components \textnormal{(\normalsize 32 references)}}
Divertor and plasma-facing-component research addresses extreme heat and particle loads and the resulting erosion. Ion energy flux is central to component survival and tritium retention. This is the high-value fusion application domain.
\begin{enumerate}[leftmargin=1.4em]
\item \textbf{Power Balance in the ITER Plasma and Divertor} (1995). D. E. Post, B. Braams, J. Mandrekas, W. Stacey, N. Putvinskaya. arXiv:plasm-ph/9512001. \url{https://arxiv.org/abs/plasm-ph/9512001}
\item \textbf{Inferring the scrape-off layer heat flux width in a divertor with a low degree of axisymmetry} (2024). C. Marsden, X. Zhang, M. Moscheni, T. K. Gray, E. Vekshina, A. Rengle, A. Scarabosio, M. Sertoli, M. Romanelli, the ST40 team. arXiv:2406.16471. \url{https://arxiv.org/abs/2406.16471}
\item \textbf{Runaway electron-induced plasma facing component damage in tokamaks} (2025). S. Ratynskaia, M. Hoelzl, E. Nardon, P. Aleynikov, F. J. Artola, others. arXiv:2506.10411. \url{https://arxiv.org/abs/2506.10411}
\item \textbf{Designing tungsten armoured plasma facing components to pulsed heat loads in magnetic fusion machines} (2024). R. Mitteau, M. Diez, M. Firdaouss. arXiv:2410.13269. \url{https://arxiv.org/abs/2410.13269}
\item \textbf{Gyrokinetic projection of the divertor heat-flux width from present tokamaks to ITER} (2017). C. S. Chang, S. Ku, A. Loarte, V. Parail, F. Kochl, M. Romanelli, R. Maingi, J.-W. Ahn, T. Gray, J. Hughes, B. LaBombard, T. Leonard, M. Makowski, J. Terry. arXiv:1701.05507. \url{https://arxiv.org/abs/1701.05507}
\item \textbf{Fast Dynamic 1D Simulation of Divertor Plasmas with Neural PDE Surrogates} (2023). Yoeri Poels, Gijs Derks, Egbert Westerhof, Koen Minartz, Sven Wiesen, Vlado Menkovski. arXiv:2305.18944. \url{https://arxiv.org/abs/2305.18944}
\item \textbf{Heat and particle flux detachment with stable plasma conditions in the Wendelstein 7-X stellarator fusion experiment} (2020). M. Jakubowski, R. Konig, O. Schmitz, Y. Feng, M. Krychowiak, others. arXiv:2001.02067. \url{https://arxiv.org/abs/2001.02067}
\item \textbf{First heat flux estimation in the lower divertor of WEST with embedded thermal measurements} (2019). J. Gaspar, Y. Corre, M. Firdaouss, J. Gardarein, J. Gerardin, J. Gunn, M. Houry, G. Laffont, T. Loarer, M. Missirlian, J. Morales, P. Moreau, C. Pocheau, E. .... arXiv:1902.10969. \url{https://arxiv.org/abs/1902.10969}
\item \textbf{Studies of Plasma Detachment Using a One Dimensional Model for Divertor Operation} (1995). R. A. Vesey, D. E. Post, G. Bateman. arXiv:plasm-ph/9507001. \url{https://arxiv.org/abs/plasm-ph/9507001}
\item \textbf{Power Deposition on Tokamak Plasma-Facing Components} (2014). Wayne Arter, Valeria Riccardo, Geoff Fishpool. arXiv:1403.7142. \url{https://arxiv.org/abs/1403.7142}
\item \textbf{Divertor heat fluxes and profiles during mitigated and unmitigated Edge Localised Modes (ELMs) on the Mega Amp Spherical Tokamak (MAST)} (2013). A. J. Thornton, A. Kirk, I. T. Chapman, J. R. Harrison, the MAST Team. arXiv:1306.6499. \url{https://arxiv.org/abs/1306.6499}
\item \textbf{Heat flux deposition on plasma-facing components using a convective model with ripple and Shafranov shift} (1999). R. Mitteau, A. Moal, J. Schlosser, D. Guilhem. arXiv:2102.11599. \url{https://arxiv.org/abs/2102.11599}
\item \textbf{Electromagnetic turbulence simulation of tokamak edge plasma dynamics and divertor heat load during thermal quench} (2023). Ben Zhu, Xue-qiao Xu, Xian-Zhu Tang. arXiv:2302.06859. \url{https://arxiv.org/abs/2302.06859}
\item \textbf{Thomson scattering diagnostics at the Globus-M2 tokamak} (2023). . arXiv:2311.18723. \url{https://arxiv.org/abs/2311.18723}
\item \textbf{Experimental observations of bifurcated power decay lengths in the near Scrape-Off Layer of tokamak plasmas} (2024). . arXiv:2405.14485. \url{https://arxiv.org/abs/2405.14485}
\item \textbf{Kinetic modeling of divertor heat load fluxes in the Alcator C-Mod and DIII-D tokamaks} (2015). . arXiv:1506.07899. \url{https://arxiv.org/abs/1506.07899}
\item \textbf{Conceptual design study for heat exhaust management in the ARC fusion pilot plant} (2018). . arXiv:1809.10555. \url{https://arxiv.org/abs/1809.10555}
\item \textbf{Prediction of Critical Heat Flux in Water-Cooled Plasma Facing Components Using Computational Fluid Dynamics} (2011). Dennis Youchison. Fusion Science \textbackslash{}\& Technology. \url{https://doi.org/10.13182/fst10-294}
\item \textbf{Analysis of heat transfer and erosion effects on ITER divertor plasma facing components induced by slow high-power transients}. G. Federici, A. R. Raffray, H. D. Pacher. Proceedings of 16th International Symposium on Fusion Engineering. \url{https://doi.org/10.1109/fusion.1995.534256}
\item \textbf{High heat flux issues for plasma-facing components in fusion reactors} (1993). Robert D. Watson. SPIE Proceedings; High Heat Flux Engineering. \url{https://doi.org/10.1117/12.140495}
\item \textbf{Fracture mechanical analysis of a tungsten monoblock-type plasma-facing component without macroscopic interlayer for high-heat-flux divertor target} (2017). Muyuan Li, Franklin Gallay, Marianne Richou. Fusion Engineering and Design. \url{https://doi.org/10.1016/j.fusengdes.2017.09.002}
\item \textbf{Comparison between FEM and high heat flux thermal fatigue testing results of ITER divertor plasma facing mock-ups} (2014). F. Crescenzi, S. Roccella, E. Visca. Fusion Engineering and Design. \url{https://doi.org/10.1016/j.fusengdes.2014.04.012}
\item \textbf{Development of tungsten armored high heat flux plasma facing components for ITER like divertor application} (2019). K. Premjit Singh, S. S. Khirwadkar, Nikunj Patel. Fusion Engineering and Design. \url{https://doi.org/10.1016/j.fusengdes.2019.03.170}
\item \textbf{High Heat Flux Performance of Plasma Facing Materials and Components under Service Conditions in Future Fusion Reactors} (2010). Jochen Linke. Fusion Science and Technology. \url{https://doi.org/10.13182/fst10-a9420}
\item \textbf{SIRHEX - A new experimental facility for high heat flux testing of plasma facing components} (2015). Andre Kunze, Bradut-Eugen Ghidersa, Flavia Bonelli. Fusion Engineering and Design. \url{https://doi.org/10.1016/j.fusengdes.2015.01.026}
\item \textbf{Erosion from Liquid-Metal Plasma-Facing Components in Future Fusion Devices} (2002). Jean P. Allain. AIP Conference Proceedings. \url{https://doi.org/10.1063/1.1516316}
\item \textbf{High heat flux testing of divertor plasma facing materials and components using the HHF test facility at IPR} (2016). Yashashri Patil, S. S. Khirwadkar, Sunil Belsare. Physica Scripta. \url{https://doi.org/10.1088/0031-8949/t167/1/014020}
\item \textbf{Simultaneous Identification of Thermophysical Properties and Surface Heat Flux on Plasma Facing Components Inside the JET Fusion Reactor} (2010). Fabrice Rigollet, Jean-Laurent Gardarein, Yann Corre. 14th International Heat Transfer Conference. \url{https://doi.org/10.1115/ihtc14-23254}
\item \textbf{A simple model to evaluate the heat flux onto the EU-DEMO plasma facing components during the current quench} (2025). G. Pautasso, M. Siccinio, F. Maviglia. Nuclear Fusion. \url{https://doi.org/10.1088/1741-4326/add675}
\item \textbf{Stability and erosion of melt layers developed on plasma facing components of tokamaks} (2014). G. Miloshevsky, A. Hassanein. Nuclear Fusion. \url{https://doi.org/10.1088/0029-5515/54/4/043016}
\item \textbf{Two-dimensional collisional particle model of the divertor sheath with electron emissive walls} (2023). . Nuclear Fusion. \url{https://iopscience.iop.org/article/10.1088/1741-4326/acdf70}
\item \textbf{Recent DIII-D progress toward validating models of tungsten erosion, re-deposition, and migration for application to next-step fusion devices} (2023). . Materials Research Express. \url{https://iopscience.iop.org/article/10.1088/2053-1591/ad0f43}
\end{enumerate}

\section*{22\quad Tokamak edge and scrape-off-layer diagnostics \textnormal{(\normalsize 37 references)}}
Edge and scrape-off-layer diagnostics characterize the plasma delivering flux to fusion walls. The literature covers probe, spectroscopic and scattering measurements at the edge. It provides the plasma context and cross-checks for the framework.
\begin{enumerate}[leftmargin=1.4em]
\item \textbf{Plasma Turbulence in the Scrape-off Layer of the ISTTOK Tokamak} (2016). Rogerio Jorge, Paolo Ricci, Federico D. Halpern, Nuno F. Loureiro, Carlos Silva. arXiv:1606.09538. \url{https://arxiv.org/abs/1606.09538}
\item \textbf{Merging dynamics of plasma blobs in the Scrape-off Layer of a tokamak} (2024). Souvik Mondal, N. Bisai, Abhijit Sen, Indranil Bandyopadhyay. arXiv:2409.15064. \url{https://arxiv.org/abs/2409.15064}
\item \textbf{Formation of convective cells in the scrape-off layer of the CASTOR tokamak} (2004). J. Stockel, P. Devynck, J. Gunn, E. Martines, G. Bonhomme, G. Van Oost, M. Hron, J. Adamek, others. arXiv:physics/0410208. \url{https://arxiv.org/abs/physics/0410208}
\item \textbf{Relaxation of potential, flows, and density in the edge plasma of CASTOR tokamak} (2004). Martin Hron, V. Weinzettl, E. Dufkova, C. Hidalgo, I. Duran, J. Stoeckel. arXiv:physics/0410225. \url{https://arxiv.org/abs/physics/0410225}
\item \textbf{Numerical investigation of Scrape Off Layer anomalous particle transport for MAST parameters} (2013). F. Militello, W. Fundamenski, V. Naulin, A. H. Nielsen. arXiv:1307.2391. \url{https://arxiv.org/abs/1307.2391}
\item \textbf{Experimental and numerical characterisation of the turbulence in the Scrape-Off Layer of MAST} (2013). F. Militello, P. Tamain, W. Fundamenski, A. Kirk, V. Naulin, A. H. Nielsen, MAST team. arXiv:1305.5064. \url{https://arxiv.org/abs/1305.5064}
\item \textbf{L-mode filament characteristics on MAST as a function of plasma current measured using visible imaging} (2016). A. Kirk, A. J. Thornton, J. R. Harrison, F. Militello, N. R. Walkden. arXiv:1602.03021. \url{https://arxiv.org/abs/1602.03021}
\item \textbf{Scaling laws for electron kinetic effects in tokamak scrape-off layer plasmas} (2022). Dominic Power, Stefan Mijin, Michael Wigram, Fulvio Militello, Robert J. Kingham. arXiv:2208.10862. \url{https://arxiv.org/abs/2208.10862}
\item \textbf{The impact of non-local parallel electron transport on plasma-impurity reaction rates in tokamak scrape-off layer plasmas} (2024). Dominic Power, Stefan Mijin, Kevin Verhaegh, Fulvio Militello, Robert J. Kingham. arXiv:2410.00651. \url{https://arxiv.org/abs/2410.00651}
\item \textbf{Impurity seeding for suppression of the near Scrape-Off Layer heat flux feature in tokamak limited plasmas} (2017). F. Nespoli, B. Labit, I. Furno, C. Theiler, U. Sheikh, C. K. Tsui, J. A. Boedo, TCV team. arXiv:1712.00951. \url{https://arxiv.org/abs/1712.00951}
\item \textbf{Analysis of high-field side plasma instabilities in tokamak edge} (2024). Maxim V. Umansky. arXiv:2409.12329. \url{https://arxiv.org/abs/2409.12329}
\item \textbf{Magnetic Fluctuations in Gyrokinetic Simulations of Tokamak Scrape-Off Layer Turbulence} (2021). Noah Mandell. arXiv:2103.16062. \url{https://arxiv.org/abs/2103.16062}
\item \textbf{Stability of scrape-off layer plasma: a modified Rayleigh-Benard problem} (2018). Fryderyk Wilczynski, David W. Hughes, Sven Van Loo, Wayne Arter, Fulvio Militello. arXiv:1810.05111. \url{https://arxiv.org/abs/1810.05111}
\item \textbf{Velocity shear, turbulent saturation, and steep plasma gradients in the scrape-off layer of inner-wall limited tokamaks} (2016). Federico D. Halpern, Paolo Ricci. arXiv:1606.08564. \url{https://arxiv.org/abs/1606.08564}
\item \textbf{Numerical turbulence simulations of intermittent fluctuations in the scrape-off layer of magnetized plasmas} (2021). Gregor Decristoforo, Audun Theodorsen, John Omotani, Thomas Nicholas, Odd Erik Garcia. arXiv:2102.04723. \url{https://arxiv.org/abs/2102.04723}
\item \textbf{Reduced modeling of scrape-off losses of runaway electrons during tokamak disruptions} (2024). Oskar Vallhagen, Lise Hanebring, Tunde Fulop, Mathias Hoppe, Istvan Pusztai. arXiv:2410.03512. \url{https://arxiv.org/abs/2410.03512}
\item \textbf{Scrape-off layer ion temperature measurements at the divertor target in MAST by retarding field energy analyser} (2013). M. Kocan, others. arXiv:1307.1563. \url{https://arxiv.org/abs/1307.1563}
\item \textbf{Fluctuation statistics in the scrape-off layer of Alcator C-Mod} (2014). O. E. Garcia, S. M. Wold, S. Theodorsen, et al.. arXiv:1410.4114. \url{https://arxiv.org/abs/1410.4114}
\item \textbf{Selected methods of electron- and ion-diagnostics in tokamak scrape-off-layer} (2015). Marek J. Sadowski. Nukleonika. \url{https://doi.org/10.1515/nuka-2015-0039}
\item \textbf{Scrape-off layer tokamak plasma turbulence} (2012). N. Bisai, R. Singh, P. K. Kaw. Physics of Plasmas. \url{https://doi.org/10.1063/1.4718714}
\item \textbf{Simulation of the scrape-off layer region of tokamak devices} (2015). Paolo Ricci. Journal of Plasma Physics. \url{https://doi.org/10.1017/s0022377814001202}
\item \textbf{Plasma turbulence in the scrape-off layer of tokamak devices} (2013). Paolo Ricci, B. N. Rogers. Physics of Plasmas. \url{https://doi.org/10.1063/1.4789551}
\item \textbf{Edge and scrape-off layer tokamak plasma turbulence simulation using two-field fluid model} (2005). Nirmal Bisai, Amita Das, Shishir Deshpande, Ratneshwar Jha, Predhiman Kaw, Abhijit Sen, Raghvendra Singh. Physics of Plasmas. \url{https://doi.org/10.1063/1.1942427}
\item \textbf{Role of neutral gas in scrape-off layer tokamak plasma} (2015). N. Bisai, R. Jha, P. K. Kaw. Physics of Plasmas. \url{https://doi.org/10.1063/1.4913429}
\item \textbf{Evidence for strong secondary electron emission in the tokamak scrape-off layer} (2012). J. P. Gunn. Plasma Physics and Controlled Fusion. \url{https://doi.org/10.1088/0741-3335/54/8/085007}
\item \textbf{Parallel velocity in a narrow scrape-off layer of a tokamak} (2012). V. Rozhansky et al.. Plasma Physics and Controlled Fusion. \url{https://doi.org/10.1088/0741-3335/54/10/102001}
\item \textbf{Simulations of edge and scrape off layer turbulence in mega ampere spherical tokamak plasmas} (2012). F. Militello, W. Fundamenski, V. Naulin, A. H. Nielsen. Plasma Physics and Controlled Fusion. \url{https://doi.org/10.1088/0741-3335/54/9/095011}
\item \textbf{Polarization currents in tokamak scrape-off layer at disruptions and edge-localized modes} (2012). A. V. Nedospasov. Nuclear Fusion. \url{https://doi.org/10.1088/0029-5515/52/3/033002}
\item \textbf{Numerical imaging of edge plasma characteristics in the high-recycling scrape-off layer of the KSTAR tokamak} (2002). Jin-Seok Ko, Ki-Hak Im, Deok-Kyu Kim, Sang Hee Hong. IEEE Transactions on Plasma Science. \url{https://doi.org/10.1109/tps.2002.1003925}
\item \textbf{Ion temperature measurements in the scrape-off layer of the Tore Supra Tokamak}. Martin Kocan. \url{https://doi.org/10.70675/ca34d9b0zab8ez4fd9za838zb098882b9615}
\item \textbf{A new technique for tokamak edge density measurement based on microwave interferometer} (2022). . Review of Scientific Instruments. \url{https://pubs.aip.org/aip/rsi/article/93/1/013502/2848727/A-new-technique-for-tokamak-edge-density}
\item \textbf{Implementation of high-resolution spectroscopy for ion (and electron) temperature measurements of the divertor plasma in the Tokamak a configuration variable} (2022). . Review of Scientific Instruments. \url{https://pubs.aip.org/aip/rsi/article/93/12/123505/2845373/Implementation-of-high-resolution-spectroscopy-for}
\item \textbf{Design of a Thomson scattering diagnostic for the SMall Aspect Ratio Tokamak (SMART)} (2024). . Review of Scientific Instruments. \url{https://pubs.aip.org/aip/rsi/article/95/9/093508/3311136/Design-of-a-Thomson-scattering-diagnostic-for-the}
\item \textbf{Langmuir probe electronics upgrade on the tokamak a configuration variable} (2019). . Review of Scientific Instruments. \url{https://pubs.aip.org/aip/rsi/article/90/8/083502/360055/Langmuir-probe-electronics-upgrade-on-the-tokamak}
\item \textbf{Diagnostics of mobile dust in scrape-off layer plasmas} (2011). . Plasma Physics and Controlled Fusion. \url{https://iopscience.iop.org/article/10.1088/0741-3335/53/7/074009}
\item \textbf{Tokamak plasma diagnosis by electrical probes} (1994). . Plasma Physics and Controlled Fusion. \url{https://iopscience.iop.org/article/10.1088/0741-3335/36/10/002}
\item \textbf{Advanced probes for edge plasma diagnostics on the CASTOR tokamak} (2007). . Journal of Physics: Conference Series. \url{https://iopscience.iop.org/article/10.1088/1742-6596/63/1/012001}
\end{enumerate}

\section*{23\quad Particle-in-cell sheath simulation \textnormal{(\normalsize 28 references)}}
Particle-in-cell simulation resolves the kinetic, non-equilibrium structure of sheaths self-consistently. It generates synthetic observables and ground-truth flux for validating the inversion. It is the computational engine for the physics constraint.
\begin{enumerate}[leftmargin=1.4em]
\item \textbf{A fully kinetic, self-consistent particle simulation model of the collisionless plasma-sheath region} (1990). R. J. Procassini, C. K. Birdsall, E. C. Morse. Physics of Fluids B: Plasma Physics. \url{https://doi.org/10.1063/1.859229}
\item \textbf{Particle-in-Cell Simulation of Plasma Sheath Dynamics With Kinetic Ions} (2015). Chun-Wei Huang, Yen-Chiao Chen, Yasutaro Nishimura. IEEE Transactions on Plasma Science. \url{https://doi.org/10.1109/tps.2014.2382876}
\item \textbf{Particle-in-Cell Simulation of Magnetized Plasma Sheath} (2023). A. Castillo, K. Hara. 2023 IEEE International Conference on Plasma Science (ICOPS). \url{https://doi.org/10.1109/icops45740.2023.10481239}
\item \textbf{Kinetic Simulation of Plasma Sheath Relevant to Cross-Field Diodes} (2022). A. M. Castillo, K. Hara. 2022 IEEE International Conference on Plasma Science (ICOPS). \url{https://doi.org/10.1109/icops45751.2022.9813327}
\item \textbf{Particle-in-cell simulation of the pulsed sheath in a trench} (1996). T. E. Sheridan. IEEE Transactions on Plasma Science. \url{https://doi.org/10.1109/27.491689}
\item \textbf{Note on boundary conditions in plasma sheath simulations using the particle-in-cell algorithm} (2001). B. Briehl, H. M. Urbassek. IEEE Transactions on Plasma Science. \url{https://doi.org/10.1109/27.964479}
\item \textbf{Particle simulation of the magnetized rf plasma sheath} (1992). A. C. J. Paes, R. D. Sydora, J. M. Dawson. Physics of Fluids B: Plasma Physics. \url{https://doi.org/10.1063/1.860139}
\item \textbf{Particle-in-cell simulation of grid sheath oscillations in a double-plasma device} (1997). A. Rohde, H. Klostermann, A. Piel. IEEE Transactions on Plasma Science. \url{https://doi.org/10.1109/27.649637}
\item \textbf{Application of the particle-in-cell method for numerical simulation of sheath plasma} (1998). D. S. Filippychev. Computational Mathematics and Modeling. \url{https://doi.org/10.1007/bf02409864}
\item \textbf{Multispecies Particle Simulation of the Hydrogen Plasma Sheath Region} (2006). S. Longo, P. Diomede, M. Capitelli. Contributions to Plasma Physics. \url{https://doi.org/10.1002/ctpp.200610075}
\item \textbf{Particle in cell calculation of plasma force on a small grain in a non-uniform collisional sheath} (2013). I. H. Hutchinson. Plasma Physics and Controlled Fusion. \url{https://doi.org/10.1088/0741-3335/55/11/115014}
\item \textbf{Self-Consistent 1D3V Kinetic Trajectory Simulation Model of Magnetized Plasma Sheath} (2018). R. Chalise, R. Khanal. 2018 IEEE International Conference on Plasma Science (ICOPS). \url{https://doi.org/10.1109/icops35962.2018.9575673}
\item \textbf{A kinetic trajectory simulation model for magnetized plasma sheath} (2012). R. Chalise, R. Khanal. Plasma Physics and Controlled Fusion. \url{https://doi.org/10.1088/0741-3335/54/9/095006}
\item \textbf{Kinetic trajectory simulation method for the multi-component magnetized plasma sheath} (2019). Suresh Basnet, Raju Khanal. Plasma Physics and Controlled Fusion. \url{https://doi.org/10.1088/1361-6587/ab1708}
\item \textbf{Effect of electron temperature in a magnetized plasma sheath using kinetic trajectory simulation} (2021). R. Chalise, S. K. Pandit, G. Thakur, R. Khanal. BIBECHANA. \url{https://doi.org/10.3126/bibechana.v18i1.29204}
\item \textbf{Photoelectron Sheath and Plasma Charging on the Lunar Surface: Semianalytic Solutions and Fully-Kinetic Particle-in-Cell Simulations} (2021). Jianxun Zhao, Xinpeng Wei, Xiaoping Du, Xiaoming He, Daoru Han. IEEE Transactions on Plasma Science. \url{https://doi.org/10.1109/tps.2021.3110946}
\item \textbf{Development of a Fully Kinetic Particle Simulation Code For Coupled Plasma-Dust Transport} (2022). J. Zhao, D. Han. 2022 IEEE International Conference on Plasma Science (ICOPS). \url{https://doi.org/10.1109/icops45751.2022.9813088}
\item \textbf{Particle-in-cell modeling of electron beam generated plasma} (2023). Shahid Rauf, D. Sydorenko, S. Jubin, W. Villafana, S. Ethier, A. Khrabrov, I. Kaganovich. arXiv:2302.14121. \url{http://arxiv.org/abs/2302.14121v1}
\item \textbf{Kinetic theory of particle-in-cell simulation plasma and the ensemble averaging technique} (2022). Michael Touati, Romain Codur, Frank Tsung, Viktor K. Decyk, Warren B. Mori, Luis O. Silva. arXiv:2208.06375. \url{http://arxiv.org/abs/2208.06375v1}
\item \textbf{Particle-in-Cell/Test-Particle Simulations of Technological Plasmas: Sputtering Transport in Capacitive Radio Frequency Discharges} (2016). . arXiv:1607.04069. \url{https://arxiv.org/abs/1607.04069}
\item \textbf{Latest progress in Hall thrusters plasma modelling} (2019). F. Taccogna, L. Garrigues. Reviews of Modern Plasma Physics. \url{https://link.springer.com/article/10.1007/s41614-019-0033-1}
\item \textbf{Hall thruster modeling with multiple simulation techniques: Model benchmarking, fluid-kinetic consistency, and experimental validation} (2026). . Journal of Applied Physics. \url{https://pubs.aip.org/aip/jap/article/139/9/093303/3381992/Hall-thruster-modeling-with-multiple-simulation}
\item \textbf{Hybrid plasma simulations of a magnetically shielded Hall thruster} (2022). . Journal of Applied Physics. \url{https://pubs.aip.org/aip/jap/article/131/10/103302/2836568/Hybrid-plasma-simulations-of-a-magnetically}
\item \textbf{Validation of retarding-field energy analyzer interpretation process in non-magnetic plasmas using particle-in-cell simulations and characterization of End-Hall 400 ion source} (2026). . ScienceDirect (Fusion Engineering and Design). \url{https://www.sciencedirect.com/science/article/pii/S0920379626003443}
\item \textbf{Sheath and presheath in ion-ion plasmas via particle-in-cell simulation}. . OSTI.GOV (journal article). \url{https://www.osti.gov/biblio/21016325}
\item \textbf{Particle in cell simulation of plasma thrusters} (2013). Konstantin Matyash. 2013 IEEE International Conference on Plasma Science (ICOPS). \url{https://doi.org/10.1109/plasma.2013.6635220}
\item \textbf{Particle Simulation Study of the Effect of Radial Electric Field on Scrape-off Layer Plasma and Sheath Formation} (2000). T. Takizuka, M. Hosokawa. Contributions to Plasma Physics. \url{https://doi.org/10.1002/1521-3986(200006)40:3/4<471::aid-ctpp471>3.0.co;2-u}
\item \textbf{Kinetic (particle-in-cell) simulation of nonlinear laser absorption in a finite-size plasma with a background inhomogeneous magnetic field} (2015). H. Mehdian, A. Kargarian, K. Hajisharifi. Physics of Plasmas. \url{https://doi.org/10.1063/1.4921934}
\end{enumerate}

\section*{24\quad Kinetic modelling of plasma sheaths \textnormal{(\normalsize 18 references)}}
Vlasov/Boltzmann-level kinetic treatments capture distribution-function effects that fluid models miss in thin, collisionless sheaths. They inform the priors and the forward operator of the inversion. They bound the validity of simpler models.
\begin{enumerate}[leftmargin=1.4em]
\item \textbf{Continuum kinetic and multi-fluid simulations of classical sheaths} (2017). P. Cagas, A. Hakim, J. Juno, B. Srinivasan. Physics of Plasmas. \url{https://doi.org/10.1063/1.4976544}
\item \textbf{Vlasov Simulation of Emissive Plasma Sheath with Energy-Dependent Secondary Emission Coefficient and Improved Modeling for Dielectric Charging Effects} (2022). Guang-Yu Sun, Shu Zhang, Bao-Hong Guo, An-Bang Sun, Guan-Jun Zhang. arXiv:2209.09567. \url{http://arxiv.org/abs/2209.09567v1}
\item \textbf{Collisional sheath solutions of a bi-species Vlasov-Poisson-Boltzmann boundary value problem} (2021). Mehdi Badsi. Kinetic \textbackslash{}\& Related Models. \url{https://doi.org/10.3934/krm.2020052}
\item \textbf{Grid-Point Requirements for Direct Kinetic Simulation of the Vlasov and Boltzmann Equations} (2022). W. H. R. Chan, I. D. Boyd. IEEE International Conference on Plasma Science (ICOPS). \url{https://doi.org/10.1109/icops45751.2022.9813347}
\item \textbf{A Kinetic Vlasov Model for Plasma Simulation Using Discontinuous Galerkin Method on Many-Core Architectures} (2016). Noah Reddell. OSTI (thesis/technical report). \url{https://doi.org/10.2172/1469331}
\item \textbf{Spectral element moment-closure methods for kinetic Vlasov models of plasma}. Christine Vaughan. Thesis (Iowa State University). \url{https://doi.org/10.31274/td-20240329-688}
\item \textbf{Plasma-sheath transition in the kinetic Tonks-Langmuir model} (2006). K.-U. Riemann. Physics of Plasmas. \url{https://doi.org/10.1063/1.2209928}
\item \textbf{Kinetic Bohm criterion in the Tonks-Langmuir model: Assumption or theorem?} (2019). . Physics of Plasmas. \url{https://pubs.aip.org/aip/pop/article/26/12/123505/253564/Kinetic-Bohm-criterion-in-the-Tonks-Langmuir-model}
\item \textbf{Plasma-material boundary conditions for discontinuous Galerkin continuum-kinetic simulations, with a focus on secondary electron emission} (2020). . Journal of Computational Physics. \url{https://www.sciencedirect.com/science/article/abs/pii/S0021999119309209}
\item \textbf{Large gyro-orbit model of ion velocity distribution in plasma near a wall in a grazing-angle magnetic field} (2020). . arXiv:2007.05806. \url{https://arxiv.org/abs/2007.05806}
\item \textbf{The Influence of Non-Extensive Electrons on the Characteristics of the Space-Charge Sheath in Tonks-Langmuir Kinetic Model} (2025). . Contributions to Plasma Physics. \url{https://onlinelibrary.wiley.com/doi/10.1002/ctpp.70047}
\item \textbf{Boltzmann and Vlasov equations in plasma physics} (2016). Gianpiero Colonna. Plasma Modeling (IOP book chapter). \url{https://doi.org/10.1088/978-0-7503-1200-4ch1}
\item \textbf{A solvable model of Vlasov-kinetic plasma turbulence in Fourier-Hermite phase space} (2018). T. Adkins, A. A. Schekochihin. Journal of Plasma Physics. \url{https://doi.org/10.1017/s0022377818000089}
\item \textbf{Ion energy distribution functions behind the sheaths of magnetized and non-magnetized radio frequency discharges} (2013). . Journal of Physics D: Applied Physics. \url{https://iopscience.iop.org/article/10.1088/0022-3727/46/8/084016}
\item \textbf{Laser-induced fluorescence measurements of ion fluctuations in electron and ion presheaths} (2020). R. Hood, S. D. Baalrud, R. L. Merlino, F. Skiff. Physics of Plasmas. \url{https://doi.org/10.1063/1.5142014}
\item \textbf{Laser induced fluorescence of argon ions in a plasma presheath} (2001). . Physics of Plasmas. \url{https://pubs.aip.org/aip/pop/article-abstract/8/5/1729/859196/Laser-induced-fluorescence-of-argon-ions-in-a}
\item \textbf{Transmission characteristics of millimeter and sub-terahertz channels through spatially ripple plasma sheath layers} (2024). Wenbo Liu, Peian Li, Da Li, Daniel M. Mittleman, Jianjun Ma. arXiv:2407.18476. \url{http://arxiv.org/abs/2407.18476v1}
\item \textbf{Communication through plasma sheaths} (2007). A. O. Korotkevich, A. C. Newell, V. E. Zakharov. arXiv:0704.3103. \url{http://arxiv.org/abs/0704.3103v2}
\end{enumerate}

\section*{25\quad Inverse problems and regularization in diagnostics \textnormal{(\normalsize 29 references)}}
This literature provides the theory of ill-posed inverse problems, regularization and reconstruction from indirect, noisy data, including plasma tomography and spectroscopic inversion. It supplies the mathematical machinery that turns partial optical data into a stable flux estimate. It is the methodological core of the framework.
\begin{enumerate}[leftmargin=1.4em]
\item \textbf{Tomographic laser absorption spectroscopy using Tikhonov regularization} (2014). . arXiv:1408.4330. \url{http://arxiv.org/abs/1408.4330}
\item \textbf{Using Poisson-regularized inversion of Bremsstrahlung emission to extract full EEDFs from x-ray pulse-height detector data} (2017). . arXiv:1712.05422. \url{http://arxiv.org/abs/1712.05422}
\item \textbf{A novel inversion algorithm for optical tomography of plasmas} (2000). . IEEE Conference Publication. \url{https://ieeexplore.ieee.org/document/854917/}
\item \textbf{Tomographic image reconstruction techniques for spectroscopic sources - I. Theory and computer simulations} (1988). . Spectrochimica Acta Part B (ScienceDirect). \url{https://www.sciencedirect.com/science/article/abs/pii/0584854788801654}
\item \textbf{Emission plasma tomography with large acceptance angle apertures relevant to laser induced plasma spectroscopy} (2011). . Spectrochimica Acta Part B (ScienceDirect). \url{https://www.sciencedirect.com/science/article/abs/pii/S0584854711000693}
\item \textbf{Abel transform of laser-induced plasma image: Overcoming challenges of noisy and insufficient data for reliable reconstruction} (2024). . Spectrochimica Acta Part B (ScienceDirect). \url{https://www.sciencedirect.com/science/article/abs/pii/S0584854724002246}
\item \textbf{Bayesian modelling of Thomson scattering and multichannel interferometer diagnostics using Gaussian processes} (2020). . Journal article (ResearchGate record). \url{https://www.researchgate.net/publication/338436458_Bayesian_modelling_of_Thomson_scattering_and_multichannel_interferometer_diagnostics_using_Gaussian_processes}
\item \textbf{Bayesian inference of plasma parameters from collective Thomson scattering technique on a gas-puff near stagnation} (2023). . Scientific Reports (Nature). \url{https://www.nature.com/articles/s41598-023-40014-x}
\item \textbf{Two-dimensional inference of divertor plasma characteristics: advancements to a multi-instrument Bayesian analysis system} (2025). . Plasma Physics and Controlled Fusion (IOP). \url{https://iopscience.iop.org/article/10.1088/1361-6587/adab1b}
\item \textbf{Radially resolved optical emission spectral imaging study of an atmospheric-pressure plasma} (2021). . Journal article (NSF PAR). \url{https://par.nsf.gov/servlets/purl/10231555}
\item \textbf{Inverse Problems in Magnetohydrodynamics: Theoretical and Experimental Aspects} (2003). Frank Stefani, Thomas Gundrum, Gunter Gerbeth, Uwe Guenther, Mingtian Xu. arXiv:physics/0312093. \url{http://arxiv.org/abs/physics/0312093}
\item \textbf{Particular solutions of the inverse problem for 1D Vlasov-Maxwell equilibria using Hermite polynomials} (2015). O. Allanson, T. Neukirch, S. Troscheit, F. Wilson. arXiv:1508.07925. \url{http://arxiv.org/abs/1508.07925}
\item \textbf{Wavelet transforms and their applications to MHD and plasma turbulence: a review} (2015). Marie Farge, Kai Schneider. arXiv:1508.05650. \url{http://arxiv.org/abs/1508.05650}
\item \textbf{General dispersion properties of magnetized plasmas with drifting bi-Kappa distributions (DIS-K solver)} (2021). R. A. Lopez, S. M. Shaaban, M. Lazar. arXiv:2102.12513. \url{http://arxiv.org/abs/2102.12513}
\item \textbf{Linear Theory of Electron-Plasma Waves at Arbitrary Collisionality} (2018). R. Jorge, P. Ricci, S. Brunner, S. Gamba, V. Konovets, N. F. Loureiro, L. M. Perrone, N. Teixeira. arXiv:1811.12855. \url{http://arxiv.org/abs/1811.12855}
\item \textbf{A note on convergence of solutions of total variation regularized linear inverse problems} (2017). Jose A. Iglesias, Gwenael Mercier, Otmar Scherzer. arXiv:1711.06495. \url{http://arxiv.org/abs/1711.06495}
\item \textbf{Variational source conditions and stability estimates for inverse electromagnetic medium scattering problems} (2015). Frederic Weidling, Thorsten Hohage. arXiv:1512.06586. \url{http://arxiv.org/abs/1512.06586}
\item \textbf{Inverse spectral problems for Sturm-Liouville operators with singular potentials} (2002). R. O. Hryniv, Ya. V. Mykytyuk. arXiv:math/0211247. \url{http://arxiv.org/abs/math/0211247}
\item \textbf{Efficient Determination of Optimal Regularization Parameter for Inverse Problem in EXAFS Spectroscopy} (2005). M. Kunicke, I. Yu. Kamensky, et al.. Physica Scripta. \url{https://doi.org/10.1238/physica.topical.115a00237}
\item \textbf{Improvement of the Tikhonov Regularization Method for Solving the Inverse Problem of Mossbauer Spectroscopy} (2018). O. M. Nemtsova, G. N. Konygin. Journal of Applied Spectroscopy. \url{https://doi.org/10.1007/s10812-018-0741-2}
\item \textbf{Maximum entropy regularization of the geomagnetic core field inverse problem} (2007). Andrew Jackson, Catherine Constable. Geophysical Journal International. \url{https://doi.org/10.1111/j.1365-246x.2007.03530.x}
\item \textbf{Iterative Regularization Techniques in Image Reconstruction} (2000). Martin Hanke. Surveys on Solution Methods for Inverse Problems (Springer). \url{https://doi.org/10.1007/978-3-7091-6296-5_3}
\item \textbf{Electron paramagnetic resonance image reconstruction with total variation and curvelets regularization} (2017). Sylvain Durand, Yves-Michel Frapart. Inverse Problems (IOP). \url{https://doi.org/10.1088/1361-6420/aa8412}
\item \textbf{Stability, reconstruction formula and regularization for an inverse source hyperbolic problem by a control method} (1995). M. Yamamoto. Inverse Problems (IOP). \url{https://doi.org/10.1088/0266-5611/11/2/013}
\item \textbf{Total variation regularization for a backward time-fractional diffusion problem} (2013). Liyan Wang, Jijun Liu. Inverse Problems (IOP). \url{https://doi.org/10.1088/0266-5611/29/11/115013}
\item \textbf{Regularization strategy for an inverse problem for a 1+1 dimensional wave equation} (2016). Jussi Korpela, Matti Lassas, others. Inverse Problems (IOP). \url{https://doi.org/10.1088/0266-5611/32/6/065001}
\item \textbf{Regularization strategies for a two-dimensional inverse heat conduction problem} (2007). Zhi Qian, Chu-Li Fu. Inverse Problems (IOP). \url{https://doi.org/10.1088/0266-5611/23/3/013}
\item \textbf{On the regularization of the inverse conductivity problem with discontinuous conductivities} (2008). Luca Rondi. Inverse Problems and Imaging (AIMS). \url{https://doi.org/10.3934/ipi.2008.2.397}
\item \textbf{A dual approach to Kohn-Vogelius regularization applied to a data completion problem} (2020). Fabien Caubet, Jeremie Darde. Inverse Problems (IOP). \url{https://doi.org/10.1088/1361-6420/ab7868}
\end{enumerate}

\section*{26\quad Bayesian inference and integrated data analysis \textnormal{(\normalsize 26 references)}}
Integrated data analysis and Bayesian inference combine heterogeneous diagnostics under a common physical model with quantified uncertainty. This is precisely the sensor-fusion paradigm the project adopts. It enables principled combination of the four optical diagnostics with physics priors.
\begin{enumerate}[leftmargin=1.4em]
\item \textbf{Simultaneous kinetic profile and magnetic equilibrium inference with Bayesian integrated data analysis in preparation for ITER} (2025). S. S. Denk, T. B. Amara, T. Odstrcil, L. Stagner, C. Akcay, T. Slendebroek, S. P. Smith, M. A. Van Zeeland, T. Akiyama, R. Nazikian. Plasma Physics and Controlled Fusion. \url{https://doi.org/10.1088/1361-6587/adc8d1}
\item \textbf{Integrated data analysis of fusion diagnostics by means of the Bayesian probability theory} (2004). R. Fischer, A. Dinklage. Review of Scientific Instruments. \url{https://doi.org/10.1063/1.1787607}
\item \textbf{Bayesian Integrated Data Analysis and Experimental Design for External Magnetic Plasma Diagnostics in DEMO} (2025). Jeffrey De Rycke, Alfredo Pironti, Marco Ariola, Antonio Quercia, Geert Verdoolaege. Physical Sciences Forum (43rd MaxEnt Workshop). \url{https://doi.org/10.3390/psf2025012013}
\item \textbf{Bayesian inference applied to electron temperature data: computational performances and diagnostics integration} (2022). A. Fassina, D. Abate, P. Franz. Journal of Instrumentation. \url{https://doi.org/10.1088/1748-0221/17/09/c09012}
\item \textbf{A Review of the Bayesian Method in Nuclear Fusion Diagnostic Research} (2024). . Journal of Fusion Energy. \url{https://doi.org/10.1007/s10894-024-00404-z}
\item \textbf{2022 Review of Data-Driven Plasma Science} (2022). R. Anirudh, R. Archibald, M. S. Asif, M. M. Becker, S. Benkadda, P.-T. Bremer, R. M. Churchill, J. Citrin, et al.. arXiv:2205.15832. \url{https://arxiv.org/abs/2205.15832}
\item \textbf{Integrated Data Analysis and Validation} (2024). . arXiv:2411.09270. \url{https://arxiv.org/abs/2411.09270}
\item \textbf{Bayesian modelling of multiple plasma diagnostics at Wendelstein 7-X} (2024). . Nuclear Fusion. \url{https://doi.org/10.1088/1741-4326/ad6e02}
\item \textbf{Bayesian inference of electron density and ion temperature profiles from neutral beam and halo Balmer-alpha emission at Wendelstein 7-X} (2024). . Plasma Physics and Controlled Fusion. \url{https://doi.org/10.1088/1361-6587/ad3c1d}
\item \textbf{Bayesian inference of spectrometric data and validation with numerical simulations of plasma sheath diagnostics of a plasma focus discharge} (2022). . Scientific Reports. \url{https://www.nature.com/articles/s41598-022-19764-7}
\item \textbf{Bayesian analysis of optical emission spectroscopy measurements coupled to a collisional-radiative model: application to low-pressure misty plasmas} (2025). . Plasma Sources Science and Technology. \url{https://doi.org/10.1088/1361-6595/ae37a2}
\item \textbf{Bayesian plasma model selection for Thomson scattering} (2024). . Review of Scientific Instruments. \url{https://pubs.aip.org/aip/rsi/article/95/4/043004/3283078/Bayesian-plasma-model-selection-for-Thomson}
\item \textbf{Thomson scattering analysis with the Bayesian probability theory} (2002). R. Fischer, others. Plasma Physics and Controlled Fusion. \url{https://www.researchgate.net/publication/27265771_Thomson_scattering_analysis_with_the_Bayesian_probability_theory}
\item \textbf{Uncertainty quantification of a multi-component Hall thruster model at varying facility pressures} (2025). . arXiv:2507.08113. \url{https://arxiv.org/abs/2507.08113}
\item \textbf{Hall thruster model improvement by multidisciplinary uncertainty quantification} (2024). . Journal of Electric Propulsion. \url{https://doi.org/10.1007/s44205-024-00079-w}
\item \textbf{Calibration of a Hall Thruster numerical model using a Bayesian approach} (2026). . Nonlinear Dynamics. \url{https://doi.org/10.1007/s11071-026-12217-0}
\item \textbf{Bayesian analysis of triple Langmuir probe measurements for the characterization of Hall thruster plasmas} (2019). . Review of Scientific Instruments. \url{https://pubs.aip.org/aip/rsi/article-abstract/90/2/023502}
\item \textbf{Deep learning for Gaussian process soft x-ray tomography model selection in the ASDEX Upgrade tokamak} (2020). . Review of Scientific Instruments. \url{https://pubs.aip.org/aip/rsi/article/91/10/103501/363552/Deep-learning-for-Gaussian-process-soft-x-ray}
\item \textbf{Gaussian Process Tomography at ASDEX Upgrade with Magnetic Equilibrium Information and Nonstationary Kernels} (2022). . Fusion Science and Technology. \url{https://doi.org/10.1080/15361055.2022.2072659}
\item \textbf{Integrated Data Analysis of Plasma Electron Density Profile Tomography for HL-3 with Gaussian Process Regression} (2025). . arXiv:2502.08882. \url{https://arxiv.org/abs/2502.08882}
\item \textbf{Integrated Data Analysis of Profile Diagnostics at ASDEX Upgrade} (2010). R. Fischer, others. Fusion Science and Technology. \url{https://www.researchgate.net/publication/46171842_Integrated_Data_Analysis_of_Profile_Diagnostics_at_ASDEX_Upgrade}
\item \textbf{Thinking Bayesian for plasma physicists} (2024). . Physics of Plasmas. \url{https://pubs.aip.org/aip/pop/article/31/5/050901/3295473/Thinking-Bayesian-for-plasma-physicists}
\item \textbf{Neural network approximation of Bayesian models for the inference of ion and electron temperature profiles at W7-X} (2019). . Plasma Physics and Controlled Fusion. \url{https://doi.org/10.1088/1361-6587/ab1d26}
\item \textbf{Modelling of JET Diagnostics Using Bayesian Graphical Models} (2011). J. Svensson, A. Werner. Contributions to Plasma Physics. \url{https://onlinelibrary.wiley.com/doi/10.1002/ctpp.201000058}
\item \textbf{Accelerated Bayesian inference of plasma profiles with self-consistent MHD equilibria at W7-X via neural networks} (2023). . arXiv:2305.12889. \url{https://arxiv.org/abs/2305.12889}
\item \textbf{Neural network surrogates of Bayesian diagnostic models for fast inference of plasma parameters} (2021). . Review of Scientific Instruments. \url{https://pubs.aip.org/aip/rsi/article/92/3/033531/1061541/Neural-network-surrogates-of-Bayesian-diagnostic}
\end{enumerate}

\section*{27\quad Gaps and Positioning}
Across the themes above, three findings are consistent. First, the physics of the sheath and the definition of ion energy flux are mature, yet the flux itself remains \emph{predicted rather than directly measured} under realistic, time-varying conditions. Second, intrusive probes perturb the sheath they sample, while each individual optical diagnostic captures only a partial slice of the plasma state, and few measurements are time-resolved. Third, the mathematical tools for combining heterogeneous data under physical constraints (inverse methods, regularization, and Bayesian integrated data analysis) are well developed but have not been assembled into a validated, non-intrusive, time-resolved reconstruction of ion energy flux at plasma-surface boundaries. The proposed framework occupies precisely this gap: it fuses four complementary non-perturbing optical diagnostics and regularizes the reconstruction with first-principles sheath and kinetic models, delivering a validated, spatially and temporally resolved estimate of ion energy flux. This is the project's central contribution.
\end{document}